\oldpage[IV-3]{116}
\section{Propriedades topológicas dos morfismos planos de apresentação finita. Critérios de planicidade}
\label{section:IV.11}

Enquanto na \S~2 consideramos os enunciados relativos à planicidade que não dependem de nenhuma hipótese de finitude, e a \S~6 estuda a noção de planicidade no âmbito dos pré-esquemas localmente noetherianos (mas sem hipótese de finitude para os morfismos), a presente seção é dedicada à noção de $f$-planicidade no caso em que o morfismo $f:X\to Y$ é \emph{localmente de apresentação finita}.
O interesse da noção de morfismo plano de apresentação finita reside em que ela parece exprimir tecnicamente, do modo mais adequado, a noção intuitiva de ``família de pré-esquemas algébricos parametrizada por um esquema $Y$'', cujo estudo é um dos principais objetos da Geometria Algébrica.
Aliás, mesmo que inicialmente só nos interessasse o caso de um esquema de base noetheriano, é indispensável, por certas razões técnicas (por exemplo, para algumas aplicações da teoria da ``descida'', que leva a introduzir esquemas não necessariamente noetherianos), não nos limitar a esse caso quando se trata de problemas de natureza essencialmente relativa ligados a morfismos localmente de apresentação finita.
Seguiremos sistematicamente esse princípio, já respaldado pelos resultados das \S\S~8 e 9, em todo o restante deste capítulo, e até mesmo no restante deste Tratado, ainda que por vezes isso sacrifique a simplicidade de algumas demonstrações, que hipóteses noetherianas às vezes permitem simplificar\footnote{Esse princípio também se inspira na necessidade de dar lugar, como ``espaços de parâmetros'' para famílias de esquemas algébricos, a espaços anelados arbitrários (e mesmo a ``topos'' anelados arbitrários), para os quais, em geral, já não cabem hipóteses noetherianas. Parece bastante claro que não se poderá evitar por muito mais tempo essa nova extensão da Geometria Algébrica, e convém desde já desenvolver as noções e técnicas de natureza ``relativa'' da teoria dos esquemas de maneira que possam adaptar-se praticamente sem alterações a esse quadro mais geral.}.
Na presente seção, isso nos leva a retomar, no contexto da ``apresentação finita'' (em particular no n$^\circ$~3), certos enunciados de planicidade já obtidos no contexto noetheriano.
O instrumento técnico essencial para efetuar a redução ao caso noetheriano é o teorema da compatibilidade da planicidade com os limites projetivos de pré-esquemas \sref{IV.11.2.6}, que completa os resultados gerais da \S~8.
Demonstramos também, de passagem, em \sref{IV.11.3.1}, um resultado frequentemente utilizado

\oldpage[IV-3]{117}

mais adiante, que implica que o conjunto dos pontos de planicidade de um morfismo localmente de apresentação finita é aberto.

Nos n$^{\mathrm{os}}$~4 a 8, estudamos a questão da ``descida'' da planicidade, que consiste em encontrar condições úteis sobre um morfismo de mudança de base $Y'\to Y$ (não plano em geral) para poder concluir que, se $X\times_YY'$ é plano sobre $Y'$, então $X$ é plano sobre $Y$.
Esses resultados, mais técnicos que os dos n$^{\mathrm{os}}$~1 a 3, são utilizados com menor frequência no restante; desempenharão, contudo, papel importante nas técnicas de construção não projetivas, no capítulo seguinte.
O único resultado dos n$^{\mathrm{os}}$~4 a 8 utilizado posteriormente no cap.~IV é o critério valuativo de planicidade \sref{IV.11.8}, que aplicaremos em \sref{IV.15.2}.

Por fim, os n$^{\mathrm{os}}$~9 e 10 são dedicados ao estudo de uma noção que torna mais precisa, na teoria dos esquemas, a densidade em sentido topológico, a saber, a noção de família de subpré-esquemas \emph{esquematicamente densa} em um pré-esquema dado, e em particular ao estudo do comportamento dessa noção sob mudança de base (plana ou arbitrária).
Essa noção é sobretudo utilizada, por ora, no estudo dos esquemas em grupos.

\subsection{Conjuntos de planicidade (caso noetheriano)}
\label{subsection:IV.11.1}

\begin{theorem}[11.1.1]
\label{IV.11.1.1}
Sejam $Y$ um pré-esquema localmente noetheriano, $f:X\to Y$ um morfismo localmente de tipo finito e $\sh{F}$ um $\sh{O}_X$-Módulo coerente. Então o conjunto $U$ dos $x\in X$ tais que $\sh{F}$ é $f$-plano no ponto $x$ é aberto em $X$.
\end{theorem}

\begin{proof}
Como a questão é evidentemente local em $X$ e em $Y$, podemos supor $Y=\operatorname{Spec}(A)$, $X=\operatorname{Spec}(B)$, com $A$ noetheriano e $B$ uma $A$-álgebra de tipo finito. Temos então $\sh{F}=\widetilde M$, sendo $M$ um $B$-módulo de tipo finito. Aplicando o critério $(\mathbf 0_{\mathrm{III}},\,9.2.6)$, basta, portanto, demonstrar a seguinte afirmação:

\begin{env}[11.1.1.1]
\label{IV.11.1.1.1}
\emph{Sejam $A$ um anel noetheriano, $B$ uma $A$-álgebra de tipo finito, $M$ um $B$-módulo de tipo finito, $\mathfrak q$ um ideal primo de $B$ e $\mathfrak p$ sua imagem inversa em $A$. Suponhamos que $M_{\mathfrak q}$ seja um $A_{\mathfrak p}$-módulo plano. Então existe $g\in B-\mathfrak q$ tal que, para todo ideal primo $\mathfrak q'\supset\mathfrak q$ de $B$ com $g\notin\mathfrak q'$, $M_{\mathfrak q'}$ é um $A_{\mathfrak p'}$-módulo plano, sendo $\mathfrak p'$ a imagem inversa de $\mathfrak q'$ em $A$ (o que equivale, aliás, a dizer $(\mathbf 0_{\mathrm I},\,6.3.1)$ que $M_{\mathfrak q'}$ é um $A$-módulo plano).}
\end{env}

Para isso, consideremos $B_{\mathfrak q'}$ como uma $A$-álgebra; temos evidentemente $\mathfrak pB_{\mathfrak q'}\subset\mathfrak q'B_{\mathfrak q'}$. Sabe-se então $(\mathbf 0_{\mathrm{III}},\,10.2.2)$ que, para $M_{\mathfrak q'}$ ser um $A$-módulo plano, é necessário e suficiente que $M_{\mathfrak q'}/\mathfrak pM_{\mathfrak q'}$ seja um $(A/\mathfrak p)$-módulo plano e que $\operatorname{Tor}_1^A(M_{\mathfrak q'},A/\mathfrak p)=0$.
Ora, $M_{\mathfrak q'}=M\otimes_BB_{\mathfrak q'}$; como $B_{\mathfrak q'}$ é plano sobre $B$, temos
\[
  M_{\mathfrak q'}/\mathfrak pM_{\mathfrak q'}
  =(M/\mathfrak pM)_{\mathfrak q'}
\]
e
\[
  \operatorname{Tor}_1^A(M_{\mathfrak q'},A/\mathfrak p)
  =(\operatorname{Tor}_1^A(M,A/\mathfrak p))_{\mathfrak q'}
\]
(definindo o Tor por meio de uma resolução projetiva de $A/\mathfrak p$); pela mesma razão, como devemos ter $g\notin\mathfrak q'$, $B_{\mathfrak q'}$ é plano sobre $B_g$, logo
\[
  (M/\mathfrak pM)_{\mathfrak q'}
  =((M/\mathfrak pM)_g)_{\mathfrak q'B_g}
\]
e
\[
  (\operatorname{Tor}_1^A(M,A/\mathfrak p))_{\mathfrak q'}
  =((\operatorname{Tor}_1^A(M,A/\mathfrak p))_g)_{\mathfrak q'B_g},
\]
onde, nessas fórmulas, $M/\mathfrak pM$ e $\operatorname{Tor}_1^A(M,A/\mathfrak p)$ são considerados como $B$-módulos. Levando em conta $(\mathbf 0_{\mathrm I},\,6.3.2)$, vemos que basta provar o seguinte lema.

\begin{lemma}[11.1.1.2]
\label{IV.11.1.1.2}
Sob as condições de \sref{IV.11.1.1.1}, existe $g\in B-\mathfrak q$ tal que:
\begin{enumerate}[label=(\roman*)]
  \item $(M/\mathfrak pM)_g$ seja um $(A/\mathfrak p)$-módulo plano;
  \item $(\operatorname{Tor}_1^A(M,A/\mathfrak p))_g=0$.
\end{enumerate}
\end{lemma}

\oldpage[IV-3]{118}

Pelo teorema da planicidade genérica \sref{IV.6.9.1}, aplicado ao domínio integral $A/\mathfrak p$, à $(A/\mathfrak p)$-álgebra de tipo finito $B/\mathfrak pB$ e ao $(B/\mathfrak pB)$-módulo de tipo finito $M/\mathfrak pM$, existe $h\in A-\mathfrak p$ tal que, denotando por $\bar h$ sua imagem canônica em $A/\mathfrak p$, $(M/\mathfrak pM)_{\bar h}$ seja um $(A/\mathfrak p)$-módulo plano.
Por outro lado, como $M_{\mathfrak q}$ é um $A_{\mathfrak p}$-módulo plano e, portanto, um $A$-módulo plano $(\mathbf 0_{\mathrm I},\,6.3.1)$, temos $\operatorname{Tor}_1^A(M_{\mathfrak q},A/\mathfrak p)=0$, o que também se escreve $(\operatorname{Tor}_1^A(M,A/\mathfrak p))_{\mathfrak q}=0$.
Como $A$ e $B$ são noetherianos, $\operatorname{Tor}_1^A(M,A/\mathfrak p)$ é um $B$-módulo de tipo finito; logo, por $(\mathbf 0_{\mathrm I},\,5.2.2)$, existe $g\in B-\mathfrak q$ tal que $(\operatorname{Tor}_1^A(M,A/\mathfrak p))_g=0$.
Além disso, $(M/\mathfrak pM)_{\bar h}=(M/\mathfrak pM)_h$ (considerando-se no segundo membro $M/\mathfrak pM$ como um $A$-módulo); ademais, como $(M/\mathfrak pM)_h$ é um $B$-módulo, $(M/\mathfrak pM)_{hg}$ continua a ser um $(A/\mathfrak p)$-módulo plano, pois se escreve $(M/\mathfrak pM)_{\overline{hg}}$, onde $\overline{hg}$ é a imagem canônica de $hg$ em $B/\mathfrak pB$, e basta aplicar $(\mathbf 0_{\mathrm I},\,6.3.2)$; por fim, $(\operatorname{Tor}_1^A(M,A/\mathfrak p))_{hg}=0$ e $hg\in B-\mathfrak q$, porque $h\in A-\mathfrak p$.
\end{proof}

\begin{corollary}[11.1.2]
\label{IV.11.1.2}
Sejam $Y$ um pré-esquema localmente noetheriano, $f:X\to Y$ um morfismo localmente de tipo finito, $\sh{F}$ e $\sh{F}'$ dois $\sh{O}_X$-Módulos coerentes, e $u:\sh{F}'\to\sh{F}$ um homomorfismo de $\sh{O}_X$-Módulos. Suponhamos que $\sh{F}$ seja $f$-plano. Então o conjunto $U$ dos $x\in X$ para os quais, pondo $y=f(x)$, o homomorfismo
\[
  u_x\otimes1:\sh{F}'_x\otimes_{\sh{O}_y}k(y)\longrightarrow
  \sh{F}_x\otimes_{\sh{O}_y}k(y)
\]
é injetivo, é aberto em $X$.
\end{corollary}

\begin{proof}
Com efeito, seja $\sh{N}$ (resp.\ $\sh{P}$) o núcleo (resp.\ o conúcleo) de $u$; apliquemos $(\mathbf 0_{\mathrm{III}},\,10.2.4)$\footnote{Eis uma demonstração de $(\mathbf 0_{\mathrm{III}},\,10.2.4)$ que não foi publicada na \emph{Algèbre commutative} de N.~Bourbaki. Tendo em conta $(\mathbf 0_{\mathrm I},\,6.1.2)$, basta mostrar que b) implica a). Ponhamos $P=\operatorname{Im}(u)$, $Q=\operatorname{Coker}(u)$ e $R=\operatorname{Ker}(u)$. A composta $M\otimes k\to P\otimes k\to N\otimes k$ é injetiva, e $M\otimes k\to P\otimes k$ é sobrejetiva; portanto, $P\otimes k\to N\otimes k$ é injetiva e $M\otimes k\to P\otimes k$ é bijetiva. A sequência exata $0\to P\to N\to Q\to0$ fornece a sequência exata
\[
  0=\operatorname{Tor}_1^A(N,k)\to\operatorname{Tor}_1^A(Q,k)
  \to P\otimes k\to N\otimes k\to Q\otimes k\to0
\]
e, como $P\otimes k\to N\otimes k$ é injetiva, temos $\operatorname{Tor}_1^A(Q,k)=0$; como $Q$ é um $B$-módulo de tipo finito, $(\mathbf 0_{\mathrm{III}},\,10.2.2)$ mostra que $Q$ é um $A$-módulo plano; então $P$ também é um $A$-módulo plano por $(\mathbf 0_{\mathrm I},\,6.1.2)$. Sendo exata a sequência $0\to R\to M\to P\to0$, também o é $0\to R\otimes k\to M\otimes k\to P\otimes k\to0$ por $(\mathbf 0_{\mathrm I},\,6.1.2)$; como $M\otimes k\to P\otimes k$ é bijetiva, temos $R\otimes k=0$; mas, como $B$ é noetheriano, $R$ é um $B$-módulo de tipo finito e, portanto, $R=0$ pelo lema de Nakayama.} aos anéis locais $\sh{O}_y$ e $\sh{O}_x$ e aos $\sh{O}_x$-módulos $\sh{F}'_x$ e $\sh{F}_x$: dizer que $u_x\otimes1$ é injetivo equivale a dizer que $\sh{N}_x=0$ e que $\sh{P}$ é $f$-plano no ponto $x$. Ora, como $\sh{N}$ e $\sh{P}$ são coerentes $(\mathbf 0_{\mathrm I},\,5.3.4)$, o conjunto dos $x$ nos quais $\sh{N}_x=0$ é aberto $(\mathbf 0_{\mathrm I},\,5.2.2)$, e o conjunto dos $x$ nos quais $\sh{P}$ é $f$-plano é aberto por \sref{IV.11.1.1}; daí a conclusão.
\end{proof}

Em particular:

\begin{corollary}[11.1.3]
\label{IV.11.1.3}
Sejam $Y$ um pré-esquema localmente noetheriano, $f:X\to Y$ um morfismo plano e localmente de tipo finito, e $g$ uma seção de $\sh{O}_X$ sobre $X$. O conjunto dos $x\in X$ tais que $g_x\otimes1$ não seja divisor de zero em $\sh{O}_x\otimes_{\sh{O}_{f(x)}}k(f(x))$ é aberto em $X$.
\end{corollary}

\begin{proof}
Basta aplicar \sref{IV.11.1.2} ao endomorfismo de $\sh{O}_X$ definido por $g$.
\end{proof}

\begin{corollary}[11.1.4]
\label{IV.11.1.4}
Sejam $Y$ um pré-esquema localmente noetheriano, $f:X\to Y$ um morfismo localmente de tipo finito e $\sh{F}$ um $\sh{O}_X$-Módulo coerente e $f$-plano. Seja $(g_i)_{1\leq i\leq n}$ uma sequência de seções de $\sh{O}_X$ sobre $X$. Então o conjunto $U$ dos $x\in X$ tais que a sequência $((g_i)_x\otimes1)$ seja $(\sh{F}_x\otimes_{\sh{O}_{f(x)}}k(f(x)))$-regular é aberto em $X$.
\end{corollary}

\begin{proof}
\oldpage[IV-3]{119}

Como $\sh{F}$ é $f$-plano, resulta de $(\mathbf 0,\,15.1.16)$ que $U$ é também o conjunto dos $x\in X$ tais que a sequência $((g_i)_x)$ é $\sh{F}_x$-regular e que o $\sh{O}_x$-módulo $\sh{G}_x$, onde $\sh{G}=\sh{F}/(\sum_{i=1}^n g_i\sh{F})$, é $f$-plano. Mas $\sh{F}$ e $\sh{G}$ são coerentes; portanto, o corolário decorre de \sref{IV.11.1.1} e de $(\mathbf 0,\,15.2.4)$.
\end{proof}

\begin{corollary}[11.1.5]
\label{IV.11.1.5}
Sejam $X$, $Y$ e $Z$ três pré-esquemas localmente noetherianos, $f:X\to Y$ e $g:Y\to Z$ dois morfismos de tipo finito, e $\sh{F}$ um $\sh{O}_X$-Módulo coerente. Então o conjunto $U$ dos $y\in Y$ tais que, para toda generização $y'$ de $y$, $\sh{F}$ é $(g\circ f)$-plano em todos os pontos de $f^{-1}(y')$ (isto é, tais que $\sh{F}\otimes_Y\operatorname{Spec}(\sh{O}_{Y,y})$ é plano relativamente ao morfismo $X\times_Y\operatorname{Spec}(\sh{O}_{Y,y})\to Z$) é aberto em $Y$.
\end{corollary}

\begin{proof}
Se $V$ é o conjunto dos $x\in X$ nos quais $\sh{F}$ é $(g\circ f)$-plano, $U$ é o conjunto dos $y\in Y$ tais que, para toda generização $y'$ de $y$, temos $f^{-1}(y')\subset V$. Ora, $V$ é aberto por \sref{IV.11.1.1} e, portanto, localmente construtível em $X$; o conjunto $W$ dos $y\in Y$ tais que $f^{-1}(y)\subset V$ é igual a $Y-f(X-V)$ e, assim, também é localmente construtível em $Y$ pelo teorema de Chevalley \sref{IV.1.8.4}. Mas decorre então de $(\mathbf 0_{\mathrm{III}},\,9.2.5)$ que os pontos de $U$ são os pontos interiores de $W$; daí a conclusão.
\end{proof}

\begin{corollary}[11.1.6]
\label{IV.11.1.6}
Sejam $A$ um anel noetheriano, $B$ uma $A$-álgebra de tipo finito, $C$ uma $B$-álgebra de tipo finito e $M$ um $C$-módulo de tipo finito. Então o conjunto dos $\mathfrak q\in\operatorname{Spec}(B)$ tais que $M_{\mathfrak q}$ é um $A$-módulo plano é aberto em $\operatorname{Spec}(B)$.
\end{corollary}

\begin{proof}
Levando em conta \sref{IV.2.1.2}, isso é consequência de \sref{IV.11.1.5}, aplicado a $Z=\operatorname{Spec}(A)$, $Y=\operatorname{Spec}(B)$, $X=\operatorname{Spec}(C)$ e $\sh{F}=\widetilde M$.
\end{proof}

Os resultados deste número serão generalizados em \sref{IV.11.3}, sem as hipóteses noetherianas.

\subsection{Planicidade de um limite projetivo de pré-esquemas}
\label{subsection:IV.11.2}

\begin{env}[11.2.1]
\label{IV.11.2.1}
Sejam $A$ um anel, $M$ e $N$ dois $A$-módulos e $A'$ uma $A$-álgebra; ponhamos $M'=M\otimes_AA'$ e $N'=N\otimes_AA'$. Recordemos $(\mathbf{III},\,6.3.8)$ que se define, para todo $i$, um homomorfismo canônico de $A$-módulos
\begin{equation}
  \psi_i:\operatorname{Tor}_i^A(M,N)\to\operatorname{Tor}_i^{A'}(M',N')
  \label{IV.11.2.1.1}
  \tag{11.2.1.1}
\end{equation}
da seguinte maneira: considera-se uma resolução à esquerda de $M$ por $A$-módulos livres
\begin{equation}
  \cdots\longrightarrow L_{i+1}\xrightarrow{f_i}L_i\longrightarrow\cdots
  \longrightarrow L_0\xrightarrow{\varepsilon}M\longrightarrow0
  \label{IV.11.2.1.2}
  \tag{11.2.1.2}
\end{equation}
da qual se deduz, por tensorização com $A'$, um complexo de $A'$-módulos
\begin{equation}
  \cdots\longrightarrow L'_{i+1}\xrightarrow{f'_i}L'_i\longrightarrow\cdots
  \longrightarrow L'_0\xrightarrow{\varepsilon'}M'\longrightarrow0
  \label{IV.11.2.1.3}
  \tag{11.2.1.3}
\end{equation}
em que se pôs $L'_i=L_i\otimes_AA'$, $f'_i=f_i\otimes1_{A'}$ e $\varepsilon'=\varepsilon\otimes1_{A'}$. Consideremos, por outro lado, uma resolução à esquerda de $M'$ por $A'$-módulos livres
\begin{equation}
  \cdots\longrightarrow L''_{i+1}\xrightarrow{f''_i}L''_i\longrightarrow\cdots
  \longrightarrow L''_0\xrightarrow{\varepsilon''}M'\longrightarrow0.
  \label{IV.11.2.1.4}
  \tag{11.2.1.4}
\end{equation}
\end{env}

\oldpage[IV-3]{120}

Como os $L'_i$ são $A'$-módulos livres, sabe-se $(\mathrm M,\,\mathrm V,\,1.1)$ que existem $A'$-homomorfismos $u_i$ que formam um diagrama comutativo
\[
\xymatrix{
\cdots\ar[r]&L'_{i+1}\ar[r]^{f'_i}\ar[d]_{u_{i+1}}&
L'_i\ar[r]\ar[d]_{u_i}&\cdots\ar[r]&
L'_0\ar[r]^{\varepsilon'}\ar[d]_{u_0}&M'\ar[r]\ar[d]_{1_{M'}}&0\\
\cdots\ar[r]&L''_{i+1}\ar[r]_{f''_i}&L''_i\ar[r]&\cdots\ar[r]&
L''_0\ar[r]_{\varepsilon''}&M'\ar[r]&0.
}
\]
Compondo o homomorfismo de complexos $u_\bullet:L'_\bullet\to L''_\bullet$ assim definido com o homomorfismo canônico $L_\bullet\to L'_\bullet$, obtém-se um homomorfismo de complexos de $A$-módulos $L_\bullet\to L''_\bullet$; observando que $(L_i\otimes_AN)\otimes_AA'=L'_i\otimes_{A'}N'$, deduz-se daí um homomorfismo de complexos de $A$-módulos $L_\bullet\otimes_AN\to L''_\bullet\otimes_{A'}N'$ e, passando à homologia, os homomorfismos canônicos \eref{IV.11.2.1.1}. Como os $u_i$ estão determinados a menos de homotopia $(\mathrm M,\,\mathrm V,\,1.1)$, os homomorfismos \eref{IV.11.2.1.1} não dependem da escolha dos $u_i$ nem da escolha das resoluções livres $L_\bullet$ e $L''_\bullet$.

\begin{sloppypar}
Como os $\operatorname{Tor}_i^{A'}(M',N')$ são $A'$-módulos, deduzem-se canonicamente de \eref{IV.11.2.1.1} $A'$-homomorfismos
\begin{equation}
\begin{gathered}
\varphi_i:\operatorname{Tor}_i^A(M,N)\otimes_AA'\\[-2pt]
\longrightarrow\operatorname{Tor}_i^{A'}(M',N').
\end{gathered}
\label{IV.11.2.1.5}
\tag{11.2.1.5}
\end{equation}
\end{sloppypar}

Consideremos agora dois homomorfismos de anéis
\[
\rho:A\to A^{(1)},\qquad \sigma:A^{(1)}\to A^{(2)},
\]
e seu composto $\sigma\circ\rho:A\to A^{(2)}$; ponhamos $M^{(j)}=M\otimes_AA^{(j)}$ e $N^{(j)}=N\otimes_AA^{(j)}$ para $j=1,2$. Então o homomorfismo canônico composto
\[
\operatorname{Tor}_i^A(M,N)\to
\operatorname{Tor}_i^{A^{(1)}}(M^{(1)},N^{(1)})\to
\operatorname{Tor}_i^{A^{(2)}}(M^{(2)},N^{(2)})
\]
coincide com o homomorfismo canônico deduzido de $\sigma\circ\rho$; isso decorre de que, se $L_\bullet^{(j)}$ é uma resolução livre de $M^{(j)}$, o diagrama
\[
\xymatrix{
L_\bullet\ar[r]&L_\bullet\otimes_AA^{(1)}\ar[r]\ar[d]&
L_\bullet^{(1)}\ar[d]\\
&L_\bullet\otimes_AA^{(2)}\ar[r]&
L_\bullet^{(1)}\otimes_{A^{(1)}}A^{(2)}\ar[r]&L_\bullet^{(2)}
}
\]
é comutativo.

\begin{env}[11.2.2]
\label{IV.11.2.2}
Com as notações de \sref{IV.11.2.1}, consideremos agora um \emph{sistema indutivo filtrante} de $A$-álgebras $(A_\alpha,g_{\beta\alpha})$ e, para cada índice $\alpha$, ponhamos $M_\alpha=M\otimes_AA_\alpha$ e $N_\alpha=N\otimes_AA_\alpha$; decorre então de \sref{IV.11.2.1} que, para cada $i$, $(\operatorname{Tor}_i^{A_\alpha}(M_\alpha,N_\alpha),\psi_{i\beta\alpha})$, onde $\psi_{i\beta\alpha}$ é o homomorfismo canônico \eref{IV.11.2.1.1} correspondente a $g_{\beta\alpha}$, é um \emph{sistema indutivo} de $A_\alpha$-módulos. Ponhamos
\[
A'=\varinjlim_\alpha A_\alpha,\qquad
M'=\varinjlim_\alpha M_\alpha=M\otimes_AA',\qquad
N'=\varinjlim_\alpha N_\alpha=N\otimes_AA';
\]
se designarmos por $g_\alpha:A_\alpha\to A'$ o homomorfismo canônico, daí se deduzem homomorfismos canônicos \eref{IV.11.2.1.1}
\[
\psi_{i\alpha}:\operatorname{Tor}_i^{A_\alpha}(M_\alpha,N_\alpha)
\to\operatorname{Tor}_i^{A'}(M',N')
\]
que, ainda em virtude de \sref{IV.11.2.1}, formam um sistema indutivo de homomorfismos; propomo-nos completar o resultado de $(\mathrm M,\,\mathrm V,\,9.4^*)$ mostrando que os
\begin{equation}
\psi_i=\varinjlim_\alpha\psi_{i\alpha}:
\varinjlim_\alpha\operatorname{Tor}_i^{A_\alpha}(M_\alpha,N_\alpha)
\to\operatorname{Tor}_i^{A'}(M',N')
\label{IV.11.2.2.1}
\tag{11.2.2.1}
\end{equation}
são \emph{isomorfismos de $A'$-módulos}. Para isso, procedemos como em $(\mathrm M,\,\mathrm V,\,9.5^*)$, associando a cada $M_\alpha$ sua \emph{resolução livre canônica}. Tudo se reduz, tendo em conta a exatidão do funtor $\varinjlim$, a demonstrar o
\end{env}

\oldpage[IV-3]{121}

\begin{lemma}[11.2.2.2]
\label{IV.11.2.2.2}
Sejam $(A_\alpha,g_{\beta\alpha})$ um sistema indutivo filtrante de anéis e $(M_\alpha,h_{\beta\alpha})$ um sistema indutivo de conjuntos; sejam $A'=\varinjlim A_\alpha$, $M'=\varinjlim M_\alpha$ e $g_\alpha:A_\alpha\to A'$, $h_\alpha:M_\alpha\to M'$ as aplicações canônicas. Para cada $\alpha$, seja $F(M_\alpha)$ o $A_\alpha$-módulo das combinações lineares formais de elementos de $M_\alpha$; seja, do mesmo modo, $F(M')$ o $A'$-módulo das combinações lineares formais de elementos de $M'$. Se $h^0_{\beta\alpha}:F(M_\alpha)\to F(M_\beta)$ (para $\alpha\leq\beta$) e $h'_\alpha:F(M_\alpha)\to F(M')$ são os $A_\alpha$-homomorfismos deduzidos respectivamente de $h_{\beta\alpha}$ e de $h_\alpha$, então $(F(M_\alpha),h^0_{\beta\alpha})$ é um sistema indutivo de $A_\alpha$-módulos e $(h'_\alpha)$ é um sistema indutivo de homomorfismos. Então
\[
h''=\varinjlim h'_\alpha:\varinjlim F(M_\alpha)\to
F(M')=F(\varinjlim M_\alpha)
\]
é um isomorfismo.
\end{lemma}

Para a demonstração, veja-se Bourbaki, \emph{Alg.}, cap.~II, 3$^{\mathrm a}$ ed., \S~6, número~6, cor. da prop.~10.

\begin{env}[11.2.3]
\label{IV.11.2.3}
Retomemos as notações de \sref{IV.11.2.1} e consideremos em particular o caso $i=1$; ponhamos $T=\operatorname{Tor}_1^A(M,N)$, $T'=T\otimes_AA'$ e $T''=\operatorname{Tor}_1^{A'}(M',N')$; então $\psi_1$ é o homomorfismo
\[
\operatorname{Ker}(f_0\otimes1_N)/\operatorname{Im}(f_1\otimes1_N)
\to
\operatorname{Ker}(f''_0\otimes1_{N'})/
\operatorname{Im}(f''_1\otimes1_{N'})
\]
deduzido, por passagem aos quocientes, da restrição $\operatorname{Ker}(f_0\otimes1_N)\to\operatorname{Ker}(f''_0\otimes1_{N'})$ de
\[
u_1\otimes1:L_1\otimes_AN\to L''_1\otimes_{A'}N'=L''_1\otimes_AN.
\]\footnote{Nota editorial: o impresso, p.~121, usa $L_0$ e $L''_0$ nesta aplicação. Mantemos aqui o grau $1$, pois $u_1$ é a componente de grau $1$ do morfismo de resoluções e $\operatorname{Ker}(f_0\otimes1_N)$ é um submódulo de $L_1\otimes_AN$. A leitura impressa é conservada nesta nota.}
Ponhamos $R=\operatorname{Ker}(\varepsilon)$ e $R''=\operatorname{Ker}(\varepsilon'')$; temos então as sequências exatas
\[
0\to R\xrightarrow{j}L_0\xrightarrow{\varepsilon}M\to0
\qquad\text{e}\qquad
0\to R''\xrightarrow{j''}L''_0\xrightarrow{\varepsilon''}M'\to0,
\]
das quais se deduzem as sequências exatas de homologia
\begin{equation}
0=\operatorname{Tor}_1^A(L_0,N)\to T\xrightarrow{\partial}R\otimes_AN
\xrightarrow{j\otimes1}L_0\otimes_AN\xrightarrow{\varepsilon\otimes1}
M\otimes_AN\to0
\label{IV.11.2.3.1}
\tag{11.2.3.1}
\end{equation}
\begin{equation}
0=\operatorname{Tor}_1^{A'}(L''_0,N')\to T''\xrightarrow{\partial''}
R''\otimes_{A'}N'\xrightarrow{j''\otimes1}L''_0\otimes_{A'}N'
\xrightarrow{\varepsilon''\otimes1}M'\otimes_{A'}N'\to0.
\label{IV.11.2.3.2}
\tag{11.2.3.2}
\end{equation}
Temos, por outro lado, um homomorfismo de $A'$-módulos
\[
v:R'=\operatorname{Ker}(\varepsilon)\otimes_AA'
\to\operatorname{Ker}(\varepsilon')\xrightarrow{u_0}
\operatorname{Ker}(\varepsilon'')=R''.
\]
\end{env}

\oldpage[IV-3]{122}

Mostremos que o diagrama
\begin{equation}
\xymatrix{
&T'\ar[r]^{\partial\otimes1}\ar[d]_{\varphi_1}&
R'\otimes_{A'}N'\ar[r]\ar[d]_{v\otimes1}&
L'_0\otimes_{A'}N'\ar[r]\ar[d]_{u_0\otimes1}&
M'\otimes_{A'}N'\ar[r]\ar@{=}[d]&0\\
0\ar[r]&T''\ar[r]_{\partial''}&R''\otimes_{A'}N'\ar[r]&
L''_0\otimes_{A'}N'\ar[r]&M'\otimes_{A'}N'\ar[r]&0
}
\label{IV.11.2.3.3}
\tag{11.2.3.3}
\end{equation}
é \emph{comutativo}. Para isso, verifica-se imediatamente $(\mathrm M,\,\mathrm{IV},\,1)$ que o homomorfismo $\partial:T\to R\otimes_AN$ provém, no caso presente, por passagem ao quociente, do homomorfismo $w:\operatorname{Ker}(f_0\otimes1_N)\to R\otimes_AN$, restrição do homomorfismo $g_0\otimes1_N:L_1\otimes_AN\to R\otimes_AN$, em que $g_0:L_1\to R$ é tal que $f_0=j\circ g_0$; o mesmo vale para $\partial''$. Basta, portanto, ver que o diagrama
\[
\xymatrix{
\operatorname{Ker}(f_0\otimes1_N)\ar[r]^w\ar[d]&R\otimes_AN\ar[r]&
R'\otimes_{A'}N'=(R\otimes_AN)\otimes_AA'\ar[d]_{v\otimes1}\\
\operatorname{Ker}(f''_0\otimes1_{N'})\ar[rr]&&R''\otimes_{A'}N'
}
\]
é comutativo, o que decorre da comutatividade do diagrama
\[
\xymatrix{
L_1\ar[r]\ar[d]&R\ar[d]\\
L''_1\ar[r]&R''.
}
\]

\begin{lemma}[11.2.4]
\label{IV.11.2.4}
Sejam $A$ um anel, $C$ uma $A$-álgebra, $M$ um $A$-módulo, $N$ um $C$-módulo e $A'$ uma $A$-álgebra. Ponhamos $C'=C\otimes_AA'$, $M'=M\otimes_AA'$ e $N'=N\otimes_AA'=N\otimes_CC'$. Suponhamos que $M\otimes_AN$ seja um $C$-módulo plano. Então o homomorfismo canônico
\[
\varphi_1:\operatorname{Tor}_1^A(M,N)\otimes_AA'
\to\operatorname{Tor}_1^{A'}(M',N')
\]
(cf.\ \eref{IV.11.2.1.5}) é \emph{sobrejetivo}.
\end{lemma}

\begin{proof}
Conservemos as notações de \sref{IV.11.2.3}; a exatidão à direita do produto tensorial mostra que a sequência $L'_1\xrightarrow{f'_0}L'_0\xrightarrow{\varepsilon'}M'\to0$ é exata; como $L'_0$ e $L'_1$ são $A'$-módulos livres, \emph{podemos supor que foram tomados} $L'_0=L''_0$ e $L'_1=L''_1$, sendo $u_0$ e $u_1$ as \emph{aplicações identidade} e $f''_0=f'_0$. Como $R=\operatorname{Im}(f_0)$ e $R''=\operatorname{Im}(f''_0)=\operatorname{Im}(f'_0)$, o homomorfismo $v$ é sobrejetivo e, portanto, também o é $v\otimes1$. A demonstração estará concluída se provarmos que a primeira linha de \eref{IV.11.2.3.3} é \emph{exata}, pois $v\otimes1$ é sobrejetivo e $u_0\otimes1$ é injetivo (Bourbaki, \emph{Alg. comm.}, cap.~I, \S~1, número~4, cor.~2 da prop.~2). Utilizemos para isso

\oldpage[IV-3]{123}

a hipótese de que $M\otimes_AN$ é um $C$-módulo plano. Pondo $Q=\operatorname{Ker}(\varepsilon\otimes1)=\operatorname{Im}(j\otimes1)$ na sequência exata \eref{IV.11.2.3.1}, temos as duas sequências exatas $0\to Q\to L_0\otimes_AN\to M\otimes_AN\to0$ e $T\to R\otimes_AN\to Q\to0$, cujos homomorfismos são homomorfismos de $C$-módulos; usando a hipótese de planicidade, deduz-se, por $(\mathbf 0_{\mathrm I},\,6.1.2)$, a sequência exata
\[
0\to Q\otimes_CC'\to(L_0\otimes_AN)\otimes_CC'
\to(M\otimes_AN)\otimes_CC'\to0
\]
e, por outro lado, como o produto tensorial é exato à direita, temos a sequência exata
\[
T\otimes_CC'\to(R\otimes_AN)\otimes_CC'\to Q\otimes_CC'\to0,
\]
e daí, finalmente, a sequência exata
\[
T\otimes_CC'\to(R\otimes_AN)\otimes_CC'\to
(L_0\otimes_AN)\otimes_CC'\to(M\otimes_AN)\otimes_CC'\to0.
\]
Mas, por definição, para todo $C$-módulo $P$, temos $P\otimes_CC'=P\otimes_AA'$, daí a conclusão.
\end{proof}

\begin{lemma}[11.2.5]
\label{IV.11.2.5}
Sejam $A$ um anel, $\mathfrak J$ um ideal de $A$, $B$ uma $A$-álgebra, $M$ um $B$-módulo e $A\to A'$ um homomorfismo de anéis. Ponhamos $\mathfrak J'=\mathfrak JA'$, $B'=B\otimes_AA'$ e $M'=M\otimes_AA'=M\otimes_BB'$. Seja $\mathfrak p'$ um ideal primo de $B'$ que contém $\mathfrak J'B'$. Suponhamos verificada uma das hipóteses seguintes:
\begin{enumerate}[label=\emph{\alph*)}]
\item $\mathfrak J$ é nilpotente.
\item $A'$ é noetheriano, $B'$ é uma $A'$-álgebra de tipo finito e $M'$ é um $B'$-módulo de tipo finito.
\end{enumerate}
Nessas condições, suponhamos verificadas as duas propriedades seguintes:
\begin{enumerate}[label=\emph{(\roman*)}]
\item $M\otimes_A(A/\mathfrak J)$ é um $(A/\mathfrak J)$-módulo plano.
\item O homomorfismo canônico composto
\[
\operatorname{Tor}_1^A(M,A/\mathfrak J)\xrightarrow{\psi_1}
\operatorname{Tor}_1^{A'}(M',A'/\mathfrak J')\xrightarrow{\theta}
(\operatorname{Tor}_1^{A'}(M',A'/\mathfrak J'))_{\mathfrak p'}
\]
(em que $\psi_1$ é o homomorfismo \eref{IV.11.2.1.1} e $\theta$ é o homomorfismo canônico de um $B'$-módulo em seu localizado em $\mathfrak p'$) é nulo.
\end{enumerate}
Então $M'_{\mathfrak p'}$ é um $A'$-módulo plano.
\end{lemma}

\begin{proof}
Notemos que, na hipótese b), $M'_{\mathfrak p'}$ é um $B'_{\mathfrak p'}$-módulo de tipo finito, $B'_{\mathfrak p'}$ é uma $A'$-álgebra noetheriana e $\mathfrak J'B'_{\mathfrak p'}$ está contido no radical de Jacobson $\mathfrak p'B'_{\mathfrak p'}$ de $B'_{\mathfrak p'}$; na hipótese a), $\mathfrak J'$ é nilpotente. Podemos, portanto, aplicar o critério de planicidade $(\mathbf 0_{\mathrm{III}},\,10.2.1)$ ou $(\mathbf 0_{\mathrm{III}},\,10.2.2)$, conforme se verifique a) ou b). Em primeiro lugar,
\[
\begin{aligned}
M'_{\mathfrak p'}\otimes_{A'}(A'/\mathfrak J')
&=(M'\otimes_{A'}(A'/\mathfrak J'))_{\mathfrak p'}\\
&=((M\otimes_A(A/\mathfrak J))\otimes_{A/\mathfrak J}
(A'/\mathfrak J'))_{\mathfrak p'}.
\end{aligned}
\]
A hipótese (i) implica, portanto, que $M'_{\mathfrak p'}\otimes_{A'}(A'/\mathfrak J')$ é um $(A'/\mathfrak J')$-módulo plano, tendo em conta $(\mathbf 0_{\mathrm I},\,6.2.1\text{ e }6.3.2)$. Resta, pois, mostrar que $\operatorname{Tor}_1^{A'}(M'_{\mathfrak p'},A'/\mathfrak J')=0$; ora, esse $B'_{\mathfrak p'}$-módulo é igual a $(\operatorname{Tor}_1^{A'}(M',A'/\mathfrak J'))_{\mathfrak p'}$, em virtude da planicidade de $B'_{\mathfrak p'}$ sobre $B'$. Pela hipótese (ii), o homomorfismo composto $\theta\circ\varphi_1:\operatorname{Tor}_1^A(M,A/\mathfrak J)\otimes_AA'\to(\operatorname{Tor}_1^{A'}(M',A'/\mathfrak J'))_{\mathfrak p'}$ é nulo; por outro lado, \sref{IV.11.2.4}, aplicado com $C=A/\mathfrak J$ e $N=C$, mostra, tendo em conta (i), que $\varphi_1$ é \emph{sobrejetivo} (pois $A'/\mathfrak J'=C\otimes_AA'$). Logo, $\theta$ é \emph{nulo}; como a imagem por $\theta$ de $\operatorname{Tor}_1^{A'}(M',A'/\mathfrak J')$ gera o $B'_{\mathfrak p'}$-módulo $(\operatorname{Tor}_1^{A'}(M',A'/\mathfrak J'))_{\mathfrak p'}$, este é nulo.
\end{proof}

\begin{theorem}[11.2.6]
\label{IV.11.2.6}
Com as notações de (8.5.1) e (8.8.1), suponhamos que $S_\alpha$ é quase-compacto e que $X_\alpha$ e $Y_\alpha$ são de apresentação finita sobre $S_\alpha$; sejam $f_\alpha:X_\alpha\to Y_\alpha$ um $S_\alpha$-morfismo e $\mathcal F_\alpha$ um $\mathcal O_{X_\alpha}$-Módulo quase-coerente de apresentação finita.

\oldpage[IV-3]{124}

\begin{enumerate}[label=\emph{(\roman*)}]
\item Seja $x$ um ponto de $X$ e $x_\lambda$ sua projeção canônica em $X_\lambda$. Para que $\mathcal F$ seja $f$-plano no ponto $x$, é necessário e suficiente que exista $\lambda\geq\alpha$ tal que $\mathcal F_\lambda$ seja $f_\lambda$-plano no ponto $x_\lambda$.
\item Para que $\mathcal F$ seja $f$-plano, é necessário e suficiente que exista $\lambda\geq\alpha$ tal que $\mathcal F_\lambda$ seja $f_\lambda$-plano.
\end{enumerate}
\end{theorem}

\begin{proof}
Podemos supor que $S_\alpha=S_0$; como $Y_0$ é de apresentação finita sobre $S_0$, $Y_0$ é quase-compacto, e $f_0:X_0\to Y_0$ é um morfismo de apresentação finita (1.6.2, (v)); portanto, podemos também nos restringir ao caso em que $S_0=Y_0$. Além disso, pela quase-compacidade de $S_0$ e porque o conjunto de índices $L$ é filtrante, podemos nos restringir ao caso em que $S_0=\operatorname{Spec}(A_0)$ é afim. Como $X_0$ também é quase-compacto, o mesmo raciocínio permite supor que $X_0=\operatorname{Spec}(B_0)$ é afim. Temos então $\mathcal F_0=\widetilde M_0$, onde $M_0$ é um $B_0$-módulo de apresentação finita, e o enunciado \sref{IV.11.2.6} equivale, nesse caso, ao seguinte (levando em conta $(\mathbf 0_{\mathrm I},\,6.3.1)$):

\begin{corollary}[11.2.6.1]
\label{IV.11.2.6.1}
Sejam $A_0$ um anel, $(A_\lambda)_{\lambda\in L}$ um sistema indutivo filtrante de $A_0$-álgebras, $B_0$ uma $A_0$-álgebra de apresentação finita e $M_0$ um $B_0$-módulo de apresentação finita. Ponhamos
\[
\begin{gathered}
B_\lambda=B_0\otimes_{A_0}A_\lambda,\qquad
M_\lambda=M_0\otimes_{A_0}A_\lambda=M_0\otimes_{B_0}B_\lambda,\\
A=\varinjlim_\lambda A_\lambda,\qquad
B=\varinjlim_\lambda B_\lambda=B_0\otimes_{A_0}A,\\
M=\varinjlim_\lambda M_\lambda
=M_0\otimes_{A_0}A=M_0\otimes_{B_0}B.
\end{gathered}
\]
\begin{enumerate}[label=\emph{(\roman*)}]
\item Seja $\mathfrak p$ um ideal primo de $B$ e, para cada $\lambda$, seja $\mathfrak p_\lambda$ sua imagem inversa em $B_\lambda$. Para que $M_{\mathfrak p}$ seja um $A$-módulo plano, é necessário e suficiente que exista $\lambda$ tal que $(M_\lambda)_{\mathfrak p_\lambda}$ seja um $A_\lambda$-módulo plano.
\item Para que $M$ seja um $A$-módulo plano, é necessário e suficiente que exista $\lambda$ tal que $M_\lambda$ seja um $A_\lambda$-módulo plano.
\end{enumerate}
\end{corollary}

Basta demonstrar que as condições são necessárias (2.1.4). Procederemos em várias etapas.

\medskip
\noindent\textbf{I)} \emph{Redução ao caso em que os $A_\lambda$ são noetherianos.} Em virtude de (8.9.1), existem um subanel $A'_0$ de $A_0$ que é uma $\mathbf Z$-álgebra de tipo finito, uma $A'_0$-álgebra de tipo finito $B'_0$ e um $B'_0$-módulo de tipo finito $M'_0$ tais que
\[
B_0=B'_0\otimes_{A'_0}A_0,\qquad
M_0=M'_0\otimes_{A'_0}A_0.
\]
Como $B_\lambda=B'_0\otimes_{A'_0}A_\lambda$ e $M_\lambda=M'_0\otimes_{A'_0}A_\lambda$, podemos substituir $A_0$, $B_0$ e $M_0$ por $A'_0$, $B'_0$ e $M'_0$ no enunciado de \sref{IV.11.2.6.1}, considerando os $A_\lambda$ como $A'_0$-álgebras. Podemos, portanto, supor desde já que $A_0$ é noetheriano.

Seja $H$ o conjunto dos pares $(\lambda,C_\lambda)$, em que $C_\lambda$ é uma sub-$A_0$-álgebra de \emph{tipo finito} de $A_\lambda$. Ordenemos $H$ pondo $(\lambda,C_\lambda)\leq(\mu,D_\mu)$ se $\lambda\leq\mu$ e se o homomorfismo $\varphi_{\mu\lambda}:A_\lambda\to A_\mu$ satisfaz $\varphi_{\mu\lambda}(C_\lambda)\subset D_\mu$. Com essa ordem, $H$ é filtrante crescente: se $(\lambda,C_\lambda)$ e $(\mu,D_\mu)$ são dois elementos quaisquer de $H$, podemos majorá-los por um $(\nu,E_\nu)$ escolhendo $\nu\geq\lambda$ e $\nu\geq\mu$ em $L$ e tomando $E_\nu$ como a sub-$A_0$-álgebra de $A_\nu$ gerada por $\varphi_{\nu\lambda}(C_\lambda)$ e $\varphi_{\nu\mu}(D_\mu)$.

Para um elemento $\xi=(\lambda,C_\lambda)$ de $H$, ponhamos $A_\xi=C_\lambda$; se $\xi\leq\eta=(\mu,D_\mu)$, então $\varphi_{\mu\lambda}(C_\lambda)\subset D_\mu$, e $\varphi_{\eta\xi}:A_\xi\to A_\eta$ é a restrição de $\varphi_{\mu\lambda}$ a $C_\lambda$, considerada como homomorfismo em $D_\mu$. É claro que se obtém assim um sistema indutivo filtrante de $A_0$-álgebras. Ponhamos $B_\xi=B_0\otimes_{A_0}A_\xi$ e $M_\xi=M_0\otimes_{A_0}A_\xi$; desta vez os $A_\xi$ são noetherianos. Além disso, a fórmula do duplo limite indutivo (Bourbaki, \emph{Alg.}, cap.~II, 3.ª ed., \S~6, n.º~4, prop.~7) demonstra que ainda temos
\[
\varinjlim_H A_\xi=A,\qquad
\varinjlim_H B_\xi=B,\qquad
\varinjlim_H M_\xi=M.
\]

Suponhamos demonstrado \sref{IV.11.2.6.1} para o sistema indutivo $(A_\xi)_{\xi\in H}$, e seja $\mathfrak p$ um ideal primo de $B$ tal que $M_{\mathfrak p}$ seja um $A$-módulo plano.

\oldpage[IV-3]{125}

Existe então $\xi\in H$ tal que, se $\mathfrak p_\xi$ é a imagem inversa de $\mathfrak p$ em $B_\xi$, o módulo $(M_\xi)_{\mathfrak p_\xi}$ é plano sobre $A_\xi$. Seja $\xi=(\lambda,C_\lambda)$; a inclusão $A_\xi=C_\lambda\to A_\lambda$ induz um homomorfismo $B_\xi=B_0\otimes_{A_0}C_\lambda\to B_\lambda=B_0\otimes_{A_0}A_\lambda$. Se $\mathfrak p_\lambda$ é a imagem inversa de $\mathfrak p$ em $B_\lambda$, então $\mathfrak p_\xi$ é a imagem inversa de $\mathfrak p_\lambda$ em $B_\xi$; consequentemente,
\[
(M_\lambda)_{\mathfrak p_\lambda}
=(M_\xi\otimes_{C_\lambda}A_\lambda)_{\mathfrak p_\lambda}
=((M_\xi)_{\mathfrak p_\xi}\otimes_{C_\lambda}A_\lambda)_{\mathfrak p_\lambda}.
\]
Logo $(M_\lambda)_{\mathfrak p_\lambda}$ é um $A_\lambda$-módulo plano $(\mathbf 0_{\mathrm I},\,6.2.1\text{ e }6.3.2)$. O caso (ii) do enunciado é tratado da mesma maneira. Podemos, portanto, supor daqui em diante que os $A_\lambda$ são noetherianos para $\lambda\in L$, embora o próprio $A$ não o seja necessariamente.

\medskip
\noindent\textbf{II)} \emph{Redução do enunciado global (ii) ao enunciado pontual (i).} Suponhamos que $\mathcal F$ seja $f$-plano. Para cada $\lambda$, seja $U_\lambda$ o conjunto dos pontos $x_\lambda\in X_\lambda$ em que $\mathcal F_\lambda$ é $f_\lambda$-plano. Sabemos que $U_\lambda$ é \emph{aberto} em $X_\lambda$, pois $S_\lambda$ é noetheriano e $f_\lambda$ é de tipo finito (11.1.1); seja $V_\lambda=v_\lambda^{-1}(U_\lambda)$ sua imagem inversa em $X$. Por hipótese, para todo $x\in X$ existe $\lambda$ tal que $\mathcal F_\lambda$ é $f_\lambda$-plano na projeção $x_\lambda$ de $x$ em $X_\lambda$; isto significa que $x\in V_\lambda$ para algum $\lambda$, ou seja, $X$ é a união dos $V_\lambda$. Além disso, por (2.1.4), se $\lambda\leq\mu$ temos $V_\lambda\subset V_\mu$. Como $X$ é quase-compacto, existe então um índice $\mu$ tal que $X=V_\mu$. Como os $X_\lambda$ são quase-compactos, decorre de (8.3.4) que existe $\nu\geq\mu$ tal que $v_{\mu\nu}^{-1}(U_\mu)=X_\nu$; por (2.1.4), isso implica que $\mathcal F_\nu$ é $f_\nu$-plano.

\medskip
\noindent\textbf{III)} \emph{Conclusão da demonstração.} Resta demonstrar (i), supondo $S_0$ afim e noetheriano. Se $y_0$ é a projeção de $x$ em $S_0$, podemos, por (2.1.4) e $(\mathbf I,\,3.6.5)$, substituir $S_0$ por $\operatorname{Spec}(\mathcal O_{S_0,y_0})$. Em outras palavras, podemos nos restringir ao caso em que $A_0$ é um anel \emph{local} noetheriano de ideal maximal $\mathfrak m_0=\mathfrak j_{y_0}$. Por definição, $\mathfrak m_0$ é a imagem inversa do ideal primo $\mathfrak p=\mathfrak j_x$ de $B$, e $M_{\mathfrak p}$ é, por hipótese, um $A$-módulo plano. Temos, portanto,
\[
\operatorname{Tor}_1^A(A/\mathfrak m_0A,M_{\mathfrak p})=0.
\]
Como os módulos $\operatorname{Tor}_1^A(A/\mathfrak m_0A,M)$ são $B$-módulos e $B_{\mathfrak p}$ é plano sobre $B$, isso também se escreve
\[
(\operatorname{Tor}_1^A(A/\mathfrak m_0A,M))_{\mathfrak p}=0.
\]
Agora, $\operatorname{Tor}_1^{A_0}(A_0/\mathfrak m_0,M_0)$ é um $B_0$-módulo \emph{de tipo finito}. Com efeito, podemos calculá-lo tomando uma resolução de $A_0/\mathfrak m_0$ por $A_0$-módulos livres \emph{de tipo finito}, pois $A_0$ é noetheriano, e tensorizando-a por $M_0$; obtemos assim $B_0$-módulos de tipo finito. Como $B_0$ é noetheriano, a homologia do complexo obtido também é formada por $B_0$-módulos de tipo finito. Sejam $(t_i^0)_{1\leq i\leq m}$ geradores do $B_0$-módulo $\operatorname{Tor}_1^{A_0}(A_0/\mathfrak m_0,M_0)$, e seja $t_i$ a imagem canônica de $t_i^0$ em $\operatorname{Tor}_1^A(A/\mathfrak m_0A,M)$ (11.2.1.1). A hipótese implica que existe $h\in B\setminus\mathfrak p$ tal que $ht_i=0$ para $1\leq i\leq m$.

Ora, por (11.2.2.1),
\[
\operatorname{Tor}_1^A(A/\mathfrak m_0A,M)
=\varinjlim_\lambda\operatorname{Tor}_1^{A_\lambda}
(A_\lambda/\mathfrak m_0A_\lambda,M_\lambda).
\]
Existe, portanto, $\lambda$ tal que, se $t_i^\lambda$ é a imagem de $t_i^0$ em $\operatorname{Tor}_1^{A_\lambda}(A_\lambda/\mathfrak m_0A_\lambda,M_\lambda)$, há $h_\lambda\in B_\lambda$ de imagem $h$ com $h_\lambda t_i^\lambda=0$ para $1\leq i\leq m$. Seja $\mathfrak p_\lambda=\mathfrak j_{x_\lambda}$ o ideal primo de $B_\lambda$ que é a imagem inversa de $\mathfrak p$; como $h_\lambda\in B_\lambda\setminus\mathfrak p_\lambda$, as imagens canônicas de todos os $t_i^0$ em
\[
(\operatorname{Tor}_1^{A_\lambda}
(A_\lambda/\mathfrak m_0A_\lambda,M_\lambda))_{\mathfrak p_\lambda}
\]
são nulas. Logo, o homomorfismo $\operatorname{Tor}_1^{A_0}(A_0/\mathfrak m_0,M_0)\to(\operatorname{Tor}_1^{A_\lambda}(A_\lambda/\mathfrak m_0A_\lambda,M_\lambda))_{\mathfrak p_\lambda}$ é nulo. As condições do Lema \sref{IV.11.2.5} estão, portanto, satisfeitas, pois $A_0/\mathfrak m_0$ é um corpo; daí $(M_\lambda)_{\mathfrak p_\lambda}$ ser um $A_\lambda$-módulo plano, o que conclui a demonstração do teorema.
\end{proof}

\begin{corollary}[11.2.7]
\label{IV.11.2.7}
Sejam $S=\operatorname{Spec}(A)$ um esquema afim, $f:X\to S$ um morfismo, $\mathcal F$ um $\mathcal O_X$-Módulo quase-coerente e $x$ um ponto de $X$. As condições seguintes são equivalentes:
\begin{enumerate}[label=\emph{\alph*)}]
\item $f$ é um morfismo de apresentação finita, $\mathcal F$ é um $\mathcal O_X$-Módulo de apresentação finita e $\mathcal F$ é $f$-plano no ponto $x$ (resp.\ $f$-plano).

\oldpage[IV-3]{126}

\item Existe um esquema afim noetheriano $S_0=\operatorname{Spec}(A_0)$, um morfismo de tipo finito $f_0:X_0\to S_0$, um $\mathcal O_{X_0}$-Módulo coerente $\mathcal F_0$ e um morfismo $S\to S_0$ tais que o $S$-pré-esquema $X_0\otimes_{S_0}S$ é $S$-isomorfo a $X$ e que, identificando $X$ com $X_0\otimes_{S_0}S$, $\mathcal F$ é isomorfo a $\mathcal F_0\otimes_{\mathcal O_{S_0}}\mathcal O_S$, sendo $\mathcal F_0$ $f_0$-plano no ponto $x_0$, projeção de $x$ em $X_0$ (resp.\ $f_0$-plano).
\item As condições de b) são satisfeitas e, além disso, $A_0$ é uma sub-$\mathbf Z$-álgebra de tipo finito de $A$, correspondendo o morfismo $S\to S_0$ à inclusão canônica $A_0\to A$.
\end{enumerate}
\end{corollary}

\begin{proof}
É claro que c) implica b) e que b) implica a) em virtude de (2.1.4). Por outro lado, podemos considerar $A$ como limite indutivo filtrante de suas sub-$\mathbf Z$-álgebras de tipo finito; por (8.9.1), existem uma tal subálgebra $A_0$ e um morfismo de tipo finito $f_0:X_0\to S_0$ tais que $X$ é $S$-isomorfo a $X_0\times_{S_0}S$ e $\mathcal F$ é isomorfo a $\mathcal F_0\otimes_{\mathcal O_{S_0}}\mathcal O_S$. Podemos, portanto, escrever $X=\varprojlim_\lambda X_\lambda$, onde $X_\lambda=X_0\otimes_{A_0}A_\lambda$ e $A_\lambda$ percorre o conjunto das sub-$\mathbf Z$-álgebras de tipo finito de $A$ que contêm $A_0$. Além disso, $\mathcal F=v_\lambda^*(\mathcal F_\lambda)$, em que $v_\lambda:X\to X_\lambda$ é a projeção canônica e $\mathcal F_\lambda=\mathcal F_0\otimes_{A_0}A_\lambda$. Decorre de \sref{IV.11.2.6} que existe $\lambda$ tal que $\mathcal F_\lambda$ é $f_\lambda$-plano no ponto $x_\lambda=v_\lambda(x)$ (resp.\ $f_\lambda$-plano); então o subanel $A_\lambda$ de $A$ satisfaz as condições de c).
\end{proof}

\begin{proposition}[11.2.8]
\label{IV.11.2.8}
Sejam $g:X\to S$ e $h:Y\to S$ dois morfismos de apresentação finita, $f:X\to Y$ um $S$-morfismo e $\mathcal F$ um $\mathcal O_X$-Módulo quase-coerente de apresentação finita. Para cada $s\in S$, sejam $X_s=g^{-1}(s)=X\times_S\operatorname{Spec}(k(s))$, $Y_s=h^{-1}(s)=Y\times_S\operatorname{Spec}(k(s))$, $f_s=f\times1:X_s\to Y_s$ e $\mathcal F_s=\mathcal F\otimes_{\mathcal O_S}k(s)$. Então o conjunto $E$ dos $s\in S$ para os quais $\mathcal F_s$ é $f_s$-plano é localmente construtível.
\end{proposition}

\begin{proof}
A propriedade cuja construtibilidade queremos demonstrar satisfaz a condição (9.2.1, (i)), em virtude de (2.2.13) e (2.5.1). Tendo em conta (9.2.3), podemos nos restringir ao caso em que $S$ é afim, noetheriano e integral, com ponto genérico $\eta$, e basta provar que $E$ ou $S-E$ é uma vizinhança de $\eta$ em $S$. Se $\eta\in S-E$, isto decorre imediatamente do Lema (9.4.7.1). Se, ao contrário, $\eta\in E$, decorre primeiro de \sref{IV.11.2.6}, aplicado pelo método de (8.1.2, a)), que existe uma vizinhança aberta $U$ de $\eta$ em $S$ tal que $\mathcal F|_{g^{-1}(U)}$ é $h^{-1}(U)$-plano. \emph{A fortiori}, por (2.1.4), $\mathcal F_s$ é $f_s$-plano para todo $s\in U$, o que conclui a demonstração.
\end{proof}

O teorema seguinte generaliza (11.2.6.1, (ii)).

\begin{theorem}[11.2.9]
\label{IV.11.2.9}
(Raynaud). Sejam $A_0$ um anel, $(A_\lambda)_{\lambda\in L}$ um sistema indutivo filtrante de $A_0$-álgebras, $B_0$ uma $A_0$-álgebra de apresentação finita, $\mathfrak J_0$ um ideal de tipo finito de $B_0$ e $M_0$ um $B_0$-módulo de apresentação finita. Ponhamos
\[
\begin{gathered}
B_\lambda=B_0\otimes_{A_0}A_\lambda,\qquad
\mathfrak J_\lambda=\mathfrak J_0B_\lambda,\qquad
M_\lambda=M_0\otimes_{A_0}A_\lambda=M_0\otimes_{B_0}B_\lambda,\\
A=\varinjlim_\lambda A_\lambda,\qquad
B=\varinjlim_\lambda B_\lambda=B_0\otimes_{A_0}A,\qquad
\mathfrak J=\mathfrak J_0B,\\
M=\varinjlim_\lambda M_\lambda
=M_0\otimes_{A_0}A=M_0\otimes_{B_0}B.
\end{gathered}
\]
Para que $\operatorname{gr}_{\mathfrak J}^*(M)$ seja um $A$-módulo plano, é necessário e suficiente que exista $\lambda$ tal que $\operatorname{gr}_{\mathfrak J_\lambda}^*(M_\lambda)$ seja um $A_\lambda$-módulo plano; os homomorfismos canônicos
\begin{equation}
\left\{
\begin{aligned}
\operatorname{gr}_{\mathfrak J_\lambda}^*(M_\lambda)
\otimes_{A_\lambda}A&\longrightarrow
\operatorname{gr}_{\mathfrak J}^*(M),\\
\operatorname{gr}_{\mathfrak J_\lambda}^*(M_\lambda)
\otimes_{A_\lambda}A_\mu&\longrightarrow
\operatorname{gr}_{\mathfrak J_\mu}^*(M_\mu)
\end{aligned}
\right.
\qquad\text{para }\lambda\leq\mu
\label{IV.11.2.9.1}
\tag{11.2.9.1}
\end{equation}
são então bijetivos, e $\operatorname{gr}_{\mathfrak J_\mu}^*(M_\mu)$ é um $A_\mu$-módulo plano para $\lambda\leq\mu$.\footnote{O impresso francês traz $A_\lambda$ neste último termo; adotamos $A_\mu$ conforme a mudança de base de (11.2.9.2).}
\end{theorem}

\begin{proof}
Mostremos primeiro que as condições são suficientes, o que equivale a demonstrar a bijetividade dos homomorfismos (11.2.9.1). Isto decorre do lema seguinte:

\oldpage[IV-3]{127}

\begin{lemma}[11.2.9.2]
\label{IV.11.2.9.2}
Sejam $A$ um anel, $B$ uma $A$-álgebra, $M$ um $B$-módulo e $\mathfrak J$ um ideal de $B$. Seja $A'$ uma $A$-álgebra; ponhamos $B'=B\otimes_AA'$, $M'=M\otimes_AA'=M\otimes_BB'$ e $\mathfrak J'=\mathfrak JB'$. Se $\operatorname{gr}_{\mathfrak J}^*(M)$ é um $A$-módulo plano, o homomorfismo canônico
\[
(\operatorname{gr}_{\mathfrak J}^*(M))\otimes_AA'
\longrightarrow\operatorname{gr}_{\mathfrak J'}^*(M')
\]
é \emph{bijetivo}.
\end{lemma}

Com efeito, por indução em $k$, a hipótese de que os módulos $\mathfrak J^kM/\mathfrak J^{k+1}M$ são planos sobre $A$ para $k\geq0$ implica primeiro, por $(\mathbf 0_{\mathrm I},\,6.1.2)$, que $M/\mathfrak J^{k+1}M$ é um $A$-módulo plano. Além disso, por $(\mathbf 0_{\mathrm I},\,6.1.2)$, a sequência
\[
0\longrightarrow(\mathfrak J^{k+1}M)\otimes_AA'
\longrightarrow M\otimes_AA'
\longrightarrow(M/\mathfrak J^{k+1}M)\otimes_AA'
\longrightarrow0
\]
é então exata; em outras palavras, $(\mathfrak J^{k+1}M)\otimes_AA'$ identifica-se com sua imagem canônica $(\mathfrak J')^{k+1}M'$ em $M'=M\otimes_AA'$. Por outro lado, ainda em virtude de $(\mathbf 0_{\mathrm I},\,6.1.2)$, a sequência
\[
0\longrightarrow(\mathfrak J^{k+1}M)\otimes_AA'
\longrightarrow(\mathfrak J^kM)\otimes_AA'
\longrightarrow(\mathfrak J^kM/\mathfrak J^{k+1}M)\otimes_AA'
\longrightarrow0
\]
é exata, o que demonstra o lema.

Para demonstrar a necessidade das condições de \sref{IV.11.2.9}, procederemos em várias etapas.

\medskip
\noindent\textbf{I)} \emph{Redução ao caso em que os $A_\lambda$ são noetherianos.} Procedemos como na redução I) de \sref{IV.11.2.6.1}, conservando suas notações; basta começar por substituir $A'_0$ por uma sub-$A'_0$-álgebra de tipo finito $A''_0$ de $A_0$ tal que, pondo $B''_0=B'_0\otimes_{A'_0}A''_0$, exista em $B''_0$ um ideal de tipo finito $\mathfrak J''_0$ com $\mathfrak J''_0B_0=\mathfrak J_0$.

Para isso, consideremos as sub-$A'_0$-álgebras de tipo finito $A'_\alpha$ de $A_0$, que formam uma família filtrante; temos $B_0=\varinjlim_\alpha B'_\alpha$, onde $B'_\alpha=B'_0\otimes_{A'_0}A'_\alpha$. Existe, portanto, um índice $\beta$ tal que um sistema finito de geradores de $\mathfrak J_0$ é formado pelas imagens em $B_0$ de elementos de $B'_\beta$. Tomaremos então $A''_0=A'_\beta$, $B''_0=B'_\beta$ e, para $\mathfrak J''_0$, o ideal gerado por esses elementos. Podemos, portanto, supor que $A_0$ e, por conseguinte, também $B_0$ são noetherianos.

Definimos em seguida, como no lugar citado, o conjunto filtrante $H$ e os $A_\xi$, $B_\xi$, $M_\xi$ para $\xi\in H$; ponhamos também $\mathfrak J_\xi=\mathfrak J''_0B_\xi$ para todo $\xi\in H$. Suponhamos demonstrado que existe $\xi=(\lambda,C_\lambda)$ tal que $\operatorname{gr}_{\mathfrak J_\xi}^*(M_\xi)$ é um $C_\lambda$-módulo plano. Como $M_\lambda=M_\xi\otimes_{C_\lambda}A_\lambda$, decorre de \sref{IV.11.2.9.2} que $\operatorname{gr}_{\mathfrak J_\lambda}^*(M_\lambda)$ é um $A_\lambda$-módulo plano.

\medskip
\noindent\textbf{II)} \emph{Lemas preliminares.}

\begin{lemma}[11.2.9.3]
\label{IV.11.2.9.3}
Sejam $A$ um anel, $B=\bigoplus_{i\geq0}B_i$ uma $A$-álgebra graduada em graus não negativos e $M=\bigoplus_{i\in\mathbf Z}M_i$ um $B$-módulo graduado. Suponhamos que $B_0$ seja um anel local e que cada $M_i$ seja um $B_0$-módulo de tipo finito. Para que $M$ seja um $B$-módulo de tipo finito, é necessário e suficiente que, sendo $\mathfrak q$ a imagem inversa em $A$ do ideal maximal $\mathfrak m$ de $B_0$, o módulo $M\otimes_Ak(\mathfrak q)$ seja de tipo finito sobre $B\otimes_Ak(\mathfrak q)$.
\end{lemma}

A condição é evidentemente necessária; provemos que é suficiente. Como $B_0$ (e, por conseguinte, $B$) é uma $A_{\mathfrak q}$-álgebra, podemos substituir $A$ por $A_{\mathfrak q}$, isto é, supor que $A$ também é um anel local cujo ideal maximal é $\mathfrak q$. Por hipótese, existe um inteiro $N$ tal que o homomorfismo canônico
\[
\left(\bigoplus_{-N\leq i\leq N}(M_i\otimes_{B_0}B)\right)
\otimes_Ak(\mathfrak q)\longrightarrow M\otimes_Ak(\mathfrak q)
\]

\oldpage[IV-3]{128}

é sobrejetivo; isso significa também que, para todo inteiro $n$, o homomorfismo canônico de $k(\mathfrak q)$-módulos
\[
\left(\bigoplus_{-N\leq i\leq N}(M_i\otimes_{B_0}B_{n-i})\right)
\otimes_Ak(\mathfrak q)\longrightarrow M_n\otimes_Ak(\mathfrak q)
\]
é sobrejetivo. Ora, $M_n$ é um $B_0$-módulo de tipo finito, $B_0$ é um anel local e $\mathfrak qB_0\subset\mathfrak m$; portanto, o lema de Nakayama demonstra que cada um dos homomorfismos canônicos
\[
\bigoplus_{-N\leq i\leq N}(M_i\otimes_{B_0}B_{n-i})
\longrightarrow M_n
\]
é sobrejetivo, de onde segue a conclusão.

\begin{corollary}[11.2.9.4]
\label{IV.11.2.9.4}
Sob as hipóteses de \sref{IV.11.2.9.3}, suponhamos, além disso, que cada $B_i$ e cada $M_i$ seja um $B_0$-módulo de apresentação finita, e que $M$ seja um $A$-módulo plano. Para que $M$ seja um $B$-módulo de apresentação finita, é necessário e suficiente que $M\otimes_Ak(\mathfrak q)$ seja um módulo de apresentação finita sobre $B\otimes_Ak(\mathfrak q)$.
\end{corollary}

Em virtude de \sref{IV.11.2.9.3}, se a condição do enunciado for satisfeita, $M$ é um $B$-módulo de tipo finito; existe, portanto, um $B$-módulo livre \emph{graduado} de tipo finito $L$ (que tem, assim, uma base finita de elementos homogêneos) e um homomorfismo graduado sobrejetivo de grau $0$, $u:L\to M$. Seja $R=\ker(u)$, que é um $B$-módulo graduado. Existem, portanto, inteiros $m_j$ em número finito ($1\leq j\leq r$) tais que, para cada inteiro $i$, $R_i$ é o núcleo de um homomorfismo sobrejetivo
\[
\bigoplus_{1\leq j\leq r}B_{i+m_j}\longrightarrow M_i;
\]
conclui-se então, pela hipótese sobre os $M_i$ e os $B_i$, que $R_i$ é um $B_0$-módulo de tipo finito (Bourbaki, \emph{Alg. comm.}, cap.~I, \S~2, n\textsuperscript{o}~8, lema 9). Para provar que $R$ é um $B$-módulo de tipo finito, observemos que, em virtude da hipótese de planicidade e de $(\mathbf 0_{\mathrm I},\,6.1.2)$, a sequência
\[
0\longrightarrow R\otimes_Ak(\mathfrak q)
\longrightarrow L\otimes_Ak(\mathfrak q)
\longrightarrow M\otimes_Ak(\mathfrak q)\longrightarrow0
\]
é exata; a hipótese acarreta, portanto (Bourbaki, \emph{loc. cit.}), que $R\otimes_Ak(\mathfrak q)$ é um módulo de tipo finito sobre $B\otimes_Ak(\mathfrak q)$; basta, pois, aplicar a $R$ o lema \sref{IV.11.2.9.3}.

\medskip
\noindent\textbf{III)} \emph{Redução ao caso em que os homomorfismos de transição $\varphi_{\mu\lambda}:A_\lambda\to A_\mu$ (para $\lambda\leq\mu$) são injetivos.} Seja $A'_\lambda$ a imagem de $A_\lambda$ pelo homomorfismo canônico $\varphi_\lambda:A_\lambda\to A$; é claro que os $A'_\lambda$ formam um sistema indutivo de subanéis noetherianos de $A$, cujo limite indutivo é $A$, e que os homomorfismos de transição $A'_\lambda\to A'_\mu$ (para $\lambda\leq\mu$) são injetivos. Ponhamos
\[
B'_\lambda=B_0\otimes_{A_0}A'_\lambda,\qquad
\mathfrak J'_\lambda=\mathfrak J_0B'_\lambda,\qquad
M'_\lambda=M_0\otimes_{A_0}A'_\lambda
=M_\lambda\otimes_{A_\lambda}A'_\lambda;
\]
temos ainda $B=\varinjlim_\lambda B'_\lambda$ e $M=\varinjlim_\lambda M'_\lambda$. Suponhamos demonstrado que existe $\lambda$ tal que, para $\mu\geq\lambda$, $\operatorname{gr}_{\mathfrak J'_\mu}^*(M'_\mu)$ seja um $A'_\mu$-módulo plano. Seja $\mathfrak a_\lambda$ o núcleo do homomorfismo $\varphi_\lambda$; trata-se, portanto, de um ideal de tipo finito do anel noetheriano $A_\lambda$. Decorre da definição de limite indutivo que existe um índice $\mu\geq\lambda$ tal que $\varphi_{\mu\lambda}(\mathfrak a_\lambda)=0$; em outras palavras, o homomorfismo $\varphi_{\mu\lambda}$ se fatora como $A_\lambda\to A'_\lambda\to A_\mu$, e podemos escrever
\[
B_\mu=B'_\lambda\otimes_{A'_\lambda}A_\mu,\qquad
\mathfrak J_\mu=\mathfrak J'_\lambda B_\mu,\qquad
M_\mu=M'_\lambda\otimes_{A'_\lambda}A_\mu.
\]
Deduzimos, portanto, de \sref{IV.11.2.9.2} que $\operatorname{gr}_{\mathfrak J_\mu}^*(M_\mu)$ é um $A_\mu$-módulo plano.

\medskip
\noindent\textbf{IV)} \emph{Redução ao caso em que $B_0=A_0[T_1,\ldots,T_n]$ e $\operatorname{gr}_{\mathfrak J_0}^*(B_0)=B_0$.} Em virtude da hipótese, existe um sistema $(t_i)_{1\leq i\leq n}$ de geradores da $A_0$-álgebra de tipo finito $B_0$, tal que os $t_i$ para $i\leq m$ geram o ideal $\mathfrak J_0$ de $B_0$. Ponhamos $B'_0=A_0[T_1,\ldots,T_n]$,

\oldpage[IV-3]{129}

álgebra de polinômios (portanto noetheriana), e seja $\mathfrak J'_0$ o ideal de $B'_0$ gerado pelos $T_i$ de índice $i\leq m$; existe, então, um $A_0$-homomorfismo sobrejetivo $u_0:B'_0\to B_0$ que leva cada $T_i$ a $t_i$ ($1\leq i\leq n$), de modo que $u_0(\mathfrak J'_0)=\mathfrak J_0$. Esse homomorfismo permite considerar $M_0$ como um $B'_0$-módulo de tipo finito. Ponhamos
\[
B'_\lambda=B'_0\otimes_{A_0}A_\lambda
=A_\lambda[T_1,\ldots,T_n],\qquad
\mathfrak J'_\lambda=\mathfrak J'_0B'_\lambda,
\]
\[
B'=B'_0\otimes_{A_0}A=A[T_1,\ldots,T_n],\qquad
\mathfrak J'=\mathfrak J'_0B',
\]
de modo que $\mathfrak J'$ é o ideal de $B'$ gerado pelos $T_i$ de índice $i\leq m$. Além disso, é claro que, para todo inteiro $k\geq0$, temos $\mathfrak J'^kM=\mathfrak J^kM$ e $\mathfrak J'_\lambda{}^kM_\lambda=\mathfrak J_\lambda^kM_\lambda$ para todo $\lambda$; portanto,
\[
\operatorname{gr}_{\mathfrak J'}^*(M)
=\operatorname{gr}_{\mathfrak J}^*(M)
\quad\text{como $A$-módulos},\qquad
\operatorname{gr}_{\mathfrak J'_\lambda}^*(M_\lambda)
=\operatorname{gr}_{\mathfrak J_\lambda}^*(M_\lambda)
\quad\text{como $A_\lambda$-módulos}.
\]
Podemos, assim, substituir $B$, $B_\lambda$, $\mathfrak J$ e $\mathfrak J_\lambda$ por $B'$, $B'_\lambda$, $\mathfrak J'$ e $\mathfrak J'_\lambda$, respectivamente, na demonstração. Observemos também que, por construção, $\operatorname{gr}_{\mathfrak J'}^*(B')$ se identifica com $B'$ e $\operatorname{gr}_{\mathfrak J'_\lambda}^*(B'_\lambda)$ com $B'_\lambda$.

\medskip
\noindent\textbf{V)} \emph{Demonstração de que $\operatorname{gr}_{\mathfrak J}^*(M)$ é um módulo de apresentação finita sobre $\operatorname{gr}_{\mathfrak J}^*(B)$.} Supomos, portanto, daqui em diante, que $A_0$ e os $A_\lambda$ são noetherianos, que os homomorfismos de transição $A_\lambda\to A_\mu$ são injetivos, que $B_0=A_0[T_1,\ldots,T_n]$ e que $\mathfrak J_0$ é gerado pelos $T_i$ de índice $i\leq m$. Assim, o anel $C_0=\operatorname{gr}_{\mathfrak J_0}^*(B_0)$ se identifica com $B_0$, e $C=\operatorname{gr}_{\mathfrak J}^*(B)$ se identifica com $B$; de início, na argumentação que segue, utilizaremos apenas o fato de que $C\cong C_0\otimes_{A_0}A$.

Precisaremos da seguinte variante de (6.9.3):

\begin{lemma}[11.2.9.5]
\label{IV.11.2.9.5}
Sejam $A_0$ um anel noetheriano, $B_0$ uma $A_0$-álgebra de tipo finito, $\mathfrak J_0$ um ideal de $B_0$ e $M_0$ um $B_0$-módulo de tipo finito. Existe, então, uma sequência $(S_{0i})_{1\leq i\leq m}$ de subesquemas de $S_0=\operatorname{Spec}(A_0)$ com as seguintes propriedades:
\begin{enumerate}[label=\arabic{enumi}$^\circ$]
\item Os espaços subjacentes aos $S_{0i}$ são dois a dois disjuntos e formam uma cobertura de $S_0$.
\item Para cada $i$, o conjunto $S_{0i}$ é aberto em $\bigcup_{j\geq i}S_{0j}$.
\item Cada esquema $S_{0i}$ é afim.
\item Se $A_{0i}$ é o anel de $S_{0i}$ e se pomos $B_{0i}=B_0\otimes_{A_0}A_{0i}$ e $\mathfrak J_{0i}=\mathfrak J_0B_{0i}$, então $\operatorname{gr}_{\mathfrak J_{0i}}^*(M_0\otimes_{A_0}A_{0i})$ é um $A_{0i}$-módulo plano.
\end{enumerate}
\end{lemma}

Procedemos como em (6.9.3), por indução noetheriana, supondo o lema verdadeiro quando se substitui $A_0$ por $A'_0=A_0/\mathfrak a_0$, onde o ideal $\mathfrak a_0$ é tal que $V(\mathfrak a_0)\neq S_0$, e se substituem $B_0$, $\mathfrak J_0$ e $M_0$, respectivamente, por $B'_0=B_0\otimes_{A_0}A'_0$, $\mathfrak J'_0=\mathfrak J_0B'_0$ e $M_0\otimes_{A_0}A'_0$. Consideremos o nilradical $\mathfrak N_0$ de $A_0$; é evidentemente suficiente demonstrar o lema substituindo $A_0$ por $A_0/\mathfrak N_0$ e $B_0$, $\mathfrak J_0$, $M_0$ pelos objetos correspondentes, como acima. Em outras palavras, podemos supor que $A_0$ é \emph{reduzido}. Por outro lado, $\operatorname{gr}_{\mathfrak J_0}^*(B_0)$ é uma $A_0$-álgebra de tipo finito, e $\operatorname{gr}_{\mathfrak J_0}^*(M_0)$ é um módulo de tipo finito sobre $\operatorname{gr}_{\mathfrak J_0}^*(B_0)$. Pelo teorema da planicidade genérica (6.9.1), existe um aberto $D(t_0)\subset S_0$ tal que, pondo $A_{01}=(A_0)_{t_0}$, o módulo $\operatorname{gr}_{\mathfrak J_0}^*(M_0)\otimes_{A_0}A_{01}$ seja plano sobre $A_{01}$. Mas, como $A_{01}$ é um $A_0$-módulo plano, esse módulo se identifica com $\operatorname{gr}_{\mathfrak J_{01}}^*(M_0\otimes_{A_0}A_{01})$, onde $B_{01}=B_0\otimes_{A_0}A_{01}$ e $\mathfrak J_{01}=\mathfrak J_0B_{01}$. O complementar de $D(t_0)$ em $S_0$ é, então, da forma $V(\mathfrak a_0)$, e concluímos aplicando a hipótese de indução noetheriana.

Apliquemos esse lema à situação atual, conservando as mesmas notações; ponhamos $U_{0i}=\bigcup_{j\leq i}S_{0j}$, que é um aberto quase-compacto de $S_0$. Existe, portanto, uma família $(f_{ik})_{1\leq i\leq m,\,1\leq k\leq n}$ de elementos de $A_0$ tal que, para cada $i$, $U_{0i}$ é a união dos $D(f_{ik})$ ($1\leq k\leq n$).

\oldpage[IV-3]{130}

O $C_0$-módulo $\operatorname{gr}_{\mathfrak J_0}^*(M_0)$ é gerado por um número finito de elementos homogêneos, de graus não necessariamente iguais.\footnote{O impresso francês diz “de grau 1”, o que não vale em geral: para $B_0=k[t]$, $\mathfrak J_0=(t)$ e $M_0=k[t]/(t)$, o graduado de $M_0$ está concentrado no grau $0$. Conserva-se aqui a leitura impressa como testemunho; a graduação da base livre é ajustada aos graus dos geradores.} Dotando $C_0^r$ da graduação em que cada vetor da base canônica tem o grau do gerador correspondente, obtemos um epimorfismo de $C_0$-módulos graduados
\[
u_0:C_0^r\longrightarrow\operatorname{gr}_{\mathfrak J_0}^*(M_0).
\]
Como $C=C_0\otimes_{A_0}A$ e o homomorfismo canônico $\operatorname{gr}_{\mathfrak J_0}^*(M_0)\otimes_{A_0}A\to\operatorname{gr}_{\mathfrak J}^*(M)$ é sobrejetivo, deduzimos de $u_0$ um epimorfismo de $C$-módulos graduados, com a graduação induzida em $C^r$,
\[
u:C^r\xrightarrow{u_0\otimes1_A}
\operatorname{gr}_{\mathfrak J_0}^*(M_0)\otimes_{A_0}A
\longrightarrow\operatorname{gr}_{\mathfrak J}^*(M),
\]
e tudo se reduz a mostrar que o $C$-módulo graduado $Q=\ker(u)$ é \emph{de tipo finito}.

\begin{lemma}[11.2.9.6]
\label{IV.11.2.9.6}
Sob as hipóteses precedentes, seja $\mathfrak R_0$ um ideal de $A_0$; ponhamos
\[
A'_0=A_0/\mathfrak R_0,\qquad
B'_0=B_0\otimes_{A_0}A'_0=B_0/\mathfrak R_0B_0,\qquad
\mathfrak J'_0=\mathfrak J_0B'_0,
\]
e $\mathfrak R=\mathfrak R_0A$, $A'=A/\mathfrak R$; suponhamos, além disso, que $\operatorname{gr}_{\mathfrak J'_0}^*(M_0\otimes_{A_0}A'_0)$ seja um $A'_0$-módulo plano. Então existe um número finito de elementos $q_h$ ($1\leq h\leq s$) de $Q$ tais que, para todo ideal primo $\mathfrak p\supset\mathfrak RB$ de $B$, as imagens canônicas dos $q_h$ em $Q_{\mathfrak p}$ geram o $C_{\mathfrak p}$-módulo $Q_{\mathfrak p}$.
\end{lemma}

Suponhamos demonstrado esse lema e observemos que também podemos aplicá-lo substituindo $A_0$ por $(A_0)_{t_0}$ e $B_0$, $\mathfrak J_0$, $M_0$, $A_\lambda$, $A$ pelas respectivas localizações $(B_0)_{t_0}$, $(\mathfrak J_0)_{t_0}$, $(M_0)_{t_0}$, $(A_\lambda)_{t_0}$, $A_{t_0}$, para qualquer $t_0\in A_0$, pois $(A_0)_{t_0}$ é um $A_0$-módulo plano. Aplicaremos, então, o lema substituindo sucessivamente $A_0$ por cada anel $(A_0)_{f_{ik}}$ e $\mathfrak R_0$ pelo ideal $\mathfrak R_{ik}$ que define o subesquema fechado de $D(f_{ik})$ induzido por $S_{0i}$ no aberto $S_{0i}\cap D(f_{ik})$ de $S_{0i}$. Como a hipótese de planicidade de \sref{IV.11.2.9.6} se verifica em cada um desses casos por \sref{IV.11.2.9.5}, obtemos, para cada par $(i,k)$, uma família finita $(a_{ikh})_{1\leq h\leq l}$ de elementos de $Q_{f_{ik}}$ cujas imagens em $Q_{\mathfrak p_{ik}}$ geram esse $C_{\mathfrak p_{ik}}$-módulo, para todo ideal primo $\mathfrak p_{ik}$ de $B_{f_{ik}}$ que contém $\mathfrak R_{ik}B_{f_{ik}}$.

Podemos escrever, para um inteiro adequado $d>0$, $a_{ikh}=b_{ikh}/f_{ik}^d$ para todos os índices $i,k,h$, com $b_{ikh}\in Q$. Como os $S_{0i}$ cobrem $S_0$, a imagem em $S_0$ de qualquer ideal primo $\mathfrak p$ de $B$ pertence a algum $S_{0i}\cap D(f_{ik})$; logo, a imagem $\mathfrak p_{ik}$ de $\mathfrak p$ em $B_{f_{ik}}$ contém $\mathfrak R_{ik}B_{f_{ik}}$. As imagens dos $b_{ikh}$ ($1\leq h\leq l$) em $Q_{\mathfrak p_{ik}}=Q_{\mathfrak p}$ geram esse $C_{\mathfrak p}$-módulo. Concluímos, portanto, que os $b_{ikh}$ ($1\leq i\leq m$, $1\leq k\leq n$, $1\leq h\leq l$) geram o $C$-módulo $Q$ (Bourbaki, \emph{Alg. comm.}, cap.~II, \S~3, n\textsuperscript{o}~3, teorema 1).

Resta demonstrar o lema \sref{IV.11.2.9.6}. Ponhamos
\[
C'=C\otimes_AA'=C/\mathfrak RC,\qquad
B'=B\otimes_AA',\qquad
\mathfrak J'=\mathfrak JB',\qquad
M'=M\otimes_AA'.
\]
Por hipótese, $\operatorname{gr}_{\mathfrak J}^*(M)$ é um $A$-módulo plano; portanto, $\operatorname{gr}_{\mathfrak J'}^*(M')$ se identifica com $\operatorname{gr}_{\mathfrak J}^*(M)\otimes_AA'$ por \sref{IV.11.2.9.2}, e temos a sequência exata $(\mathbf 0_{\mathrm I},\,6.1.2)$
\begin{equation}
0\longrightarrow Q\otimes_AA'
\longrightarrow C'^{r}
\xrightarrow{u\otimes1_{A'}}
\operatorname{gr}_{\mathfrak J'}^*(M')
\longrightarrow0.
\label{IV.11.2.9.7}
\tag{11.2.9.7}
\end{equation}
Ponhamos $C'_0=C_0\otimes_{A_0}A'_0=C_0/\mathfrak R_0C_0$ e consideremos o epimorfismo de $C'_0$-módulos graduados
\[
u'_0:C'_0{}^{r}
\xrightarrow{u_0\otimes1_{A'_0}}
\operatorname{gr}_{\mathfrak J_0}^*(M_0)\otimes_{A_0}A'_0
\longrightarrow\operatorname{gr}_{\mathfrak J'_0}^*(M_0\otimes_{A_0}A'_0).
\]
Seja $Q_0=\ker(u'_0)$; como $\operatorname{gr}_{\mathfrak J'_0}^*(M_0\otimes_{A_0}A'_0)$ é um $A'_0$-módulo plano, seu produto tensorial sobre $A'_0$ por $A'$ se identifica com $\operatorname{gr}_{\mathfrak J'}^*(M')$ por \sref{IV.11.2.9.2}, e temos a sequência exata $(\mathbf 0_{\mathrm I},\,6.1.2)$
\begin{equation}
0\longrightarrow Q_0\otimes_{A'_0}A'
\longrightarrow C'^{r}
\xrightarrow{u\otimes1_{A'}}
\operatorname{gr}_{\mathfrak J'}^*(M')
\longrightarrow0.
\label{IV.11.2.9.8}
\tag{11.2.9.8}
\end{equation}

\oldpage[IV-3]{131}

Da comparação com (11.2.9.7) resulta, portanto, um isomorfismo de $C'$-módulos graduados
\[
Q\otimes_AA'\xrightarrow{\sim}Q_0\otimes_{A'_0}A'.
\]
Como $C'_0$ é noetheriano, $Q_0$ é um $C'_0$-módulo de tipo finito; portanto, $Q_0\otimes_{A'_0}A'$ é um $C'$-módulo graduado de tipo finito, e o mesmo vale para $Q\otimes_AA'$. Sejam $q_h$ ($1\leq h\leq s$) elementos homogêneos de $Q$ cujas imagens em $Q\otimes_AA'$ geram esse $C'$-módulo.

Consideremos agora um ideal primo qualquer $\mathfrak p\supset\mathfrak RB$ de $B$. Pela planicidade da localização, temos $C_{\mathfrak p}=\operatorname{gr}_{\mathfrak J_{\mathfrak p}}^*(B_{\mathfrak p})$, e $\operatorname{gr}_{\mathfrak J_{\mathfrak p}}^0(B_{\mathfrak p})=B_{\mathfrak p}/\mathfrak J_{\mathfrak p}$ é o anel nulo ou um anel local; para demonstrar o lema, basta evidentemente considerar o segundo caso. É claro, então, que cada $\operatorname{gr}_{\mathfrak J_{\mathfrak p}}^i(B_{\mathfrak p})$ é um módulo de apresentação finita sobre $B_{\mathfrak p}/\mathfrak J_{\mathfrak p}$. Por outro lado, para todo índice $i$,
\[
M/\mathfrak J^iM
\simeq(M_0/\mathfrak J_0^iM_0)\otimes_{A_0}A,
\]
e é, portanto, um $B$-módulo de apresentação finita. Por indução em $i$, a hipótese de que $\operatorname{gr}_{\mathfrak J}^*(M)$ é um $A$-módulo plano implica, por $(\mathbf 0_{\mathrm I},\,6.1.2)$, que $M/\mathfrak J^iM$ é um $A$-módulo plano. A aplicação de \sref{IV.11.2.6.1}, substituindo $M_0$ por $M_0/\mathfrak J_0^{k+1}M_0$ para $0\leq k\leq i$\footnote{Na p.~131 do impresso, a aplicação de (11.2.6.1) é feita a $M_0/\mathfrak J_0^kM_0$ para $k\leq i$, mas a conclusão requer os quocientes de expoente $k+1$ até $i+1$. Explicitamos essa aplicação aos expoentes $k+1$, para $0\leq k\leq i$; a leitura impressa fica preservada nesta nota.}, mostra que existe um índice $\lambda_i$ tal que, para $\mu\geq\lambda_i$, cada $M_\mu/\mathfrak J_\mu^{k+1}M_\mu$ é um $A_\mu$-módulo plano; consequentemente, por $(\mathbf 0_{\mathrm I},\,6.1.2)$, cada $\operatorname{gr}_{\mathfrak J_\mu}^k(M_\mu)$ também é plano sobre $A_\mu$ para $k\leq i$. Deduzimos de \sref{IV.11.2.9.2} que cada $\operatorname{gr}_{\mathfrak J}^i(M)$ é um módulo de apresentação finita sobre $B/\mathfrak J$ e, assim, $\operatorname{gr}_{\mathfrak J_{\mathfrak p}}^i(M_{\mathfrak p})$ é um módulo de apresentação finita sobre $B_{\mathfrak p}/\mathfrak J_{\mathfrak p}$.

Além disso, as imagens dos $q_h$ em
\[
Q_{\mathfrak p}\otimes_{B_{\mathfrak p}}k(\mathfrak p)
=(Q\otimes_AA')_{\mathfrak p}\otimes_{A'}k(\mathfrak p)
\]
geram esse módulo sobre $C_{\mathfrak p}\otimes_{B_{\mathfrak p}}k(\mathfrak p)$, pois
\[
C_{\mathfrak p}\otimes_{B_{\mathfrak p}}k(\mathfrak p)
=C'_{\mathfrak p}\otimes_{A'}k(\mathfrak p).
\]
Como temos a sequência exata
\[
Q_{\mathfrak p}\otimes_{B_{\mathfrak p}}k(\mathfrak p)
\longrightarrow C_{\mathfrak p}^{r}\otimes_{B_{\mathfrak p}}k(\mathfrak p)
\longrightarrow
\operatorname{gr}_{\mathfrak J_{\mathfrak p}}^*(M_{\mathfrak p})
\otimes_{B_{\mathfrak p}}k(\mathfrak p)\longrightarrow0,
\]
concluímos que $\operatorname{gr}_{\mathfrak J_{\mathfrak p}}^*(M_{\mathfrak p})\otimes_{B_{\mathfrak p}}k(\mathfrak p)$ é um módulo de apresentação finita sobre $C_{\mathfrak p}\otimes_{B_{\mathfrak p}}k(\mathfrak p)$. Aplicando \sref{IV.11.2.9.4}, com $A$, $B$ e $M$ substituídos, respectivamente, por $B_{\mathfrak p}$, $C_{\mathfrak p}$ e $\operatorname{gr}_{\mathfrak J_{\mathfrak p}}^*(M_{\mathfrak p})$, concluímos que $Q_{\mathfrak p}$ é um $C_{\mathfrak p}$-módulo de tipo finito. O lema de Nakayama (e o fato de que $C=B$) mostra então que as imagens dos $q_h$ em $Q_{\mathfrak p}$ geram esse $C_{\mathfrak p}$-módulo. Isso conclui a demonstração do lema \sref{IV.11.2.9.6} e prova também que $\operatorname{gr}_{\mathfrak J}^*(M)$ é um $C$-módulo de apresentação finita.

\medskip
\noindent\textbf{VI)} \emph{Fim da demonstração.} Para abreviar, ponhamos $C_\lambda=C_0\otimes_{A_0}A_\lambda$ (na verdade, igual a $B_\lambda$), $N_\lambda=\operatorname{gr}_{\mathfrak J_\lambda}^*(M_\lambda)$ e $N=\operatorname{gr}_{\mathfrak J}^*(M)$. Observemos primeiro que, para todo inteiro $k$, $\mathfrak J^kM$ se identifica com $\varinjlim_\lambda\mathfrak J_\lambda^kM_\lambda$ pela exatidão do funtor $\varinjlim$: a imagem de $\mathfrak J_\lambda^kM_\lambda$ em $M_\mu$ para $\lambda\leq\mu$ (respectivamente, em $M$) gera $\mathfrak J_\mu^kM_\mu$ (respectivamente, $\mathfrak J^kM$). Usando novamente a exatidão de $\varinjlim$, concluímos que $N$ se identifica canonicamente com $\varinjlim_\lambda N_\lambda$. Feita essa identificação, demonstraremos primeiro:
\begin{equation}
\text{\emph{Para $\lambda$ suficientemente grande, o homomorfismo canônico
$v_\lambda:N_\lambda\otimes_{A_\lambda}A\to N$ é bijetivo.}}
\label{IV.11.2.9.9}
\tag{11.2.9.9}
\end{equation}
Como os $C_\lambda$ são noetherianos e $N_\lambda$ é um $C_\lambda$-módulo de tipo finito, os $C$-módulos $N_\lambda\otimes_{A_\lambda}A$ são de apresentação finita e formam um sistema indutivo filtrante cujo limite indutivo se identifica canonicamente com $N$, pois os limites indutivos filtrantes comutam com os produtos tensoriais.
Além disso, os homomorfismos de transição
$v_{\mu\lambda}:N_\lambda\otimes_{A_\lambda}A
\to N_\mu\otimes_{A_\mu}A$

\oldpage[IV-3]{132}

para $\lambda\leq\mu$ e os homomorfismos $v_\lambda$ são \emph{sobrejetivos}.
Para um $\lambda$ fixado, consideremos os sub-$C$-módulos
$\ker(v_{\mu\lambda})=N'_\mu$ de $N_\lambda\otimes_{A_\lambda}A$; pela
definição de limite indutivo, eles formam uma família filtrante crescente
de sub-$C$-módulos de $\ker(v_\lambda)$, cuja reunião é
$\ker(v_\lambda)$. Mas vimos em V) que $N$ é um $C$-módulo de
apresentação finita; portanto, pelo lema 9 de Bourbaki
(\emph{Alg. comm.}, cap.~I, \S~2, n$^\circ$~8),
$\ker(v_\lambda)$ é um $C$-módulo de tipo finito. Existe, assim, um índice
$\mu\geq\lambda$ para o qual $N'_\mu=\ker(v_\lambda)$; como
$v_{\mu\lambda}$ é sobrejetivo, isso significa que $v_\mu$ é
\emph{injetivo}. Sendo também sobrejetivo, prova-se (11.2.9.9).

Após substituir $A_0$ e $M_0$ por $A_\lambda$ e $M_\lambda$ para
$\lambda$ suficientemente grande, podemos, portanto, supor que o
homomorfismo canônico $v_\lambda$ é bijetivo para \emph{todo} $\lambda$.
Ponhamos então
$P_\lambda=N_0\otimes_{C_0}C_\lambda=N_0\otimes_{A_0}A_\lambda$,
de modo que $N=\varinjlim P_\lambda$. Como $C_0$ é uma $A_0$-álgebra
de apresentação finita e $P_0$ é um $C_0$-módulo de apresentação
finita, podemos aplicar a $C$ e $N$ o resultado de (11.2.6.1).
Vemos, assim, que existe um índice $\lambda$ tal que, para
$\mu\geq\lambda$, $P_\mu$ é um $A_\mu$-módulo plano.
Ora, para $\mu\geq\lambda$, temos um diagrama comutativo
\[
\xymatrix{
P_\mu \ar[r] \ar[d]^{w_\mu} & N \ar@{=}[d] \\
N_\mu \ar[r] & N
}
\]
no qual $w_\mu$ é \emph{sobrejetivo}, como decorre das definições.
Como, por III), os homomorfismos $A_\mu\to A$ são injetivos e
$P_\mu$ é um $A_\mu$-módulo plano, o homomorfismo canônico
$P_\mu\to N=P_\mu\otimes_{A_\mu}A$ é \emph{injetivo}.
Concluímos, portanto, do diagrama comutativo precedente que
$w_\mu$ também é injetivo, logo bijetivo. Por conseguinte,
$N_\mu$ é um $A_\mu$-módulo plano para $\mu\geq\lambda$,
o que conclui a demonstração de (11.2.9).
\end{proof}

\begin{remark}[11.2.10]
\label{IV.11.2.10}
Não sabemos se é válida a generalização de (11.2.6, (i))
análoga ao teorema de Raynaud.
\end{remark}

\subsection{Aplicação à eliminação de hipóteses noetherianas}
\label{subsection:IV.11.3}
\label{IV.11.3}

\begin{theorem}[11.3.1]
\label{IV.11.3.1}
Sejam $f:X\to Y$ um morfismo localmente de apresentação finita e
$\sh{F}$ um $\sh{O}_X$-Módulo quase-coerente de apresentação finita.
Então o conjunto $U$ dos pontos $x\in X$ nos quais $\sh{F}$ é
$f$-plano no ponto $x$ é aberto em $X$.
Se, além disso, $\operatorname{Supp}(\sh{F})=X$,
$f|_U$ é um morfismo aberto.
\end{theorem}

\begin{proof}
A segunda afirmação já foi demonstrada em (2.4.6) e é incluída
aqui apenas para lembrá-la.

Como a questão é local em $X$ e em $Y$, podemos supor que ambos
são afins e que $f$ é um morfismo de apresentação finita.
Seja $x\in X$ um ponto no qual $\sh{F}$ é $f$-plano.
Aplicando (11.2.7), podemos supor que
\[
X=X_0\times_{Y_0}Y,\qquad f=f_0\times1,\qquad
\sh{F}=\sh{F}_0\otimes_{\sh{O}_{Y_0}}\sh{O}_Y,
\]
onde $Y_0$ é noetheriano, $f_0:X_0\to Y_0$ é um morfismo
de tipo finito e $\sh{F}_0$ é um $\sh{O}_{X_0}$-Módulo coerente.
Além disso, se $x_0$ é a projeção canônica de $x$ em $X_0$,
podemos supor que $\sh{F}_0$ é $f_0$-plano no ponto $x_0$.
Então, por (11.1.1), o conjunto $U_0$ dos pontos de $X_0$
nos quais $\sh{F}_0$ é $f_0$-plano é uma vizinhança de $x_0$.
Logo, $\sh{F}$ é $f$-plano nos pontos da imagem recíproca de
$U_0$ em $X$ (2.1.4), e isso prova que $U$ é uma vizinhança de $x$.
\end{proof}

\oldpage[IV-3]{133}

\begin{corollary}[11.3.2]
\label{IV.11.3.2}
Sejam $f:X\to Y$ e $g:Y\to Z$ dois morfismos de apresentação finita,
e seja $\sh{F}$ um $\sh{O}_X$-Módulo quase-coerente de apresentação
finita. Então o conjunto $U$ dos $y\in Y$ tais que, para toda
generização $y'$ de $y$, $\sh{F}$ é $(g\circ f)$-plano em todos os
pontos de $f^{-1}(y')$ (isto é, tais que
$\sh{F}\otimes_Y\operatorname{Spec}(\sh{O}_{Y,y})$ é plano relativamente
ao morfismo $X\times_Y\operatorname{Spec}(\sh{O}_{Y,y})\to Z$)
é aberto em $Y$.
\end{corollary}

\begin{proof}
O mesmo raciocínio de (11.1.5) mostra que, se $V$ é o conjunto dos
$x\in X$ nos quais $\sh{F}$ é $(g\circ f)$-plano, então $U$ é o conjunto
dos $y\in Y$ cujas generizações pertencem todas a
$E=Y-f(X-V)$. Como $g\circ f$ é de apresentação finita, $V$ é aberto
em $X$ por (11.3.1), logo ind-construtível (1.9.6); por conseguinte,
$X-V$ é pró-construtível. Como $f$ é quase-compacto, $f(X-V)$ é
pró-construtível em $Y$ (1.9.5, (vii)); portanto, $E$ é
ind-construtível em $Y$. Decorre então de (1.10.1) que $U$ é o
interior de $E$, logo é aberto em $Y$.
\end{proof}

\begin{corollary}[11.3.3]
\label{IV.11.3.3}
Sejam $A$ um anel, $B$ uma $A$-álgebra de apresentação finita,
$C$ uma $B$-álgebra de apresentação finita e $M$ um $C$-módulo
de apresentação finita. Então o conjunto dos
$\mathfrak q\in\operatorname{Spec}(B)$ tais que $M_{\mathfrak q}$
é um $A$-módulo plano é aberto em $\operatorname{Spec}(B)$.
\end{corollary}

\begin{proposition}[11.3.4]
\label{IV.11.3.4}
Sejam $f:X\to S$ um morfismo localmente de apresentação finita,
$\mathcal J$ um Ideal quase-coerente de tipo finito de $\sh{O}_X$,
$Y$ o subpré-esquema fechado de $X$ definido por $\mathcal J$
e $\sh{F}$ um $\sh{O}_X$-Módulo de apresentação finita.
Suponhamos que $\operatorname{gr}_{\mathcal J}^*(\sh{F})$
é $f$-plano. Nessas condições:
\begin{enumerate}[label=\emph{(\roman*)}]
\item $\sh{F}$ é $f$-plano nos pontos de $Y$.
\item Se $\operatorname{gr}_{\mathcal J}^*(\sh{O}_X)$ é $f$-plano,
$\operatorname{gr}_{\mathcal J}^*(\sh{O}_X)$ é uma
$(\sh{O}_X/\mathcal J)$-Álgebra de apresentação finita, e
$\operatorname{gr}_{\mathcal J}^*(\sh{F})$ é um
$(\operatorname{gr}_{\mathcal J}^*(\sh{O}_X))$-Módulo de
apresentação finita.
\item Suponhamos que $\operatorname{gr}_{\mathcal J}^*(\sh{O}_X)$
é $f$-plano (o que implica que $\sh{O}_X/\mathcal J$ é
$f$-plano). Então o conjunto $W$ dos $y\in Y$ tais que
$(\operatorname{gr}_{\mathcal J}^*(\sh{F}))_y$ é um
$(\sh{O}_X/\mathcal J)_y$-módulo plano é aberto em $Y$.
\item Suponhamos que $\operatorname{gr}_{\mathcal J}^*(\sh{O}_X)$
é $f$-plano. Seja $S'\to S$ um morfismo qualquer; sejam
$X'=X\times_SS'$, $p:X'\to X$ a projeção canônica e
$Y'=p^{-1}(Y)=Y\times_SS'$ o subpré-esquema fechado de $X'$
definido por $\mathcal J'=p^*(\mathcal J)\sh{O}_{X'}$.
Ponhamos $\sh{F}'=p^*(\sh{F})
=\sh{F}\otimes_{\sh{O}_S}\sh{O}_{S'}$. Então, se
$W'=p^{-1}(W)$, temos
$\operatorname{gr}_{\mathcal J'}^*(\sh{F}')|W'
\simeq p^*(\operatorname{gr}_{\mathcal J}^*(\sh{F})|W)$,
e $\operatorname{gr}_{\mathcal J'}^*(\sh{F}')$ é um
$(\sh{O}_{X'}/\mathcal J')$-Módulo plano nos pontos de $W'$.
\end{enumerate}
\end{proposition}

\begin{proof}
Como as questões são locais em $X$ e em $S$, podemos supor que
$S=\operatorname{Spec}(A)$ e $X=\operatorname{Spec}(B)$ são afins,
com $B$ uma $A$-álgebra de apresentação finita; que
$\sh{F}=\widetilde M$, com $M$ um $B$-módulo de apresentação
finita; e que $\mathcal J=\widetilde{\mathfrak J}$, com
$\mathfrak J$ um ideal de tipo finito de $B$.
Por (8.9.1) e (8.5.11), existem um subanel noetheriano $A_0$
de $A$ e uma $A_0$-álgebra de tipo finito $B_0$ tais que
$B=B_0\otimes_{A_0}A$; um ideal $\mathfrak J_0$ de $B_0$
tal que $\mathfrak J=\mathfrak J_0B$; e um $B_0$-módulo de
tipo finito $M_0$ tal que
$M=M_0\otimes_{A_0}A=M_0\otimes_{B_0}B$.
Além disso, $A$ é o limite indutivo de suas sub-$A_0$-álgebras
de tipo finito $A_\lambda$. Ponhamos
$B_\lambda=B_0\otimes_{A_0}A_\lambda$,
$M_\lambda=M_0\otimes_{A_0}A_\lambda
=M_0\otimes_{B_0}B_\lambda$ e
$\mathfrak J_\lambda=\mathfrak J_0B_\lambda$,
de modo que $B=\varinjlim B_\lambda$ e
$M=\varinjlim M_\lambda$.
Por hipótese, $\operatorname{gr}_{\mathfrak J}^*(M)$ é um
$A$-módulo plano; logo, por (11.2.9), existe $\lambda$ tal que
$\operatorname{gr}_{\mathfrak J_\lambda}^*(M_\lambda)$ é um
$A_\lambda$-módulo plano. Para provar que $\sh{F}$ é $f$-plano
nos pontos de $Y$, podemos, portanto, restringir-nos ao caso em
que $S$ é noetheriano. Mas então
$(\mathbf 0_{\mathrm I},\,6.1.2)$, aplicado por indução em $n$,
mostra que os $\sh{F}/\mathcal J^{n+1}\sh{F}$ são $f$-planos,
e concluímos por $(\mathbf 0_{\mathrm{III}},\,10.2.6)$.

\oldpage[IV-3]{134}

A demonstração de (ii) reduz-se, do mesmo modo, ao caso em que
$S$ é noetheriano, utilizando (11.2.9). A conclusão é então
evidente, pois $\operatorname{gr}_{\mathfrak J}^*(B)$ é, nesse caso,
uma $(B/\mathfrak J)$-álgebra de tipo finito, e
$\operatorname{gr}_{\mathfrak J}^*(M)$ é um
$\operatorname{gr}_{\mathfrak J}^*(B)$-módulo de tipo finito.

Para provar (iii), reduzimo-nos, como em (i), ao caso em que
$S$ e $X$ são afins. Com as mesmas notações, podemos supor, por
(11.2.9), que $\operatorname{gr}_{\mathfrak J_\lambda}^*(B_\lambda)$
e $\operatorname{gr}_{\mathfrak J_\lambda}^*(M_\lambda)$ são
$A_\lambda$-módulos planos e que
\[
\operatorname{gr}_{\mathfrak J}^*(B)
=\operatorname{gr}_{\mathfrak J_\lambda}^*(B_\lambda)
\otimes_{A_\lambda}A
\quad\text{e}\quad
\operatorname{gr}_{\mathfrak J}^*(M)
=\operatorname{gr}_{\mathfrak J_\lambda}^*(M_\lambda)
\otimes_{A_\lambda}A.
\]
Consequentemente, $\operatorname{gr}_{\mathfrak J}^*(B)$ é uma
$(B/\mathfrak J)$-álgebra de apresentação finita, e
$\operatorname{gr}_{\mathfrak J}^*(M)$ é um
$\operatorname{gr}_{\mathfrak J}^*(B)$-módulo de apresentação
finita. Se $W^m$ é o conjunto dos
$\mathfrak p\in\operatorname{Spec}(B/\mathfrak J)$ tais que
$\operatorname{gr}_{\mathfrak J}^m(M)_{\mathfrak p}$ é um
$(B/\mathfrak J)$-módulo plano, $W^m$ é aberto em
$\operatorname{Spec}(B/\mathfrak J)$ por (11.3.1), e temos
$W=\bigcap_mW^m$; logo, $W$ é estável por generização.
A afirmação (iii) decorre então de (11.3.2), aplicado com
$X=\operatorname{Spec}(\operatorname{gr}_{\mathfrak J}^*(B))$,
$Y=Z=\operatorname{Spec}(B/\mathfrak J)$ e
$\sh{F}=(\operatorname{gr}_{\mathfrak J}^*(M))^{\sim}$.

Por fim, (iv) decorre imediatamente de (11.2.9.2).
\end{proof}

Generalizando a definição de (6.10.1), dizemos que, para um
pré-esquema $X$, um subpré-esquema fechado $Y$ de $X$ definido
por um Ideal quase-coerente $\mathcal J$ de $\sh{O}_X$ e um
$\sh{O}_X$-Módulo quase-coerente $\sh{F}$, o módulo $\sh{F}$
é \emph{normalmente plano ao longo de $Y$ no ponto $y$} se
$(\operatorname{gr}_{\mathcal J}^*(\sh{F}))_y$ é um
$\sh{O}_{Y,y}$-módulo plano. Dizemos que $\sh{F}$ é
\emph{normalmente plano ao longo de $Y$} se o é em todos os
pontos de $Y$.

\begin{corollary}[11.3.5]
\label{IV.11.3.5}
Sob as hipóteses gerais de (11.3.4), suponhamos que
$\operatorname{gr}_{\mathcal J}^*(\sh{O}_X)$ e
$\operatorname{gr}_{\mathcal J}^*(\sh{F})$ sejam $f$-planos.
Então o conjunto $W$ dos $y\in Y$ tais que $\sh{F}$ é
normalmente plano ao longo de $Y$ no ponto $y$ é aberto em
$Y$ e, com as notações de (11.3.4, (iv)), $\sh{F}'$ é
normalmente plano ao longo de $Y'$ em todos os pontos de
$W'=p^{-1}(W)$.
\end{corollary}

\begin{proposition}[11.3.6]
\label{IV.11.3.6}
Com as notações de (8.5.1) e (8.8.1), suponhamos que
$S_\alpha$ é quase-compacto, $X_\alpha$ é de apresentação
finita sobre $S_\alpha$ e $Y_\alpha$ é um subpré-esquema fechado
de $X_\alpha$ definido por um Ideal quase-coerente
$\mathcal J_\alpha$ de tipo finito de $\sh{O}_{X_\alpha}$
tal que $\operatorname{gr}_{\mathcal J_\alpha}^*
(\sh{O}_{X_\alpha})$ é plano sobre $S_\alpha$.
Suponhamos, por fim, que $\sh{F}_\alpha$ é um
$\sh{O}_{X_\alpha}$-Módulo de apresentação finita.
Para que $\sh{F}$ seja normalmente plano ao longo de $Y$,
é necessário e suficiente que exista $\lambda$ tal que
$\sh{F}_\lambda$ seja normalmente plano ao longo de
$Y_\lambda$.
\end{proposition}

\begin{proof}
Observemos que $Y$ (resp.\ $Y_\lambda$) é o subpré-esquema
fechado de $X$ (resp.\ $X_\lambda$) definido por
$\mathcal J=\mathcal J_\alpha\sh{O}_X$
(resp.\ $\mathcal J_\lambda
=\mathcal J_\alpha\sh{O}_{X_\lambda}$).
Pela hipótese e por (11.2.9.2), temos
\[
\operatorname{gr}_{\mathcal J}^*(\sh{O}_X)
=\operatorname{gr}_{\mathcal J_\alpha}^*
(\sh{O}_{X_\alpha})
\otimes_{\sh{O}_{S_\alpha}}\sh{O}_S
\]
e
\[
\operatorname{gr}_{\mathcal J_\lambda}^*(\sh{O}_{X_\lambda})
=\operatorname{gr}_{\mathcal J_\alpha}^*
(\sh{O}_{X_\alpha})
\otimes_{\sh{O}_{S_\alpha}}\sh{O}_{S_\lambda}
\quad\text{para }\lambda\geq\alpha.
\]
Além disso, $\operatorname{gr}_{\mathcal J}^*(\sh{O}_X)$
(resp.\ $\operatorname{gr}_{\mathcal J_\lambda}^*
(\sh{O}_{X_\lambda})$) é plano sobre $S$
(resp.\ $S_\lambda$), o que implica que $Y$ é plano
sobre $S$ (resp.\ $Y_\lambda$ é plano sobre $S_\lambda$).
Se $\sh{F}_\lambda$ é normalmente plano ao longo de
$Y_\lambda$, então
$\operatorname{gr}_{\mathcal J_\lambda}^*
(\sh{F}_\lambda)|Y_\lambda$ é plano sobre $Y_\lambda$,
logo também sobre $S_\lambda$; como
$(\operatorname{gr}_{\mathcal J_\lambda}^*
(\sh{F}_\lambda))_{x_\lambda}=0$ para
$x_\lambda\notin Y_\lambda$,
$\operatorname{gr}_{\mathcal J_\lambda}^*
(\sh{F}_\lambda)$ é plano sobre $S_\lambda$.
Concluímos por (11.2.9.2) que
\[
\operatorname{gr}_{\mathcal J}^*(\sh{F})
=\operatorname{gr}_{\mathcal J_\lambda}^*(\sh{F}_\lambda)
\otimes_{\sh{O}_{S_\lambda}}\sh{O}_S,
\]
e, portanto, $\sh{F}$ é normalmente plano ao longo de $Y$.
Para provar a recíproca, podemos supor que $S_\alpha$ e
$X_\alpha$ são afins e adotar as notações de (11.2.9).
Como $\operatorname{gr}_{\mathfrak J}^*(B)$ é um
$A$-módulo plano, o mesmo vale para
$\operatorname{gr}_{\mathfrak J}^*(M)$; assim, por (11.2.9),
existe $\lambda\geq\alpha$ tal que
$\operatorname{gr}_{\mathfrak J_\lambda}^*(M_\lambda)$
é um $A_\lambda$-módulo plano e
\[
\operatorname{gr}_{\mathfrak J}^*(M)
=\operatorname{gr}_{\mathfrak J_\lambda}^*(M_\lambda)
\otimes_{A_\lambda}A.
\]
Além disso, por (11.3.4, (ii)),
$\operatorname{gr}_{\mathfrak J_\lambda}^*(M_\lambda)$
é um $\operatorname{gr}_{\mathfrak J_\lambda}^*
(B_\lambda)$-módulo de apresentação finita, e
$\operatorname{gr}_{\mathfrak J_\lambda}^*(B_\lambda)$
é uma $(B_\lambda/\mathfrak J_\lambda)$-álgebra
de apresentação finita. A conclusão decorre de (11.2.6).
\end{proof}

\oldpage[IV-3]{135}

\begin{proposition}[11.3.7]
\label{IV.11.3.7}
Sejam $f:X\to Y$ um morfismo localmente de apresentação finita,
$\sh{F}$ e $\sh{F}'$ dois $\sh{O}_X$-Módulos quase-coerentes
de apresentação finita, $u:\sh{F}'\to\sh{F}$ um
$\sh{O}_X$-homomorfismo, $x$ um ponto de $X$ e $y=f(x)$.
Suponhamos que $\sh{F}$ é $f$-plano no ponto $x$.
As condições seguintes são então equivalentes:
\begin{enumerate}[label=\emph{\alph*)}]
\item Temos $(\ker u)_x=0$, e $\operatorname{Coker}u$ é um
$\sh{O}_X$-Módulo $f$-plano no ponto $x$.
\item O homomorfismo
\[
(u\otimes1)_x:
(\sh{F}'\otimes_{\sh{O}_Y}k(y))_x
\longrightarrow
(\sh{F}\otimes_{\sh{O}_Y}k(y))_x
\]
é injetivo.
\end{enumerate}
Se, além disso, $\sh{F}$ é $f$-plano, o conjunto dos pontos
$x$ que satisfazem as condições equivalentes precedentes é
aberto em $X$.
\end{proposition}

\begin{proof}
A condição a) implica b) por
$(\mathbf 0_{\mathrm I},\,6.1.2)$, sem nenhuma hipótese
sobre $\sh{F}$. Para demonstrar a recíproca, podemos
restringir-nos ao caso em que $X$ e $Y$ são afins e,
em seguida, por (8.9.1) e (11.2.7), supor que
\[
X=X_0\times_{Y_0}Y,\qquad f=f_0\times1,\qquad
\sh{F}=\sh{F}_0\otimes_{\sh{O}_{Y_0}}\sh{O}_Y,
\]
\[
\sh{F}'=\sh{F}'_0\otimes_{\sh{O}_{Y_0}}\sh{O}_Y,
\qquad u=u_0\otimes1_Y,
\]
onde $Y_0$ é noetheriano, $f_0:X_0\to Y_0$ é um
morfismo de tipo finito, $\sh{F}_0$ e $\sh{F}'_0$
são dois $\sh{O}_{X_0}$-Módulos coerentes, e
$u_0:\sh{F}'_0\to\sh{F}_0$ é um homomorfismo.
Além disso, se $x_0$ é a projeção de $x$ em $X_0$,
podemos supor que $\sh{F}_0$ é $f_0$-plano no ponto
$x_0$. Ponhamos $y_0=f_0(x_0)$.
Pela transitividade das fibras $(\mathbf I,\,3.6.4)$,
a projeção $f^{-1}(y)\to f_0^{-1}(y_0)$ é
fielmente plana (2.2.13); como
\[
(u\otimes1)_x=((u_0\otimes1)_{x_0})\otimes1_{k(y)},
\]
a hipótese de que $(u\otimes1)_x$ é injetivo implica
que $(u_0\otimes1)_{x_0}$ também o é (2.2.7).
Ora, aplicando $(\mathbf 0_{\mathrm{III}},\,10.2.4)$
aos anéis locais noetherianos $\sh{O}_{Y_0,y_0}$ e
$\sh{O}_{X_0,x_0}$ e aos $\sh{O}_{X_0,x_0}$-módulos
$(\sh{F}'_0)_{x_0}$ e $(\sh{F}_0)_{x_0}$, dos quais
o segundo é plano sobre $\sh{O}_{Y_0,y_0}$,
obtemos $(\ker u_0)_{x_0}=0$ e que
$\operatorname{Coker}u_0$ é $f_0$-plano no ponto
$x_0$. Daí deduzimos, primeiro, que
$\operatorname{Coker}u$ é $f$-plano no ponto $x$
(2.1.4). Além disso, pela exatidão à direita do
produto tensorial, temos
\[
\operatorname{Coker}u
=(\operatorname{Coker}u_0)
\otimes_{\sh{O}_{Y_0}}\sh{O}_Y.
\]
Aplicando então $(\mathbf 0_{\mathrm I},\,6.1.2)$
à sequência, exata por hipótese,
\[
0\longrightarrow(\sh{F}'_0)_{x_0}
\longrightarrow(\sh{F}_0)_{x_0}
\longrightarrow(\operatorname{Coker}u_0)_{x_0}
\longrightarrow0,
\]
concluímos que $u_x$ é um homomorfismo injetivo.

Por fim, decorre de (11.1.2) que o conjunto $U_0$
dos pontos $z_0\in X_0$ nos quais
$(u_0\otimes1)_{z_0}$ é injetivo é aberto.
Por planicidade, para todo $z\in X$ situado
sobre $z_0\in U_0$, o homomorfismo
$(u\otimes1)_z$ é injetivo. Portanto, o conjunto
desses pontos contém a imagem recíproca $U$ de
$U_0$ em $X$, que é uma vizinhança de $x$.
Isso conclui a demonstração.
\end{proof}

\begin{theorem}[11.3.8]
\label{IV.11.3.8}
Sejam $f:X\to Y$ um morfismo localmente de apresentação
finita, $\sh{F}$ um $\sh{O}_X$-Módulo quase-coerente de
apresentação finita e $(g_i)_{1\leq i\leq n}$ uma sequência
de seções de $\sh{O}_X$ sobre $X$. Ponhamos
\[
\sh{F}_i=\sh{F}\big/\sum_{j=1}^{i}g_j\sh{F}
\quad\text{para }0\leq i\leq n
\quad(\text{com }\sh{F}_0=\sh{F}).
\]
Sejam $x$ um ponto de $X$, $y=f(x)$, e ponhamos
$X_y=f^{-1}(y)=X\times_Y\operatorname{Spec}(k(y))$ e
$\sh{F}_y=\sh{F}\otimes_{\sh{O}_Y}k(y)$,
que é um $\sh{O}_{X_y}$-Módulo.
Suponhamos que os $(g_i)_x$ pertencem ao ideal maximal
$\mathfrak m_x$. As condições seguintes são equivalentes:
\begin{enumerate}[label=\emph{\alph*)}]
\item A sequência $((g_i)_x)$ é $\sh{F}_x$-regular,
e os $\sh{F}_i$ ($0\leq i\leq n$) são $f$-planos
no ponto $x$.
\item A sequência $((g_i)_x)$ é $\sh{F}_x$-regular,
e $\sh{F}_n$ é $f$-plano no ponto $x$.
\item[\emph{b$'$)}] Existe uma vizinhança $U$ de $x$
tal que a sequência $(g_i|U)$ é
$(\sh{F}|U)$-quase-regular, e $\sh{F}_n$ é
$f$-plano no ponto $x$.

\oldpage[IV-3]{136}

\item $\sh{F}$ é $f$-plano no ponto $x$, e a
sequência $((g_i\otimes1)_x)$ de elementos de
$\sh{O}_{X_y,x}$, imagens dos $(g_i)_x$, é
$(\sh{F}_y)_x$-regular.
\item $\sh{F}$ é $f$-plano no ponto $x$ e, para todo
morfismo $Y'\to Y$ e todo ponto
$x'\in X'=X\times_YY'$ situado sobre $x$,
se pomos $\sh{F}'=\sh{F}\otimes_YY'$,
a sequência $(g_i\otimes1)_{x'}$ de elementos de
$\sh{O}_{X',x'}$ é $\sh{F}'_{x'}$-regular.
\end{enumerate}
Além disso, o conjunto dos $x\in X$ que satisfazem
essas condições é aberto no conjunto $Z$ dos $x$
tais que $g_i(x)=0$ para todo $i$.
\end{theorem}

\begin{proof}
O fato de que a) implica d) decorre imediatamente de
$(\mathbf 0,\,15.1.15)$, e c) é um caso particular de d);
além disso, a) implica b) trivialmente. Mostremos que
b) ou c) implica a). Os $\sh{F}_i$ são de apresentação
finita; o fato de que c) implica a) decorre então
imediatamente de (11.3.7) por indução em $n$, e isso
mostra também que o conjunto dos $x\in X$ que
satisfazem c) é aberto em $Z$.

Para mostrar que b) implica a), a indução em $n$
reduz-nos imediatamente ao caso $n=1$; escreveremos
$g$ em lugar de $g_1$. Como a questão é local em $X$
e em $Y$, podemos supor que
$Y=\operatorname{Spec}(A)$ e $X=\operatorname{Spec}(B)$
são afins, com $B$ uma $A$-álgebra de apresentação
finita e $\sh{F}=\widetilde M$, onde $M$ é um
$B$-módulo de apresentação finita.
Podemos então, por (8.9.1), escrever
$B=B_0\otimes_{A_0}A$ e $M=M_0\otimes_{A_0}A$,
onde $A_0$ é um subanel noetheriano de $A$,
$B_0$ é uma $A_0$-álgebra de tipo finito e
$M_0$ é um $B_0$-módulo de tipo finito.
Além disso, $A$ é o limite indutivo filtrante de
seus subanéis $A_\lambda$ que são $A_0$-álgebras
de tipo finito e, portanto, noetherianos.
Se pusermos $B_\lambda=B_0\otimes_{A_0}A_\lambda$
e $M_\lambda=M_0\otimes_{A_0}A_\lambda$,
então $B_\lambda$ é uma $A_\lambda$-álgebra de
tipo finito, $M_\lambda$ é um $B_\lambda$-módulo
de tipo finito, e temos $B=\varinjlim B_\lambda$,
$M=\varinjlim M_\lambda$. Existe, portanto, um
índice $\lambda$ e um elemento $g_\lambda\in B_\lambda$
tais que $g=g_\lambda\otimes1$.
Voltando às notações geométricas e pondo
$Y_\lambda=\operatorname{Spec}(A_\lambda)$,
$X_\lambda=\operatorname{Spec}(B_\lambda)$ e
$\sh{F}_\lambda=\widetilde{M_\lambda}$,
basta provar que existe $\mu\geq\lambda$ tal que,
no ponto $x_\mu\in X_\mu$ projeção de $x$,
$(g_\mu)_{x_\mu}$ é
$(\sh{F}_\mu)_{x_\mu}$-regular e
$\sh{F}_\mu/g_\mu\sh{F}_\mu$ é plano sobre
$Y_\mu$ no ponto $x_\mu$. Com efeito, deduziremos
então de $(\mathbf 0,\,15.1.16)$ que
$\sh{F}_\mu$ é plano sobre $Y_\mu$ no ponto
$x_\mu$ e, portanto, $\sh{F}$ é plano sobre
$Y$ no ponto $x$.

O fato de que b) implica a) decorrerá, pois, da
proposição seguinte:

\begin{proposition}[11.3.9]
\label{IV.11.3.9}
Com as notações e hipóteses gerais de (8.5.1) e
(8.8.1), suponhamos que
$f_\alpha:X_\alpha\to Y_\alpha$ seja um morfismo
localmente de apresentação finita. Sejam
$\sh{F}_\alpha$ e $\sh{G}_\alpha$ dois
$\sh{O}_{X_\alpha}$-módulos quase-coerentes de
apresentação finita,
$h_\alpha:\sh{F}_\alpha\to\sh{G}_\alpha$ um
$\sh{O}_{X_\alpha}$-homomorfismo e
$\sh{H}_\alpha$ seu conúcleo. Sejam, por fim, $x$
um ponto de $X$ e $x_\lambda$ sua projeção em
$X_\lambda$ para $\lambda\geq\alpha$.
Para que $h_x$ seja injetivo e
$\sh{H}=\operatorname{Coker}h$ seja $f$-plano no
ponto $x$, é necessário e suficiente que exista
$\lambda\geq\alpha$ tal que
$(h_\lambda)_{x_\lambda}$ seja injetivo e
$\sh{H}_\lambda=\operatorname{Coker}h_\lambda$
seja $f_\lambda$-plano no ponto $x_\lambda$.
Além disso, o conjunto dos $x\in X$ que têm
as propriedades precedentes é aberto em $X$.
\end{proposition}

\begin{proof}
Recordemos que, em virtude da exatidão à direita
do produto tensorial, temos
\[
\sh{H}_\mu=v_{\lambda\mu}^*(\sh{H}_\lambda)
\quad\text{para }\lambda\leq\mu,
\qquad
\sh{H}=v_\lambda^*(\sh{H}_\lambda),
\]
o que justifica as notações.

A suficiência da condição decorre do seguinte:
se a sequência
\[
0\longrightarrow(\sh{F}_\lambda)_{x_\lambda}
\longrightarrow(\sh{G}_\lambda)_{x_\lambda}
\longrightarrow(\sh{H}_\lambda)_{x_\lambda}
\longrightarrow0
\]
é exata e $(\sh{H}_\lambda)_{x_\lambda}$
é um $\sh{O}_{f_\lambda(x_\lambda)}$-módulo
plano, então $\sh{H}_x$ é um
$\sh{O}_{f(x)}$-módulo plano por mudança de
base (2.1.4), e a sequência
$0\to\sh{F}_x\to\sh{G}_x\to\sh{H}_x\to0$
é exata (2.1.8). Para provar que

\oldpage[IV-3]{137}

a condição é necessária, observemos que
$\sh{H}$ é de apresentação finita. Como a questão
é local em $X$ e em $Y$, podemos supor que $X$
e $Y$ sejam afins e, em virtude de
\sref{IV.11.3.1}, que $\sh{H}$ seja $f$-plano.
Consideremos agora o lema seguinte.

\begin{lemma}[11.3.9.1]
\label{IV.11.3.9.1}
Sejam $f:X\to Y$ um morfismo localmente de
apresentação finita, $\sh{G}$ e $\sh{H}$
dois $\sh{O}_X$-módulos quase-coerentes de
apresentação finita tais que $\sh{H}$ é
$f$-plano, e $p:\sh{G}\to\sh{H}$ um
homomorfismo sobrejetivo. Então $\ker(p)$
é um $\sh{O}_X$-módulo de apresentação finita.
\end{lemma}

\begin{proof}
Podemos, com efeito, supor que $X$ e $Y$ sejam
afins e que existam um morfismo $Y\to Y_0$,
com $Y_0$ noetheriano, um morfismo
$f_0:X_0\to Y_0$ de tipo finito tais que
$X=X_0\times_{Y_0}Y$ e $f=(f_0)_{(Y)}$,
dois $\sh{O}_{X_0}$-módulos coerentes
$\sh{G}_0$ e $\sh{H}_0$ e um homomorfismo
$p_0:\sh{G}_0\to\sh{H}_0$ tais que
$\sh{G}$, $\sh{H}$ e $p$ se obtêm de
$\sh{G}_0$, $\sh{H}_0$ e $p_0$ por mudança
de base (8.9.1 e 8.5.2). Além disso, podemos
supor que $p_0$ seja sobrejetivo (8.5.7)
e $\sh{H}_0$ seja $f_0$-plano
\sref{IV.11.2.7}. Então, se
$\sh{K}_0=\ker(p_0)$, $\sh{K}_0$ é um
$\sh{O}_{X_0}$-módulo coerente
$(\mathbf 0_{\mathrm I},\,5.3.4)$ e, em
virtude de $(\mathbf 0_{\mathrm I},\,6.1.2)$,
$\sh{K}=\ker(p)$ se obtém de $\sh{K}_0$
por mudança de base; portanto, é de
apresentação finita.
\end{proof}

Demonstrado esse lema, ponhamos
$\sh{K}=\ker(\sh{G}\to\sh{H})$, que é,
portanto, de apresentação finita. Temos um
homomorfismo canônico $q:\sh{F}\to\sh{K}$
e, por hipótese,
$q_x:\sh{F}_x\to\sh{K}_x$ é um isomorfismo.
Por conseguinte, por
$(\mathbf 0_{\mathrm I},\,5.2.7)$, existe
uma vizinhança $U$ de $x$ tal que $q|U$
é um isomorfismo; restringindo $X$, podemos
supor que a sequência
\[
0\longrightarrow\sh{F}\xrightarrow{h}\sh{G}
\longrightarrow\sh{H}\longrightarrow0
\]
é \emph{exata}. Nessas condições, decorre de
\sref{IV.11.2.6} que, para $\lambda$
suficientemente grande, $\sh{H}_\lambda$
é $f_\lambda$-plano. Ponhamos
$\sh{K}_\lambda=\ker(\sh{G}_\lambda
\to\sh{H}_\lambda)$ e
$\sh{K}_\mu=v_{\lambda\mu}^*
(\sh{K}_\lambda)$ para $\mu\geq\lambda$;
decorre de $(\mathbf 0_{\mathrm I},\,6.1.2)$
que
\[
\sh{K}_\mu=\ker(\sh{G}_\mu\to\sh{H}_\mu)
\quad\text{e}\quad
\sh{K}=v_\lambda^*(\sh{K}_\lambda)
=\ker(\sh{G}\to\sh{H})=\sh{F}.
\]
Por outro lado, para todo $\mu\geq\lambda$,
temos um homomorfismo canônico
$q_\mu:\sh{F}_\mu\to\sh{K}_\mu$ com
$q_\mu=v_{\lambda\mu}^*(q_\lambda)$ e
$q=v_\mu^*(q_\mu)$. Como $q$ é um
isomorfismo, decorre de (8.5.2.4) e
\sref{IV.11.3.9.1} que existe
$\mu\geq\lambda$ tal que $q_\mu$ é um
isomorfismo e, por conseguinte,
$h_\mu:\sh{F}_\mu\to\sh{G}_\mu$
é injetivo.\footnote{Nota de fonte (p.~137): a edição impressa afirma que $\sh{H}_\lambda$ é um $\sh{O}_{X_\lambda}$-módulo plano e denomina o homomorfismo final $u_\mu$. Mantemos no texto a planicidade relativa $f_\lambda$-plana, exigida pela proposição e obtida de (11.2.6), e $h_\mu$, da família definida por $h_\alpha$. As duas leituras impressas ficam conservadas nesta nota.}
\end{proof}

Voltemos à demonstração de \sref{IV.11.3.8}.

Como o conjunto dos $x\in X$ que satisfazem b)
é aberto em $X$, é claro que b) implica b$'$.
Mostremos, por fim, que b$'$ implica c).
Em primeiro lugar, $\sh{F}_n$ é $f$-plano
em uma vizinhança de $x$ \sref{IV.11.3.1};
podemos, portanto, reduzir-nos ao caso em que
$U=X$ e $\sh{F}_n$ é $f$-plano. Como, por
definição, $\operatorname{gr}_{\mathcal J}^*
(\sh{F})$ é isomorfo a
$\sh{F}_n\otimes_{\sh{O}_Y}
\sh{O}_Y[T_1,\ldots,T_n]$
$(\mathbf 0,\,15.2.2)$, ele é $f$-plano;
deduzimos então de \sref{IV.11.3.4}, (i),
que a própria $\sh{F}$ é $f$-plana em uma
vizinhança de $x$. Por outro lado, se
$(\mathcal J_y)_x$ é o ideal de
$\sh{O}_{X_y,x}$ gerado pelos
$(g_i\otimes1)_x$, decorre de
\sref{IV.11.2.9.2} que, no diagrama
\[
\xymatrix{
(\operatorname{gr}_{\mathcal J_x}^0(\sh{F}_x)
 [T_1,\ldots,T_n])\otimes_{\sh{O}_y}k(y)
 \ar[r]\ar[d] &
\operatorname{gr}_{\mathcal J_x}^*(\sh{F}_x)
 \otimes_{\sh{O}_y}k(y)\ar[d]\\
\operatorname{gr}_{(\mathcal J_y)_x}^0
 ((\sh{F}_y)_x)[T_1,\ldots,T_n]\ar[r] &
\operatorname{gr}_{(\mathcal J_y)_x}^*
 ((\sh{F}_y)_x)
}
\]

\oldpage[IV-3]{138}

as setas verticais são isomorfismos.
Como a primeira seta horizontal é um
isomorfismo, a segunda também o é;
portanto, a sequência $((g_i\otimes1)_x)$
é $(\sh{F}_y)_x$-quase-regular e,
por conseguinte, também
$(\sh{F}_y)_x$-regular, pois $X_y$ é
localmente de tipo finito sobre $k(y)$.
\end{proof}

\begin{theorem}[11.3.10]
\label{IV.11.3.10}
(Critério de planicidade pelas fibras). ---
Sejam $S$ um pré-esquema,
$g:X\to S$ e $h:Y\to S$ dois morfismos,
$f:X\to Y$ um $S$-morfismo,
$\sh{F}$ um $\sh{O}_X$-módulo quase-coerente,
$x$ um ponto de $X$, $y=f(x)$ e
$s=h(y)=g(x)$. Suponhamos que se verifique
uma das duas hipóteses seguintes:
\begin{enumerate}[label=\arabic{enumi}\textsuperscript{o}]
\item $S$, $X$ e $Y$ são localmente noetherianos
e $\sh{F}$ é coerente.
\item $g$ e $h$ são localmente de apresentação
finita e $\sh{F}$ é de apresentação finita.
\end{enumerate}
Então, com as notações de (9.4.1), se
$\sh{F}_x\neq0$, as condições seguintes
são equivalentes:
\begin{enumerate}[label=\emph{\alph*)}]
\item $\sh{F}$ é $g$-plano no ponto $x$ e
$\sh{F}_s$ é $f_s$-plano no ponto $x$.
\item $h$ é plano no ponto $y$ e $\sh{F}$
é $f$-plano no ponto $x$.
\end{enumerate}
Além disso, sob a hipótese
2\textsuperscript{o}, o conjunto dos
$x\in X$ que satisfazem a condição b)
é aberto em $X$.
\end{theorem}

\begin{proof}
A última afirmação decorre de
\sref{IV.11.3.1} aplicado, por um lado,
a $\sh{O}_Y$ e ao morfismo $h$ e,
por outro, a $\sh{F}$ e ao morfismo $f$
(que é localmente de apresentação finita).

Como $g=h\circ f$, é claro que b) implica
a) sem supor 1\textsuperscript{o} nem
2\textsuperscript{o} (2.1.6 e 2.1.4).
Para provar que a) implica b), podemos
reduzir-nos ao caso em que $S$, $X$ e $Y$
são afins; sob a hipótese
2\textsuperscript{o}, aplicando
\sref{IV.11.2.7}, reduzimo-nos também
ao caso em que $S$ é noetheriano,
isto é, ao caso em que se verifica a
hipótese 1\textsuperscript{o}. A afirmação
equivale então ao lema seguinte, que melhora
$(\mathbf 0_{\mathrm{III}},\,10.2.5)$:

\begin{lemma}[11.3.10.1]
\label{IV.11.3.10.1}
Sejam $\rho:A\to B$ e $\sigma:B\to C$
homomorfismos locais de anéis locais
noetherianos, $k$ o corpo residual de $A$
e $M$ um $C$-módulo não nulo de tipo
finito. Então as condições seguintes
são equivalentes:
\begin{enumerate}[label=\emph{\alph*)}]
\item $M$ é um $A$-módulo plano e
$M\otimes_Ak$ é um
$(B\otimes_Ak)$-módulo plano.
\item $B$ é um $A$-módulo plano e
$M$ é um $B$-módulo plano.
\end{enumerate}
\end{lemma}

Estabeleceremos primeiro o lema mais geral
seguinte:

\begin{lemma}[11.3.10.2]
\label{IV.11.3.10.2}
Sejam $A$ um anel comutativo, $B$ uma
$A$-álgebra comutativa, $\mathfrak J$
um ideal de $A$ e $M$ um $B$-módulo.
Consideremos, por um lado, as condições
seguintes:
\begin{enumerate}[label=\emph{(\roman*)}]
\item $\mathfrak J$ é nilpotente.
\item $B$ é noetheriano e $M$ é
idealmente separado para a topologia
$\mathfrak JB$-pré-ádica
$(\mathbf 0_{\mathrm{III}},\,10.2.1)$.
\item $\mathfrak JB$ está contido no
radical de Jacobson de $B$.
\end{enumerate}
Consideremos, por outro lado, as quatro
propriedades seguintes:
\begin{enumerate}[label=\emph{\alph*)}]
\item $M$ é um $B$-módulo plano.
\item $B$ é um $A$-módulo plano.
\item $M$ é um $A$-módulo plano e
$M/\mathfrak JM$ é um
$(B/\mathfrak JB)$-módulo plano.
\item $B/\mathfrak JB$ é um
$(A/\mathfrak J)$-módulo plano e,
para todo ideal maximal
$\mathfrak m\supset\mathfrak JB$ de $B$,
temos $\mathfrak mM\neq M$.
\end{enumerate}
Então:
\begin{enumerate}[label=\arabic{enumi}\textsuperscript{o}]
\item Se uma das condições (i) ou (ii)
se verifica, a conjunção de a) e b)
implica c), e c) implica a).

\oldpage[IV-3]{139}

\item Se a condição (i) ou a conjunção
de (ii) e (iii) se verifica, a
conjunção de c) e d) implica a
conjunção de a) e b).
\end{enumerate}
\end{lemma}

\begin{proof}
\noindent 1\textsuperscript{o} A primeira
afirmação é imediata
$(\mathbf 0_{\mathrm I},\,6.2.1)$.
Suponhamos, pois, que c) se verifica
e provemos a). Consideremos os anéis
graduados $\operatorname{gr}_{\mathfrak J}^*(A)$
e $\operatorname{gr}_{\mathfrak J}^*(B)$,
e o módulo graduado
$\operatorname{gr}_{\mathfrak J}^*(M)$
(tanto sobre $\operatorname{gr}_{\mathfrak J}^*(A)$
como sobre $\operatorname{gr}_{\mathfrak J}^*(B)$),
relativos às filtrações
$\mathfrak J$-pré-ádicas, bem como as
aplicações canônicas sobrejetivas
$(\mathbf 0_{\mathrm{III}},\,10.1.1.2)$
\begin{align*}
u&:\operatorname{gr}^0(B)
\otimes_{\operatorname{gr}^0(A)}
\operatorname{gr}^*(A)
\longrightarrow\operatorname{gr}^*(B),\\
v&:\operatorname{gr}^0(M)
\otimes_{\operatorname{gr}^0(A)}
\operatorname{gr}^*(A)
\longrightarrow\operatorname{gr}^*(M),\\
w&:\operatorname{gr}^0(M)
\otimes_{\operatorname{gr}^0(B)}
\operatorname{gr}^*(B)
\longrightarrow\operatorname{gr}^*(M).
\end{align*}
É claro que temos um diagrama comutativo
\begin{equation}
\xymatrix{
\operatorname{gr}^0(M)
\otimes_{\operatorname{gr}^0(A)}
\operatorname{gr}^*(A)
\ar[rr]^v\ar[dr]_{1\otimes u}
&&\operatorname{gr}^*(M)\\
&\operatorname{gr}^0(M)
\otimes_{\operatorname{gr}^0(B)}
\operatorname{gr}^*(B)\ar[ur]_w&
}
\label{IV.11.3.10.3}
\tag{11.3.10.3}
\end{equation}
A hipótese c) implica que $v$ é bijetiva
$(\mathbf 0_{\mathrm{III}},\,10.2.1)$;
como as outras duas aplicações do diagrama
são sobrejetivas, também elas são
bijetivas. Mas, como, em virtude da
hipótese c), $M/\mathfrak JM$ é um
$(B/\mathfrak JB)$-módulo plano,
decorre de $(\mathbf 0_{\mathrm{III}},\,10.2.1)$
que $M$ é um $B$-módulo plano.

\noindent 2\textsuperscript{o} Cada uma das
condições (i) e (iii) implica que
todo ideal maximal de $B$ contém
$\mathfrak JB$. Decorre, portanto, de
1\textsuperscript{o} e da conjunção de
c) e d) que $M$ é um $B$-módulo
\emph{fielmente plano} e, por conseguinte,
$\operatorname{gr}^0(M)$ é um
$\operatorname{gr}^0(B)$-módulo
\emph{fielmente plano}
$(\mathbf 0_{\mathrm I},\,6.2.1)$.
Vimos em 1\textsuperscript{o} que a
hipótese c) implica que as três
aplicações do diagrama (11.3.10.3)
são bijetivas; o fato de que
$\operatorname{gr}^0(M)$ é um
$\operatorname{gr}^0(B)$-módulo fielmente
plano implica, portanto, que $u$
também é bijetiva
$(\mathbf 0_{\mathrm I},\,6.4.1)$.
Por outro lado, as condições (ii) e
(iii) implicam que $B$ é um
$A$-módulo idealmente separado para
a filtração $\mathfrak J$-pré-ádica
(Bourbaki, \emph{Alg. comm.}, cap.~III,
\S~5, n\textsuperscript{o}~4, prop.~2);
deduzimos então novamente de
$(\mathbf 0_{\mathrm{III}},\,10.2.1)$
que, se a condição (i) ou a conjunção
de (ii) e (iii) se verifica, $B$ é
um $A$-módulo plano.
\end{proof}

\phantomsection\label{IV.11.3.10.4}
\noindent\textbf{(11.3.10.4)} Estabelecido
o Lema \sref{IV.11.3.10.2}, dele
deduzimos o Lema \sref{IV.11.3.10.1}
tomando para $\mathfrak J$ o ideal
maximal de $A$ e observando que,
então, as condições (ii) e (iii) de
\sref{IV.11.3.10.2} se verificam
(Bourbaki, \emph{Alg. comm.}, cap.~III,
\S~5, n\textsuperscript{o}~4, prop.~2).
Isso encerra também a demonstração
de \sref{IV.11.3.10}.\footnote{Nota de fonte (p.~139): o impresso remete aqui a (11.3.7). A conclusão encerra o Teorema (11.3.10), por meio dos Lemas (11.3.10.1) e (11.3.10.2); conservamos a remissão impressa nesta nota.}
\end{proof}

\begin{corollary}[11.3.11]
\label{IV.11.3.11}
Sejam $g:X\to S$ e $h:Y\to S$ dois
morfismos localmente de apresentação
finita, e $f:X\to Y$ um $S$-morfismo.
As condições seguintes são equivalentes:
\begin{enumerate}[label=\emph{\alph*)}]
\item $g$ é plano e
$f_s:X_s\to Y_s$ é plano para todo
$s\in S$.
\item $h$ é plano em todos os pontos
de $f(X)$ e $f$ é plano.
\end{enumerate}
\end{corollary}

\begin{proof}
Basta aplicar \sref{IV.11.3.10} com
$\sh{F}=\sh{O}_X$.
\end{proof}

\begin{remark}[11.3.12]
\label{IV.11.3.12}
Seria interessante poder dar demonstrações
de \sref{IV.11.3.2} e \sref{IV.11.3.10}
que não utilizassem a passagem ao limite
indutivo; para isso bastaria demonstrar o
critério seguinte:

Sejam $A$ um anel, $B$ uma $A$-álgebra
de apresentação finita, $M$ um
$B$-módulo de apresentação finita,
$\mathfrak J$ um

\oldpage[IV-3]{140}

ideal de $A$ e $\mathfrak p$ um ideal
primo de $B$ que contém $\mathfrak JB$.
Então as duas condições seguintes são
equivalentes:
\begin{enumerate}[label=\emph{\alph*)}]
\item $M_{\mathfrak p}$ é um
$A$-módulo plano.
\item $M_{\mathfrak p}/\mathfrak J
M_{\mathfrak p}$ é um
$(A/\mathfrak J)$-módulo plano e
$\operatorname{Tor}_1^A
(A/\mathfrak J,M_{\mathfrak p})=0$.
\end{enumerate}
Quando $A$ é \emph{noetheriano},
isso decorre de
$(\mathbf 0_{\mathrm{III}},\,10.2.2)$
aplicado à $A$-álgebra noetheriana
$B_{\mathfrak p}$; pode-se deduzir daí
o enunciado geral mediante uma passagem
ao limite indutivo?

Convém observar, a esse respeito, que
tal generalização certamente não é
possível quando se substitui a condição
b) acima por uma das condições c), d)
de $(\mathbf 0_{\mathrm{III}},\,10.2.1)$,
com $M$ substituído por $M_{\mathfrak p}$.
Tomemos, por exemplo, para $A$ um anel
local cujo ideal maximal $\mathfrak J$
é principal e tal que a interseção
\[
\mathfrak R=\bigcap_{n\geq0}
\mathfrak J^{n+1}
\]
não se reduza a $0$ (por exemplo, o
anel dos germes de funções numéricas
infinitamente diferenciáveis em uma
vizinhança de $0$ em $\mathbf R$).
Tomemos $B=A$, $\mathfrak p=\mathfrak J$
e $M=A/\mathfrak N$, onde
$\mathfrak N$ é um submódulo não nulo
de $\mathfrak R$ gerado por um único
elemento; $M$ é, portanto, de
apresentação finita. É claro que
$M$ não é um $A$-módulo plano,
pois, sendo de apresentação finita,
seria livre (Bourbaki,
\emph{Alg. comm.}, cap.~II, \S~3,
n\textsuperscript{o}~2, cor.~2 da
prop.~5), o que é absurdo, já que
$\mathfrak N\neq0$. Entretanto,
\[
\mathfrak J^kM/\mathfrak J^{k+n}M
\simeq A/\mathfrak J^n
\]
quaisquer que sejam os inteiros
positivos $k$ e $n$; portanto,
$M$ satisfaz as condições c) e d)
de $(\mathbf 0_{\mathrm{III}},\,10.2.1)$,
pois $A/\mathfrak J$ é um corpo.
\end{remark}

\begin{proposition}[11.3.13]
\label{IV.11.3.13}
Sejam $f:X\to Y$ um morfismo de
pré-esquemas e $x$ um ponto de $X$
tal que $f$ é plano no ponto $x$;
ponhamos $y=f(x)$.
\begin{enumerate}[label=\emph{(\roman*)}]
\item Se $X$ é reduzido (resp.\ normal)
no ponto $x$, então $Y$ é reduzido
(resp.\ normal) no ponto $y$.
\item Suponhamos, além disso, que $f$
é de apresentação finita no ponto $x$.
Então, se $Y$ é reduzido (resp.\ normal)
no ponto $y$ e o morfismo $f$ é
reduzido (resp.\ normal) no ponto $x$
(6.8.1), $X$ é reduzido (resp.\ normal)
no ponto $x$.
\end{enumerate}
\end{proposition}

\begin{proof}
A primeira afirmação é recordada apenas
por completude, pois foi demonstrada em
(2.1.13). Para provar (ii), podemos
reduzir-nos ao caso em que
$X=\operatorname{Spec}(B)$ e
$Y=\operatorname{Spec}(A)$, sendo $A$
um anel local reduzido (resp.\ um domínio de
integridade local integralmente fechado) e $B$ uma
$A$-álgebra de apresentação finita.
Podemos então escrever, por (8.9.1),
\[
B=B_0\otimes_{A_0}A,
\]
onde $A_0$ é uma sub-$\mathbf Z$-álgebra
de tipo finito de $A$ e $B_0$ uma
$A_0$-álgebra de tipo finito. Seja
$(C_\alpha)$ a família filtrante
crescente das sub-$A_0$-álgebras de
tipo finito de $A$, que são, portanto,
$\mathbf Z$-álgebras de tipo finito;
temos $A=\varinjlim C_\alpha$.
Distinguamos agora os dois casos:

\noindent 1\textsuperscript{o} Suponhamos
que $A$ é reduzido e $f$ é reduzido
no ponto $x$. Se $\mathfrak m$ é o
ideal maximal de $A$, seja
$\mathfrak p_\alpha=\mathfrak m\cap
C_\alpha$ o ideal primo correspondente
e ponhamos
$A'_\alpha=(C_\alpha)_{\mathfrak p_\alpha}$;
assim, temos também
$A=\varinjlim A'_\alpha$ (5.13.3).
Ponhamos
$Y_\alpha=\operatorname{Spec}(A'_\alpha)$
e
$X_\alpha=\operatorname{Spec}
(B_0\otimes_{A_0}A'_\alpha)$;
sejam $f_\alpha:X_\alpha\to Y_\alpha$,
$x_\alpha$ (resp.\ $y_\alpha$) a
projeção de $x$ (resp.\ $y$) em
$X_\alpha$ (resp.\ $Y_\alpha$).
Como $f$ é reduzido no ponto $x$,
existe $\alpha_0$ tal que $f_\alpha$
é reduzido no ponto $x_\alpha$ para
$\alpha\geq\alpha_0$
((6.7.8) e \sref{IV.11.2.6});
como $Y_\alpha$ e $X_\alpha$ são
noetherianos e $A'_\alpha$ é
reduzido (pois $C_\alpha$ o é),
deduzimos de (3.3.5) que $X_\alpha$
é reduzido no ponto $x_\alpha$.
Mas, como
$B=\varinjlim
(B_0\otimes_{A_0}A'_\alpha)$,
temos também
$\sh{O}_{X,x}=\varinjlim
\sh{O}_{X_\alpha,x_\alpha}$ (5.13.3)
e, por conseguinte,
$\sh{O}_{X,x}$ é reduzido (5.13.2).

\noindent 2\textsuperscript{o} Suponhamos
que $A$ seja integralmente fechado
e $f$ seja normal no ponto $x$.
Como $C_\alpha$ é universalmente
japonês (7.7.4), seu fecho integral
$C'_\alpha$ é uma $C_\alpha$-álgebra
finita, evidentemente contida em $A$.
Seja $\mathfrak p'_\alpha$ o ideal primo
$\mathfrak m\cap C'_\alpha$, e ponhamos
$A''_\alpha=(C'_\alpha)_{\mathfrak p'_\alpha}$;
assim, $A''_\alpha$ é um domínio de integridade
local noetheriano integralmente
fechado, e $A=\varinjlim A''_\alpha$
(5.13.3). Ponhamos
$Y'_\alpha=\operatorname{Spec}(A''_\alpha)$,
$X'_\alpha=\operatorname{Spec}
(B_0\otimes_{A_0}A''_\alpha)$, e sejam
$f'_\alpha:X'_\alpha\to Y'_\alpha$,
$x'_\alpha$ (resp.\ $y'_\alpha$) a
projeção de $x$ (resp.\ $y$) em
$X'_\alpha$ (resp.\ $Y'_\alpha$).
Como $f$ é normal no ponto $x$,
existe $\alpha_0$ tal que, para
$\alpha\geq\alpha_0$, $f'_\alpha$
é normal no ponto $x'_\alpha$
((6.7.8) e \sref{IV.11.2.6});

\oldpage[IV-3]{141}

como $X'_\alpha$ e $Y'_\alpha$ são
noetherianos e $A''_\alpha$ é
integralmente fechado, deduzimos
de (6.5.4) que $X'_\alpha$ é normal
no ponto $x'_\alpha$. Mas os morfismos
$\operatorname{Spec}(A''_\beta)\to
\operatorname{Spec}(A''_\alpha)$, para
$\alpha\leq\beta$, são dominantes;
como, em virtude de \sref{IV.11.3.1},
podemos supor que os $f'_\alpha$
são planos para $\alpha\geq\alpha_0$,
resulta de (2.3.7) que toda componente
irredutível de $X'_\beta$ domina uma
componente irredutível de $X'_\alpha$.
Concluímos então, por (5.13.4), aplicado
aos pré-esquemas
\[
\operatorname{Spec}
(\sh{O}_{X_0,x_0}\otimes_{A_0}A''_\alpha),
\]
que $X$ é normal no ponto $x$.
\end{proof}

\begin{corollary}[11.3.14]
\label{IV.11.3.14}
Sejam $f:X\to Y$ um morfismo
localmente de apresentação finita,
$x$ um ponto de $X$ e $y=f(x)$.
Se $Y$ é geometricamente unirramificado
no ponto $y$ (6.15.1) e o morfismo
$f$ é normal no ponto $x$ (6.8.1),
então $X$ é geometricamente
unirramificado no ponto $x$.
\end{corollary}

\begin{proof}
Podemos evidentemente nos limitar ao
caso em que $Y=\operatorname{Spec}(A)$,
sendo $A$ um domínio de integridade local
geometricamente unirramificado e $y$
o ponto fechado de $Y$. Seja $A'$ o
fecho integral de $A$; ponhamos
$Y'=\operatorname{Spec}(A')$ e seja
$y'$ o ponto fechado de $Y'$. Assim,
o morfismo $g:Y'\to Y$ é radicial
no ponto $y$ (6.15.3), inteiro e
birracional. Se pusermos
$X'=X\times_Y Y'$ e se
$f':X'\to Y'$ e $g':X'\to X$
forem as projeções, então $g'$ é
inteiro e radicial no ponto $x$
(6.15.3.1). Seja, portanto, $x'$
o único ponto de $g'^{-1}(x)$
situado acima de $y'$. O morfismo
$f'$ é de apresentação finita e
normal no ponto $x'$ (6.7.8);
por conseguinte, o anel local
$\sh{O}_{X',x'}$ é um domínio de integridade
integralmente fechado
\sref{IV.11.3.13}. Por outro lado,
podemos supor que $f$ é plano
\sref{IV.11.3.1}; logo, $g'$ é um
morfismo birracional (6.15.4.1).
Consequentemente,
$\operatorname{Spec}(\sh{O}_{X,x})$
é irredutível e, como
$\sh{O}_{X,x}$ é reduzido
\sref{IV.11.3.13}, esse anel é
um domínio de integridade geometricamente
unirramificado em virtude de
(6.15.5).
\end{proof}

\begin{proposition}[11.3.15]
\label{IV.11.3.15}
Sejam $A$ um anel, $B$ uma
$A$-álgebra de apresentação finita
e $M$ um $B$-módulo de apresentação
finita que é plano como $A$-módulo.
Então existe uma sequência finita
$(f_i)_{1\leq i\leq n}$ de elementos
de $A$ tal que o ideal gerado pelos
$f_i$ é igual a $A$ e, para
$0\leq i<n$,
\[
M_{f_{i+1}}\Big/
\Big(\sum_{j=1}^{i}f_jM_{f_{i+1}}\Big)
\]
é um módulo livre sobre
$A_{f_{i+1}}\big/
\big(\sum_{j=1}^{i}f_jA_{f_{i+1}}\big)$.
\end{proposition}

Podemos dizer também que os $D(f_i)$
formam um recobrimento aberto de
$X=\operatorname{Spec}(A)$ e que, se
pusermos $Z_1=D(f_1)$ e, em seguida,
por recorrência,
\[
Z_{i+1}=D(f_{i+1})\cap
V(f_1A+\cdots+f_iA),
\]
os $Z_i$ formam uma \emph{partição}
de $X$ em conjuntos localmente
fechados, sendo
$A_{f_{i+1}}/
(\sum_{j=1}^{i}f_jA_{f_{i+1}})$
o anel de um subesquema afim de
$X$ cujo espaço subjacente é
$Z_{i+1}$.

\begin{proof}
Para demonstrar a proposição, podemos
primeiro, em virtude de
\sref{IV.11.2.7}, supor que existem
um subanel noetheriano $A_0$ de $A$,
uma $A_0$-álgebra de tipo finito $B_0$
e um $B_0$-módulo de tipo finito $M_0$,
plano sobre $A_0$, tais que
$B=B_0\otimes_{A_0}A$ e
$M=M_0\otimes_{A_0}A$. É claro que
basta demonstrar a proposição para
$A_0$, $B_0$ e $M_0$, pois, se os
elementos $f_i\in A_0$ satisfizerem,
nesse caso, as condições do enunciado,
também as satisfarão para $A$, $B$ e
$M$, visto que
\[
M_{f_{i+1}}\Big/
\Big(\sum_{j=1}^{i}f_jM_{f_{i+1}}\Big)
=\left((M_0)_{f_{i+1}}\Big/
\Big(\sum_{j=1}^{i}
f_j(M_0)_{f_{i+1}}\Big)\right)
\otimes_{A_0}A
\]
(Bourbaki, \emph{Alg. comm.}, cap.~II,
\S~2, n\textsuperscript{o}~7,
prop.~18). Podemos, portanto,
limitar-nos daqui em diante ao caso
em que $A$ é noetheriano.

Notemos agora que, se $C$ é um anel
noetheriano, $\mathfrak N$ seu
nilradical e $P$ um $C$-módulo
plano, então, por
$(\mathbf 0_{\mathrm{III}},10.1.2)$,
para que $P$ seja um $C$-módulo livre,

\oldpage[IV-3]{142}

é necessário e suficiente que
$P\otimes_C(C/\mathfrak N)$ seja um
$(C/\mathfrak N)$-módulo livre.
Notemos, por outro lado, que, se
$C$ é um anel noetheriano reduzido,
existe $g\in C$ tal que $C_g$ é
um domínio de integridade. Utilizemos agora
o lema (6.9.2): podemos definir,
por recorrência, uma sequência
$(f_i)_{i\geq1}$ de elementos de
$A$ da seguinte maneira:
\begin{enumerate}
\item[$1^\circ$] $f_1$ é tal que
$A_1=(A_{\mathrm{red}})_{f_1}$
é um domínio de integridade e $M\otimes_A A_1$
é um $A_1$-módulo livre;
\item[$2^\circ$] se o ideal
$\mathfrak J_i$ gerado por
$f_1,\ldots,f_i$ é $A$, então
$f_{i+1}=f_i$;
\item[$3^\circ$] se, ao contrário,
$\mathfrak J_i\neq A$, então
$f_{i+1}$ é um elemento que não
pertence a $\mathfrak J_i$ e tal que
$A_{i+1}=
((A/\mathfrak J_i)_{\mathrm{red}})
_{f_{i+1}}$ é um domínio de integridade e
$M\otimes_A A_{i+1}$ é um
$A_{i+1}$-módulo livre.
\end{enumerate}
Como $A$ é noetheriano, a sequência
crescente $(\mathfrak J_i)$ é
estacionária; portanto, existe $n$
tal que $f_1,\ldots,f_n$ geram o
ideal $A$, e os $f_i$ satisfazem as
condições pedidas, pois
\[
((A/\mathfrak J_i)_{\mathrm{red}})
_{f_{i+1}}
=((A/\mathfrak J_i)_{f_{i+1}})
_{\mathrm{red}}.
\]
\end{proof}

\begin{proposition}[11.3.16]
\label{IV.11.3.16}
Sejam $f:X\to Y$ um morfismo
fielmente plano de apresentação
finita e $g:Y\to Z$ um morfismo
tal que $g\circ f:X\to Z$ é de
tipo finito (resp.\ de apresentação
finita). Então $g$ é de tipo finito
(resp.\ de apresentação finita).
\end{proposition}

\begin{proof}
Como $f$ é sobrejetivo e $g\circ f$
é quase-compacto, $g$ é quase-compacto
(1.1.3). Mostremos primeiro que,
se $g\circ f$ é de apresentação
finita, então $g$ é quase-separado.
Com efeito, consideremos um aberto
afim $W\subset Z$ e seja $(V_i)$
um recobrimento aberto afim finito
de $g^{-1}(W)$; basta mostrar que
os $V_i\cap V_j$ são quase-compactos
(1.2.6 e 1.2.7). Para cada $i$,
seja $(U_{ih})$ um recobrimento aberto
afim finito de $f^{-1}(V_i)$; como
$g\circ f$ é de apresentação finita,
todos os $U_{ih}\cap U_{jk}$ são
quase-compactos.\footnote{Nota de fonte (p.~142): o impresso menciona apenas $f$ nesta justificativa. No caso da prova em discussão, usamos $g\circ f$, cuja apresentação finita torna quase-separada a pré-imagem do aberto afim $W$; a leitura impressa fica conservada nesta nota.} Ora, como $f$ é
sobrejetivo, $V_i\cap V_j$ é a união
das imagens $f(U_{ih}\cap U_{jk})$,
quando $h$ e $k$ variam; portanto,
é quase-compacto, pois $f$ é
contínuo.

A questão é, portanto, local em $Y$,
e podemos supor que
$Y=\operatorname{Spec}(B)$ e
$Z=\operatorname{Spec}(A)$ são afins,
sendo $X$ reunião finita de abertos
afins $X_i=\operatorname{Spec}(C_i)$,
em que os $C_i$ são $B$-álgebras
de tipo finito (resp.\ de apresentação
finita). Seja $X'$ a soma disjunta
dos $X_i$ e seja $p:X'\to X$ o
morfismo cuja restrição a cada
$X_i$ é a imersão aberta canônica.
O morfismo $p$ é fielmente plano
e de apresentação finita (1.6.5);
logo, $g\circ f\circ p$ é de tipo
finito (resp.\ de apresentação
finita) e $f\circ p$ é fielmente
plano e de apresentação finita.
Podemos, por conseguinte, supor
também que $X=\operatorname{Spec}(C)$
é afim, reduzindo-nos assim a
demonstrar o corolário seguinte.
\end{proof}

\begin{corollary}[11.3.17]
\label{IV.11.3.17}
Sejam $A$ um anel, $B$ uma
$A$-álgebra e $C$ uma $B$-álgebra
de apresentação finita que é
fielmente plana como $B$-módulo.
Se $C$ é uma $A$-álgebra de tipo
finito (resp.\ de apresentação
finita), então $B$ é uma
$A$-álgebra de tipo finito
(resp.\ de apresentação finita).
\end{corollary}

\begin{proof}
Suponhamos primeiro que $C$ seja
uma $A$-álgebra de tipo finito.
Seja $(B_\lambda)$ a família
filtrante crescente das
sub-$A$-álgebras de tipo finito
de $B$. Em virtude de (8.8.2),
existe um índice $\alpha$ tal que
$C=C_\alpha\otimes_{B_\alpha}B$,
em que $C_\alpha$ é uma
$B_\alpha$-álgebra de apresentação
finita. Além disso, por (8.10.5)
e \sref{IV.11.2.6}, podemos supor
que $C_\alpha$ é um $B_\alpha$-módulo
fielmente plano. Para
$\lambda\geq\alpha$, ponhamos
$C_\lambda=C_\alpha\otimes_{B_\alpha}
B_\lambda$; temos então
$C=C_\lambda\otimes_{B_\lambda}B$,
sendo $C_\lambda$ um
$B_\lambda$-módulo fielmente plano.
Como a aplicação $B_\lambda\to B$
é injetiva, deduzimos que também
$C_\lambda\to C$ é injetiva.
Além disso, como
$C=\varinjlim C_\lambda$, todo
elemento de $C$ pertence a algum
$C_\lambda$; como $C$ é uma
$A$-álgebra de tipo finito, existe
$\lambda$ tal que a aplicação
$C_\lambda\to C$

\oldpage[IV-3]{143}

é bijetiva.\footnote{Nota de fonte (p.~143): o impresso justifica esta conclusão dizendo que C é uma B-álgebra de tipo finito. A hipótese usada é que C é uma A-álgebra de tipo finito, cujos geradores pertencem todos a um mesmo estágio; conservamos a leitura impressa nesta nota.} Mas então a aplicação
$B_\lambda\to B$ é bijetiva por
planicidade fiel; portanto,
$B=B_\lambda$ é uma $A$-álgebra
de tipo finito.

Suponhamos agora que $C$ seja
uma $A$-álgebra de apresentação
finita. Pela primeira parte do
argumento, temos
\[
B=A[T_1,\ldots,T_n]/
\widetilde{\mathfrak J}
\]
para certo ideal
$\widetilde{\mathfrak J}$.
Seja $(\widetilde{\mathfrak J}_\lambda)$
a família filtrante dos ideais
de tipo finito de
$A[T_1,\ldots,T_n]$ contidos em
$\widetilde{\mathfrak J}$, de
modo que
$\widetilde{\mathfrak J}
=\varinjlim\widetilde{\mathfrak J}_\lambda$
e $B=\varinjlim B_\lambda$, com
$B_\lambda=A[T_1,\ldots,T_n]/
\widetilde{\mathfrak J}_\lambda$.
Aplicando, como acima, (8.8.2),
(8.10.5) e \sref{IV.11.2.6},
temos $C=C_\alpha\otimes_{B_\alpha}B$,
em que $C_\alpha$ é uma
$B_\alpha$-álgebra de apresentação
finita e um $B_\alpha$-módulo
fielmente plano. Ponhamos ainda
$C_\lambda=C_\alpha\otimes_{B_\alpha}
B_\lambda$ para
$\lambda\geq\alpha$; assim,
pela planicidade,
\[
C_\lambda=C_\alpha\Big/
\Big(C_\alpha\otimes_{B_\alpha}
(\widetilde{\mathfrak J}_\lambda/
\widetilde{\mathfrak J}_\alpha)\Big),
\]
e, do mesmo modo,
\[
C=C_\alpha\Big/
\Big(C_\alpha\otimes_{B_\alpha}
(\widetilde{\mathfrak J}/
\widetilde{\mathfrak J}_\alpha)\Big).
\]
Como, por hipótese, $C$ e
$C_\alpha$ são $A$-álgebras
de apresentação finita,
$C_\alpha\otimes_{B_\alpha}
(\widetilde{\mathfrak J}/
\widetilde{\mathfrak J}_\alpha)$
é um ideal de tipo finito de
$C_\alpha$ (1.4.4), e os
$C_\alpha\otimes_{B_\alpha}
(\widetilde{\mathfrak J}_\lambda/
\widetilde{\mathfrak J}_\alpha)$
identificam-se com ideais de
tipo finito de $C_\alpha$ cuja
união é
$C_\alpha\otimes_{B_\alpha}
(\widetilde{\mathfrak J}/
\widetilde{\mathfrak J}_\alpha)$.
Existe, por conseguinte,
$\lambda\geq\alpha$ tal que
\[
C_\alpha\otimes_{B_\alpha}
(\widetilde{\mathfrak J}/
\widetilde{\mathfrak J}_\alpha)
=C_\alpha\otimes_{B_\alpha}
(\widetilde{\mathfrak J}_\lambda/
\widetilde{\mathfrak J}_\alpha),
\]
donde
$\widetilde{\mathfrak J}
=\widetilde{\mathfrak J}_\lambda$
por planicidade fiel, o que
prova que $B$ é uma $A$-álgebra
de apresentação finita.
\end{proof}

\subsection{Descida da planicidade por
morfismos quaisquer: caso de um
pré-esquema de base artiniano}
\label{subsection:IV.11.4}

\begin{theorem}[11.4.1]
\label{IV.11.4.1}
Sejam $A$ um anel, $\mathfrak J$
um ideal nilpotente de $A$ e
$u_\lambda:A\to B_\lambda$
$(\lambda\in L)$ uma família de
homomorfismos de anéis tal que a
interseção dos núcleos dos
$u_\lambda$ se reduza a $0$.
Seja $M$ um $A$-módulo tal que,
para todo $\lambda\in L$,
$M\otimes_A B_\lambda$ seja um
$B_\lambda$-módulo livre e
$M\otimes_A(A/\mathfrak J)$ seja
um $(A/\mathfrak J)$-módulo livre.
Então $M$ é um $A$-módulo livre.
Se, além disso, o conjunto de
índices $L$ é finito, pode-se
substituir em todo o enunciado
anterior \emph{livre} por
\emph{plano}.
\end{theorem}

\begin{proof}
Em ambos os casos, bastará provar
que $M$ é um $A$-módulo
\emph{plano}: com efeito, quando
$M\otimes_A(A/\mathfrak J)$ é
um $(A/\mathfrak J)$-módulo livre,
resulta de
$(\mathbf 0_{\mathrm{III}},10.1.2)$
que $M$ é um $A$-módulo livre.

Usaremos o lema seguinte, que
generaliza
$(\mathbf 0_{\mathrm{III}},10.2.1)$:

\begin{lemma}[11.4.1.1]
\label{IV.11.4.1.1}
Seja $A$ um anel munido de uma
filtração finita
$(\mathfrak J_n)_{0\leq n\leq N+1}$,
com $\mathfrak J_0=A$ e
$\mathfrak J_{N+1}=0$. Seja $M$
um $A$-módulo munido da filtração
$(\mathfrak J_nM)_{0\leq n\leq N+1}$,
e designemos por
$\operatorname{gr}(A)$ e
$\operatorname{gr}(M)$ o anel e o
módulo graduados correspondentes.
Suponhamos que
$M\otimes_A(A/\mathfrak J_1)$
seja um $(A/\mathfrak J_1)$-módulo
plano e que o homomorfismo canônico
\begin{equation}
\operatorname{gr}_0(M)
\otimes_{\operatorname{gr}_0(A)}
\operatorname{gr}(A)
\longrightarrow\operatorname{gr}(M)
\label{IV.11.4.1.2}
\tag{11.4.1.2}
\end{equation}
seja injetivo. Então $M$ é um
$A$-módulo plano.
\end{lemma}

O homomorfismo canônico (11.4.1.2)
define-se do mesmo modo que o de
$(\mathbf 0_{\mathrm{III}},10.1.1)$:
em grau $n$, é o homomorfismo
\[
(M/\mathfrak J_1M)\otimes
(\mathfrak J_n/\mathfrak J_{n+1})
=(M\otimes\mathfrak J_n)/
(\operatorname{Im}
(M\otimes\mathfrak J_{n+1})
+\operatorname{Im}
(\mathfrak J_1M\otimes\mathfrak J_n))
\longrightarrow
\mathfrak J_nM/\mathfrak J_{n+1}M
\]
obtido do homomorfismo canônico
$M\otimes\mathfrak J_n\to
\mathfrak J_nM$ por passagem aos
quocientes. O lema demonstra-se
por recorrência em $N$, pois nada
há a demonstrar para $N=0$.

\oldpage[IV-3]{144}

As condições impostas a $M$
acarretam, pela hipótese de
recorrência, que
$M\otimes_A(A/\mathfrak J_N)$
é um $(A/\mathfrak J_N)$-módulo
plano. Notemos agora que
$\mathfrak J_N^2\subset
\mathfrak J_{N+1}=0$ se $N\geq1$,
e que
\[
(M/\mathfrak J_1M)\otimes
(\mathfrak J_N/\mathfrak J_{N+1})
=(M/\mathfrak J_NM)\otimes
(\mathfrak J_N/\mathfrak J_{N+1})
=(M/\mathfrak J_NM)\otimes
\mathfrak J_N;
\]
logo, o homomorfismo canônico
\[
(M/\mathfrak J_NM)\otimes\mathfrak J_N
\longrightarrow
\mathfrak J_NM
=\mathfrak J_NM/\mathfrak J_N^2M
\]
é injetivo. Aplicando
$(\mathbf 0_{\mathrm{III}},10.2.1)$
à filtração $\mathfrak J_N$-pré-ádica,
concluímos que $M$ é um $A$-módulo
plano.

Para aplicar esse lema a
\sref{IV.11.4.1}, designaremos por
$\mathfrak J_n$ o ideal de $A$
obtido como interseção das imagens
inversas
$u_\lambda^{-1}(\mathfrak J^n
B_\lambda)$; é imediato que
$\mathfrak J_0=A$ e
$\mathfrak J_m\mathfrak J_n\subset
\mathfrak J_{m+n}$ para
$m\geq0$, $n\geq0$. Além disso,
se $\mathfrak J^{N+1}=0$, temos
também $\mathfrak J_{N+1}=0$,
pois a interseção dos núcleos dos
$u_\lambda$ se reduz a $0$.

Munamos $A$ da filtração
$(\mathfrak J_n)$ e $M$ da
filtração $(\mathfrak J_nM)$; para
cada $\lambda$, munamos
$B_\lambda$ e
$N_\lambda=M\otimes_A B_\lambda$
das filtrações $\mathfrak J$-pré-ádicas.
Consideremos, para cada $\lambda$,
o diagrama comutativo
\[
\xymatrix{
\operatorname{gr}^0(M)
\otimes_{\operatorname{gr}^0(A)}
\operatorname{gr}(A)
 \ar[r]^{\varphi} \ar[d]_{f_\lambda}
 & \operatorname{gr}(M)
 \ar[d]^{g_\lambda} \\
\operatorname{gr}_{\mathfrak J}^0(N_\lambda)
\otimes_{\operatorname{gr}_{\mathfrak J}^0(B_\lambda)}
\operatorname{gr}_{\mathfrak J}(B_\lambda)
 \ar[r]_{\varphi_\lambda}
 & \operatorname{gr}_{\mathfrak J}(N_\lambda)
}
\]
em que as setas horizontais são
os homomorfismos canônicos
(11.4.1.2), e as setas verticais
se deduzem dos homomorfismos
canônicos $A\to B_\lambda$ e
$M\to N_\lambda$. A hipótese de
que $N_\lambda$ é um
$B_\lambda$-módulo plano implica
que $\varphi_\lambda$ é injetivo
$(\mathbf 0_{\mathrm{III}},10.2.1)$;
portanto,
$\operatorname{Ker}(\varphi)\subset
\operatorname{Ker}(f_\lambda)$.
Pondo $A_0=A/\mathfrak J_1$ e
$M_0=M\otimes_A A_0$, notemos que
\[
\operatorname{gr}_{\mathfrak J}^0
(N_\lambda)
=N_\lambda/\mathfrak JN_\lambda
=(M/\mathfrak JM)\otimes_A B_\lambda
=(M/\mathfrak J_1M)\otimes_A
(B_\lambda/\mathfrak JB_\lambda),
\]
pois esse último produto tensorial
é igual a
\[
(M\otimes_A B_\lambda)\Big/
\big(\operatorname{Im}
(M\otimes_A\mathfrak JB_\lambda)
+\operatorname{Im}
(\mathfrak J_1M\otimes_A B_\lambda)\big),
\]
e temos
$\operatorname{Im}
(\mathfrak J_1M\otimes_A B_\lambda)
=\operatorname{Im}
(M\otimes_A\mathfrak J_1B_\lambda)
\subset\operatorname{Im}
(M\otimes_A\mathfrak JB_\lambda)$,
pois $\mathfrak J_1B_\lambda
\subset\mathfrak JB_\lambda$
por definição. Por fim, essa
inclusão mostra também que
\[
(M/\mathfrak J_1M)\otimes_A
(B_\lambda/\mathfrak JB_\lambda)
=(M/\mathfrak J_1M)
\otimes_{A/\mathfrak J_1}
(B_\lambda/\mathfrak JB_\lambda);
\]
portanto, obtemos
\[
\operatorname{gr}_{\mathfrak J}^0
(N_\lambda)
=M_0\otimes_{A_0}
(B_\lambda/\mathfrak JB_\lambda)
\]
e, por conseguinte,
\[
\operatorname{gr}_{\mathfrak J}^0
(N_\lambda)
\otimes_{\operatorname{gr}_{\mathfrak J}^0(B_\lambda)}
\operatorname{gr}_{\mathfrak J}
(B_\lambda)
=M_0\otimes_{A_0}
\operatorname{gr}_{\mathfrak J}
(B_\lambda).
\]
O homomorfismo $f_\lambda$
pode, portanto, escrever-se
\[
1\otimes\operatorname{gr}(u_\lambda):
M_0\otimes_{A_0}\operatorname{gr}(A)
\longrightarrow
M_0\otimes_{A_0}
\operatorname{gr}_{\mathfrak J}(B_\lambda).
\]

\oldpage[IV-3]{145}

Como $M_0$ é, por hipótese, um
$A_0$-módulo \emph{plano}, o núcleo
de $f_\lambda$ é igual a
$M_0\otimes_{A_0}R_\lambda$,
onde $R_\lambda$ é o núcleo de
$\operatorname{gr}(u_\lambda):
\operatorname{gr}(A)\to
\operatorname{gr}_{\mathfrak J}
(B_\lambda)$. Tudo se reduz,
portanto, a demonstrar que
\[
\bigcap_{\lambda\in L}
(M_0\otimes_{A_0}R_\lambda)=0.
\]
Ora, pela definição dos
$\mathfrak J_n$, a interseção dos
núcleos dos homomorfismos
\[
\mathfrak J_nA/\mathfrak J_{n+1}A
\longrightarrow
\mathfrak J^nB_\lambda/
\mathfrak J^{n+1}B_\lambda,
\]
quando $\lambda$ percorre $L$,
se reduz a $0$; em outras
palavras,
$\bigcap_{\lambda\in L}R_\lambda=0$.
Suponhamos primeiro que $M_0$
seja um $A_0$-módulo livre.
Escolhendo uma base de $M_0$,
vemos imediatamente que
\[
M_0\otimes_{A_0}
\Big(\bigcap_{\lambda\in L}R_\lambda\Big)
=\bigcap_{\lambda\in L}
(M_0\otimes_{A_0}R_\lambda),
\]
donde se obtém a afirmação
nesse caso. Quando $L$ é finito,
a igualdade anterior continua
válida sob a única hipótese de
que $M_0$ seja um $A_0$-módulo
plano $(\mathbf 0_{\mathrm I},6.1.3)$,
o que conclui a demonstração.
\end{proof}

\begin{remark}[11.4.2]
\label{IV.11.4.2}
A conclusão de \sref{IV.11.4.1}
pode falhar se $L$ for infinito
e se se supuser apenas que
$M\otimes_A(A/\mathfrak J)$ é um
$(A/\mathfrak J)$-módulo plano.
Por exemplo, sejam $V$ um anel
de valoração discreta, $K$ seu
corpo de frações e
$A=V[T]/(T^2)$ ($T$ uma
indeterminada). Tomemos por
$\mathfrak J$ a imagem de $(T)$
em $A$, de modo que
$A/\mathfrak J=V$, e tomemos
$M=K$, que é um
$(A/\mathfrak J)$-módulo e,
portanto, é igual a
$M\otimes_A(A/\mathfrak J)$.
Além disso, $M$ é plano sobre
$A/\mathfrak J$, mas não sobre
$A$: a inclusão
$\mathfrak J\hookrightarrow A$
torna-se, após tensorização por $K$,
o mapa nulo
$\mathfrak J\otimes_A K\simeq K\longrightarrow K$,
que não é injetivo, pois $K\neq0$.\footnote{Nota de fonte (p.~145): o impresso deduziria da planicidade de $M$ que ele seria livre, invocando $(\mathbf 0_{\mathrm{III}},10.1.2)$. Aqui $M=K$ não é livre sobre $A/\mathfrak J=V$, por isso essa justificativa não se aplica. Substituímos somente o argumento de não planicidade pela falha de injetividade após tensorização; o exemplo e a leitura impressa são preservados.}
Consideremos, por outro lado,
o ideal maximal $\mathfrak m$
do anel local noetheriano $A$.
Temos $\mathfrak mK=K$ e, assim,
$M\otimes_A(A/\mathfrak m^n)=0$
para todo inteiro $n$.
Os $(A/\mathfrak m^n)$-módulos
$M\otimes_A(A/\mathfrak m^n)$
são, portanto, planos para todo
$n$, e a interseção dos
$\mathfrak m^n$ se reduz a $0$.
\end{remark}

\begin{corollary}[11.4.3]
\label{IV.11.4.3}
Sejam $A$ um anel semilocal cujo
radical de Jacobson $\mathfrak J$ é
nilpotente (por exemplo, um
anel artiniano) e
$u_\lambda:A\to B_\lambda$
$(\lambda\in L)$ uma família de
homomorfismos de anéis tal que
a interseção dos núcleos dos
$u_\lambda$ se reduz a $0$.
Um $A$-módulo $M$ é plano se, e
somente se, para todo
$\lambda\in L$,
$M\otimes_A B_\lambda$ é um
$B_\lambda$-módulo plano.
\end{corollary}

\begin{proof}
Como $A/\mathfrak J$ é um
produto direto de um número
finito de corpos (Bourbaki,
\emph{Alg. comm.}, cap.~II,
\S~3, n$^\circ$~5, prop.~16)
e $\mathfrak J$ é nilpotente,
$A$ é um produto direto de um
número finito de anéis locais
$A_i$ cujo radical de Jacobson é nilpotente
(\emph{loc. cit.}, \S~4,
n$^\circ$~3, cor.~1 da
prop.~15). Por conseguinte,
$M$ é uma soma direta de
$A_i$-módulos $M_i$, cada
$M_i$ sendo anulado pelos
$A_j$ de índice $j\neq i$.
O módulo $M$ é plano sobre
$A$ se, e somente se, cada
$M_i$ é plano sobre $A_i$.
Além disso, a interseção dos
núcleos dos homomorfismos
\[
A_i\longrightarrow
A\xrightarrow{u_\lambda}B_\lambda
\]
se reduz a $0$, e
$M_i\otimes_{A_i}B_\lambda$
é um somando direto de
$M\otimes_A B_\lambda$.
Podemos, portanto, nos
restringir ao caso em que $A$
também é local. Então
$A/\mathfrak J$ é um corpo,
de modo que
$M\otimes_A(A/\mathfrak J)$
é um $(A/\mathfrak J)$-módulo
livre; em virtude de
\sref{IV.11.4.1}, basta mostrar
que, para todo $\lambda$,
$M\otimes_A B_\lambda$ é um
$B_\lambda$-módulo livre.
Ora, pondo
$\mathfrak J_\lambda=
\mathfrak JB_\lambda$,
\[
(M\otimes_A B_\lambda)
\otimes_{B_\lambda}
(B_\lambda/\mathfrak J_\lambda)
=
(M\otimes_A(A/\mathfrak J))
\otimes_{A/\mathfrak J}
(B_\lambda/\mathfrak J_\lambda)
\]
é um
$(B_\lambda/\mathfrak J_\lambda)$-módulo
livre, e $\mathfrak J_\lambda$
é nilpotente. A conclusão
decorre, pois, da hipótese
de que $M\otimes_A B_\lambda$
é um $B_\lambda$-módulo plano
e de
$(\mathbf 0_{\mathrm{III}},10.1.2)$.
\end{proof}

\begin{corollary}[11.4.4]
\label{IV.11.4.4}
Sejam $A$ um anel e $M$ um
$A$-módulo. Suponhamos que
exista um ideal nilpotente
$\mathfrak J$ de $A$ tal que
$M\otimes_A(A/\mathfrak J)$
seja um $(A/\mathfrak J)$-módulo
livre. Então o conjunto dos
ideais $\mathfrak K$ de $A$
para os quais
$M\otimes_A(A/\mathfrak K)$
é um $(A/\mathfrak K)$-módulo
livre admite um menor elemento
$\mathfrak K_0$, que também
é o menor dos ideais
$\mathfrak K$ para os quais
$M\otimes_A(A/\mathfrak K)$
é um $(A/\mathfrak K)$-módulo
plano.
\end{corollary}

\oldpage[IV-3]{146}

\begin{proof}
Para que um homomorfismo
$u:A\to A'$ seja tal que
$M\otimes_A A'$ seja um
$A'$-módulo livre
(respectivamente, plano), é
necessário e suficiente que
$u$ se fatorize como
$A\to A/\mathfrak K_0\to A'$
(ou, equivalentemente, que
$\mathfrak K_0A'=0$).

O fato de a interseção
$\mathfrak K_0$ dos ideais
$\mathfrak K$ para os quais
$M\otimes_A(A/\mathfrak K)$
é um $(A/\mathfrak K)$-módulo
livre ser o menor desses ideais
decorre de \sref{IV.11.4.1}
aplicado ao anel
$A/\mathfrak K_0$, ao seu ideal
nilpotente
$\mathfrak J/\mathfrak K_0$
e aos homomorfismos
$A/\mathfrak K_0\to
A/\mathfrak K$, cujos núcleos
têm interseção nula.
Se $A'$ é uma
$(A/\mathfrak K_0)$-álgebra,
temos
\[
M\otimes_A A'
=
(M\otimes_A(A/\mathfrak K_0))
\otimes_{A/\mathfrak K_0}A',
\]
de modo que $M\otimes_A A'$
é um $A'$-módulo livre.
Reciprocamente, se $A'$ é
uma $A$-álgebra tal que
$M\otimes_A A'$ é um
$A'$-módulo livre e se
$\mathfrak K$ é o núcleo do
homomorfismo $A\to A'$, segue
de \sref{IV.11.4.1}, aplicado
ao anel $A/\mathfrak K$, ao
$(A/\mathfrak K)$-módulo
$M\otimes_A(A/\mathfrak K)$,
ao ideal nilpotente
$(\mathfrak J+\mathfrak K)/
\mathfrak K$ e ao homomorfismo
injetivo $A/\mathfrak K\to A'$,
que
$M\otimes_A(A/\mathfrak K)$
é um $(A/\mathfrak K)$-módulo
livre; portanto,
$\mathfrak K\supset
\mathfrak K_0$.
O fato de podermos substituir
a condição de ser livre pela
de ser plano no que precede
(conservando, naturalmente, a
hipótese de que
$M\otimes_A(A/\mathfrak J)$
é um $(A/\mathfrak J)$-módulo
livre) decorre de
$(\mathbf 0_{\mathrm{III}},10.1.2)$,
como vimos no início da
demonstração de
\sref{IV.11.4.1}.
\end{proof}

\begin{proposition}[11.4.5]
\label{IV.11.4.5}
Sejam $Y$ um pré-esquema
irredutível,
$f:X\to Y$ um morfismo de
apresentação finita e
$\sh{F}$ um
$\sh{O}_X$-módulo de
apresentação finita.
Existem então um aberto não
vazio $U$ em $Y$ e um
subesquema fechado $Z$ de
$U$, de apresentação finita
sobre $U$, com a seguinte
propriedade: para toda
mudança de base $U'\to U$,
se pusermos
$X'=X\times_YU'$,
$f'=f_{(U')}$ e
$\sh{F}'=\sh{F}
\otimes_{\sh{O}_Y}\sh{O}_{U'}$,
então $\sh{F}'$ é
$f'$-plano se, e somente se,
o morfismo $U'\to U$ se
fatoriza como
$U'\to Z\to U$.
Tal esquema $Z$ tem o mesmo
espaço subjacente que $U$.
Suponhamos, além disso, que
$Y$ seja afim, e seja
$(W_i)$ uma cobertura finita
de $X$ por abertos afins.
Podemos então escolher $U$ de
modo que, se
$U'=\operatorname{Spec}(A')$
é um esquema afim e $U'\to U$
é um morfismo que se
fatoriza como $U'\to Z\to U$,
cada
$\Gamma(W_i\times_YU',\sh{F}')$
seja um $A'$-módulo livre.
\end{proposition}

\begin{proof}
Podemos evidentemente nos
restringir ao caso em que
$Y=\operatorname{Spec}(A)$
é afim. Utilizando
\sref{IV.8.9.1}, existem um
subanel noetheriano $A_1$
de $A$, um morfismo de tipo
finito
$f_1:X_1\to
Y_1=\operatorname{Spec}(A_1)$
e um $\sh{O}_{X_1}$-módulo
coerente $\sh{F}_1$ tais que
$f=(f_1)_{(Y)}$ e
$\sh{F}=\sh{F}_1
\otimes_{\sh{O}_{Y_1}}\sh{O}_Y$.
Podemos ainda supor que os
$W_i$ sejam imagens inversas
de abertos afins de $X_1$.
Notemos também que $Y_1$ é
irredutível, pois o morfismo
$Y\to Y_1$ é dominante
$(\mathbf I,1.2.7)$.
Suponhamos demonstrada a
proposição para $Y_1$,
$X_1$ e $\sh{F}_1$, e sejam
$U_1$ o aberto de $Y_1$ e
$Z_1$ o subesquema fechado
de $U_1$ com as propriedades
desejadas. Sejam
$U=U_1\times_{Y_1}Y$ e
$Z=Z_1\times_{Y_1}Y$ suas
imagens inversas.
Então, se $U'\to U$ é uma
mudança de base tal que
$\sh{F}'=
\sh{F}\otimes_{\sh{O}_Y}
\sh{O}_{U'}=
\sh{F}_1\otimes_{\sh{O}_{Y_1}}
\sh{O}_{U'}$
é $f'$-plano, o morfismo
$U'\to U_1$ se fatoriza como
$U'\to Z_1\to U_1$; mas,
como $U'\to U_1$ também se
fatoriza como
$U'\to U\to U_1$, a
definição do produto de
pré-esquemas mostra que
$U'\to U$ se fatoriza como
$U'\to Z\to U$.

Podemos, portanto, nos
restringir ao caso em que
$A$ é \emph{noetheriano}.
Seja $\mathfrak N$ seu
nilradical, que é nilpotente,
e ponhamos
$A_0=A_{\mathrm{red}}=
A/\mathfrak N$,
$Y_0=\operatorname{Spec}(A_0)
=Y_{\mathrm{red}}$,
$X_0=X\times_YY_0$,
$\sh{F}_0=\sh{F}
\otimes_{\sh{O}_Y}\sh{O}_{Y_0}$
e $W_{i0}=W_i\times_YY_0$.
Se $M_i=\Gamma(W_i,\sh{F})$,
então
$M_{i0}=
\Gamma(W_{i0},\sh{F}_0)
=M_i\otimes_A A_0$.
Como $A_0$ é um
\emph{domínio de integridade}, pelo teorema
da planicidade genérica
\sref{IV.6.9.2} podemos,
substituindo $Y$ por um aberto
não vazio se necessário,
supor que os $M_{i0}$ sejam
$A_0$-módulos livres.
Em virtude de
\sref{IV.11.4.4}, existe,
portanto, para cada $i$,
um ideal
$\mathfrak J_i\subset
\mathfrak N$ tal

\oldpage[IV-3]{147}

que as $A$-álgebras $A'$
para as quais os
$M_i\otimes_A A'$ são
$A'$-módulos livres (ou
planos) são exatamente
aquelas para as quais
$\mathfrak J_iA'=0$.
Seja
$\mathfrak J=
\sum_i\mathfrak J_i$,
um ideal contido em
$\mathfrak N$.
Dizer que
$\sh{F}\otimes_A A'$ é
$A'$-plano equivale a dizer
que todos os
$M_i\otimes_A A'$ são
$A'$-módulos planos, isto é,
que $\mathfrak JA'=0$.
Segue-se que, tomando
$Z=V(\mathfrak J)$,
resolvemos a questão, pois,
para que um morfismo
$U'\to U$ tenha a
propriedade do enunciado,
é necessário e suficiente
que, para todo aberto afim
$V'\subset U'$, o morfismo
$V'\to U$ tenha a mesma
propriedade.
\end{proof}

\begin{corollary}[11.4.6]
\label{IV.11.4.6}
Sejam $A$ um anel tal que
$\operatorname{Spec}(A)$ é
irredutível, $\mathfrak p$
seu único ideal primo minimal,
$B$ uma $A$-álgebra de
apresentação finita e $M$ um
$B$-módulo de apresentação
finita. Suponhamos que
$M_\mathfrak p$ seja um
$A_\mathfrak p$-módulo plano.
Então existe
$t\in A-\mathfrak p$ tal que
$M_t$ é um $A_t$-módulo
livre.
\end{corollary}

\begin{proof}
Aplicando \sref{IV.11.4.5}
com
$Y=\operatorname{Spec}(A)$,
$X=\operatorname{Spec}(B)$
e $\sh{F}=\widetilde M$,
podemos, substituindo $A$ por
$A_h$ se necessário, com
$h\in A-\mathfrak p$, supor
que exista um ideal de tipo
finito $\mathfrak J$ de $A$
tal que as $A$-álgebras $A'$
para as quais
$M\otimes_A A'$ é um
$A'$-módulo livre (ou plano)
são exatamente aquelas para
as quais $\mathfrak JA'=0$.
Em particular, a hipótese de
que $M_\mathfrak p$ é plano
sobre $A_\mathfrak p$
implica
$\mathfrak JA_\mathfrak p=0$,
ou seja,
$\mathfrak J_\mathfrak p=0$.
Como $\mathfrak J$ é de tipo
finito, existe
$t\in A-\mathfrak p$ tal que
$\mathfrak J_t=0$, ou,
equivalentemente,
$\mathfrak JA_t=0$; logo,
$M_t$ é um $A_t$-módulo
livre.
\end{proof}

\begin{proposition}[11.4.7]
\label{IV.11.4.7}
Sejam $A$ um anel
noetheriano, $\mathfrak J$
um ideal de $A$ e $M$ um
$A$-módulo idealmente
separado
$(\mathbf 0_{\mathrm{III}},10.2.1)$
para a topologia
$\mathfrak J$-pré-ádica.
Seja
$(\mathfrak p_i)_{1\leq i\leq r}$
uma família finita de ideais
primos de $A$ que contêm
$\mathfrak J$. Para todo
inteiro $n\geq0$, seja
$\mathfrak p_i^{(n)}$ a
$n$-ésima potência simbólica
de $\mathfrak p_i$, isto é, o
núcleo do homomorfismo
$A\to
A_{\mathfrak p_i}/
\mathfrak p_i^nA_{\mathfrak p_i}$.
Ponhamos
$\mathfrak J_n=
\bigcap_{i=1}^{r}
\mathfrak p_i^{(n)}$, de modo
que
$\mathfrak J_n\supset
\mathfrak J^n$, e suponhamos
que a topologia definida pela
filtração $(\mathfrak J_n)$
seja idêntica à topologia
$\mathfrak J$-pré-ádica
(em outras palavras, para
todo $n$ existe $m$ tal que
$\mathfrak J^n\supset
\mathfrak J_m$).
O módulo $M$ é plano sobre
$A$ se, e somente se,
$M\otimes_A(A/\mathfrak J)$
é plano sobre $A/\mathfrak J$
e, para todo $i$,
$M\otimes_A A_{\mathfrak p_i}
=M_{\mathfrak p_i}$
é plano sobre
$A_{\mathfrak p_i}$.
\end{proposition}

\begin{proof}
Como $M$ é idealmente
separado, basta, em virtude
de
$(\mathbf 0_{\mathrm{III}},10.2.1)$,
mostrar que, para todo
$n\geq0$,
$M\otimes_A(A/\mathfrak J^n)$
é plano sobre
$A/\mathfrak J^n$.
Como todo $\mathfrak J^n$
contém algum $\mathfrak J_m$,
isso equivale a provar que,
para todo $n$,
$M\otimes_A(A/\mathfrak J_n)$
é plano sobre
$A/\mathfrak J_n$.
Ora, como
$\mathfrak J_1\supset
\mathfrak J$,
$M\otimes_A(A/\mathfrak J_1)$
é plano sobre
$A/\mathfrak J_1$.
No anel $A/\mathfrak J_n$,
o ideal
$\mathfrak J_1/\mathfrak J_n$
é nilpotente; além disso,
a interseção dos núcleos
dos homomorfismos
\[
A/\mathfrak J_n
\longrightarrow
A_{\mathfrak p_i}/
\mathfrak p_i^n
A_{\mathfrak p_i}
\]
é nula em
$A/\mathfrak J_n$, pela
definição de
$\mathfrak J_n$.
Por \sref{IV.11.4.1},
basta, portanto, verificar
que
\[
M\otimes_A
(A_{\mathfrak p_i}/
\mathfrak p_i^n
A_{\mathfrak p_i})
\]
é plano sobre
$A_{\mathfrak p_i}/
\mathfrak p_i^n
A_{\mathfrak p_i}$;
isso decorre da hipótese
de que
$M\otimes_A A_{\mathfrak p_i}$
é plano sobre
$A_{\mathfrak p_i}$.
\end{proof}

\begin{remark}[11.4.8]
\label{IV.11.4.8}
A hipótese de
\sref{IV.11.4.7} sobre a
topologia definida pelos
$\mathfrak J_n$ é satisfeita
se, para todo $n$
suficientemente grande,
$\operatorname{Ass}
(A/\mathfrak J^n)$ está
contido no conjunto dos
$\mathfrak p_i$.
De fato, $\mathfrak J^n$
é então uma interseção de
ideais primários
correspondentes aos
$\mathfrak p_i$, cada um
dos quais contém uma
potência simbólica do
respectivo $\mathfrak p_i$;
daí a conclusão. Em
particular:
\end{remark}

\begin{corollary}[11.4.9]
\label{IV.11.4.9}
Sejam $A$ um anel
noetheriano, $\mathfrak J$
um ideal nilpotente de $A$
e $M$ um $A$-módulo.
O módulo $M$ é plano
(respectivamente, livre)
sobre $A$ se, e somente se,
$M\otimes_A(A/\mathfrak J)$
é plano (respectivamente,
livre) sobre $A/\mathfrak J$
e, para todo ideal primo
$\mathfrak p\in
\operatorname{Ass}(A)$,
$M_\mathfrak p$ é plano
sobre $A_\mathfrak p$.
\end{corollary}

\oldpage[IV-3]{148}

A afirmação relativa ao caso
em que
$M\otimes_A(A/\mathfrak J)$
é livre também decorre da
afirmação relativa ao caso
em que esse módulo é plano,
por
$(\mathbf 0_{\mathrm{III}},10.1.2)$.

\begin{corollary}[11.4.10]
\label{IV.11.4.10}
Sejam $A$ um anel
noetheriano, $\mathfrak J$
um ideal de $A$ e $M$ um
$A$-módulo. Suponhamos que
$M$ seja idealmente separado
para a topologia
$\mathfrak J$-pré-ádica e
que
$\operatorname{gr}_{\mathfrak J}(A)$
seja plano sobre
$A/\mathfrak J$.
O módulo $M$ é plano sobre
$A$ se, e somente se,
$M\otimes_A(A/\mathfrak J)$
é plano sobre $A/\mathfrak J$
e, para todo
$\mathfrak p\in
\operatorname{Ass}
(A/\mathfrak J)$,
$M_\mathfrak p$ é plano
sobre $A_\mathfrak p$.
\end{corollary}

\begin{proof}
Tendo em conta
\sref{IV.11.4.8}, basta
mostrar que
$\operatorname{Ass}
(A/\mathfrak J^n)$ está
contido em
$\operatorname{Ass}
(A/\mathfrak J)$ para
todo $n$. Ora, se
$a\in A$ não pertence a
nenhum dos
$\mathfrak p\in
\operatorname{Ass}
(A/\mathfrak J)$, a
multiplicação por $a$ é
injetiva em $A/\mathfrak J$.
Como cada
$\mathfrak J^k/
\mathfrak J^{k+1}$ é um
$(A/\mathfrak J)$-módulo
plano, $a$ também é um
elemento regular em
$\mathfrak J^k/
\mathfrak J^{k+1}$;
logo, $a$ é regular em
$A/\mathfrak J^n$ para
todo $n$ (Bourbaki,
\emph{Alg. comm.}, cap.~III,
\S~2, n$^\circ$~8,
cor.~1 do teorema~1).
Por conseguinte, $a$ não
pertence a nenhum ideal
primo associado a
$A/\mathfrak J^n$, donde
o corolário (Bourbaki,
\emph{Alg. comm.}, cap.~II,
\S~1, n$^\circ$~1,
prop.~2).
\end{proof}

\begin{proposition}[11.4.11]
\label{IV.11.4.11}
Sejam $A$ um anel local
artiniano, de ideal maximal
$\mathfrak m$,
$Y=\operatorname{Spec}(A)$,
$y$ o único ponto de $Y$,
$f:X\to Y$ um morfismo
localmente de tipo finito e
$\sh{F}$ um
$\sh{O}_X$-módulo coerente.
Sejam $(A'_\alpha)$ uma
família de anéis locais e,
para cada $\alpha$,
$u_\alpha:A\to A'_\alpha$
um homomorfismo de anéis
(necessariamente local).
Ponhamos
$Y'_\alpha=
\operatorname{Spec}(A'_\alpha)$,
$X'_\alpha=
X\times_Y Y'_\alpha$ e
$\sh{F}'_\alpha=
\sh{F}\otimes_A A'_\alpha$.
Seja $x$ um ponto de $X$,
e suponhamos satisfeitas as
seguintes condições:
\begin{enumerate}[label=(\roman*)]
\item A interseção dos
núcleos dos $u_\alpha$
se reduz a $0$.
\item A extensão $k(x)$
do corpo residual $k$ de
$A$ é primária
(4.3.1)
(condição automaticamente
satisfeita se $k$ é
separavelmente fechado).
\item Para todo $\alpha$,
existe um ponto
$x'_\alpha\in X'_\alpha$
cujas projeções respectivas
em $X$ e $Y'_\alpha$
são $x$ e o ponto fechado
$y'_\alpha$ de
$Y'_\alpha$, e tal que
$\sh{F}'_\alpha$ é
$Y'_\alpha$-plano no ponto
$x'_\alpha$.
\end{enumerate}
Nessas condições, $\sh{F}$
é $f$-plano no ponto $x$.
\end{proposition}

\begin{proof}
Como a questão é local em
$X$, podemos nos restringir
ao caso em que $f$ é um
morfismo de tipo finito e
supor $\sh{F}_x\neq0$.
Procederemos em várias
etapas.

\medskip
\noindent\textbf{I)}
\emph{Redução ao caso em
que $A'_\alpha$ é um anel
local de um pré-esquema de
tipo finito sobre $Y$ e
em que o corpo residual
$k'_\alpha$ de
$A'_\alpha$ é uma extensão
finita de $k$.}

\medskip
Como a redução se faz
separadamente para cada
$A'_\alpha$, podemos
omitir o índice $\alpha$
nesta parte. Seja
$\mathfrak m'$ o ideal
maximal de $A'$.
Consideremos $A'$ como
limite indutivo de suas
sub-$A$-álgebras de tipo
finito $B_\lambda$, e
ponhamos
$\mathfrak n_\lambda=
\mathfrak m'\cap B_\lambda$.
O anel $A'$ é também limite
indutivo de seus subanéis
locais
$B'_\lambda=
(B_\lambda)_{\mathfrak n_\lambda}$
(5.13.3), e temos
evidentemente
$\mathfrak m'\cap B'_\lambda
=\mathfrak n'_\lambda$,
ideal maximal de
$B'_\lambda$.
Existe, portanto, por
\sref{IV.11.2.6}, um
índice $\lambda$ tal que,
pondo
$Z'_\lambda=
\operatorname{Spec}(B'_\lambda)$,
$\sh{F}\otimes_A B'_\lambda$
é $Z'_\lambda$-plano no
ponto $x'_\lambda$,
projeção de $x'$, sendo
a projeção de
$x'_\lambda$ em
$Z'_\lambda$ o ponto
fechado $z'_\lambda$
de $Z'_\lambda$.

Podemos, portanto, supor
que exista um esquema afim
$Z'$ de tipo finito sobre
$Y$ e um ponto $z'$ de
$Z'$ tais que
$A'=\sh{O}_{Z',z'}$ e
que, pondo
$T'=X\times_Y Z'$,
exista um ponto
$t'\in T'$ cujas projeções
em $Z'$ e $X$ são $z'$
e $x$, e tal que
$\sh{F}\otimes_Y Z'$
é $Z'$-plano em $t'$.
Sejam $Z'_1$
(respectivamente, $X_1$)
um subpré-esquema fechado
de $Z'$ (respectivamente,
de $X$) cujo espaço
subjacente é o fecho de
$\{z'\}$
(respectivamente, de
$\{x\}$), e ponhamos
$T'_1=X_1\times_Y Z'_1$.
O conjunto $U$ dos pontos
de $T'$ em que
$\sh{F}\otimes_Y Z'$
é $Z'$-plano é aberto em
$T'$ \sref{IV.11.1.1} e
contém $t'$; portanto,

\oldpage[IV-3]{149}

$V=U\cap T'_1$ é um
aberto não vazio de
$T'_1$. Sendo artiniano,
$A$ é um anel de Jacobson;
logo, por
(10.4.6),
$Z'_1$ e $T'_1$ são
pré-esquemas de Jacobson.
Há, portanto, em $V$ um
ponto $t'_0$ fechado em
$V$, cuja imagem $z'_0$
em $Z'_1$ é um ponto
fechado de $Z'_1$
(10.4.7).
Sejam $f_1$ a restrição
de $f$ a $X_1$ e
$f'_1:T'_1\to Z'_1$ a
projeção canônica.
O conjunto
$V\cap f_1'^{-1}(z'_0)$
é um aberto não vazio de
$f_1^{-1}(y)\otimes_{k(y)}
k(z'_0)$; como este
último pré-esquema é
plano sobre
$f_1^{-1}(y)$, um ponto
maximal $t'_1$ de
$V\cap f_1'^{-1}(z'_0)$
está necessariamente
sobre o único ponto
maximal $x$ de $X_1$
(2.3.4).
Por fim, $k(z'_0)$ é
uma extensão finita
de $k$ $(\mathbf I,6.4.2)$,
e o homomorfismo
$A=\sh{O}_{Y,y}\to
A'=\sh{O}_{Z',z'}$
se fatoriza como
$\sh{O}_{Y,y}\to
\sh{O}_{Z',z'_0}\to
\sh{O}_{Z',z'}$;
seu núcleo contém,
portanto, o núcleo de
$\sh{O}_{Y,y}\to
\sh{O}_{Z',z'_0}$.
Isso conclui a redução
anunciada.

\medskip
\noindent\textbf{II)}
\emph{Redução ao caso em que
os $A'_\alpha$ são em número
finito e são $A$-álgebras
finitas.} --- Seja
$\mathfrak m'_\alpha$ o ideal
maximal de $A'_\alpha$.
Como $A'_\alpha$ é um anel
local noetheriano, a
interseção dos
$\mathfrak m'^n_\alpha$ é
$0$ $(\mathbf 0_{\mathrm{I}},7.3.5)$.
A interseção dos
$u_\alpha^{-1}(\mathfrak m'^n_\alpha)$,
para todos
os índices $n$ e $\alpha$,
é, portanto, igual à
interseção dos núcleos dos
$u_\alpha$ e, pela hipótese
(i), reduz-se a $0$.
Como $A$ é artiniano, há
um número finito desses
ideais cuja interseção já
é $0$; denotemo-los por
$u_i^{-1}(\mathfrak m_i'^{\,n_i})$
$(1\leq i\leq r)$.
Como os $A'_i$ são
noetherianos, os
$\mathfrak m_i'^{\,h}/
\mathfrak m_i'^{\,h+1}$
são módulos de comprimento
finito sobre
$A'_i/\mathfrak m'_i$; e,
como $A'_i/\mathfrak m'_i$
é um $k$-espaço vetorial
de dimensão finita, vê-se
que $A''_i=A'_i/
\mathfrak m_i'^{\,n_i}$ é
uma $A$-álgebra finita e
um anel local artiniano.
A redução anunciada decorre,
pois, de (2.1.4).

\medskip
\noindent\textbf{III)}
\emph{Conclusão da
demonstração.} --- Suponhamos
agora que os $A'_i$
$(1\leq i\leq r)$ sejam
em número finito e sejam
$A$-álgebras finitas.
Para cada $i$, o corpo
residual $k_i$ de $A'_i$
é uma extensão finita de
$k$; usando a hipótese
(ii), concluímos que a
imagem inversa de $x$ em
$X'_i$ se reduz ao único
ponto $x'_i$ (4.3.2).
Seja $Y'$ a soma dos
pré-esquemas $Y'_i$, seja
$X'=X\times_Y Y'$ a soma
dos $X'_i$, e ponhamos
$\sh{F}'=\sh{F}\otimes_Y Y'$.
A hipótese implica que
$\sh{F}'$ é $Y'$-plano em
todos os pontos da imagem
inversa de $x$ pela
projeção $p:X'\to X$.
Como $p$ é de tipo finito,
existe, portanto, um
aberto $U'\supset
p^{-1}(x)$ tal que
$\sh{F}'$ é $Y'$-plano em
todos os pontos de $U'$
\sref{IV.11.1.1}.
Além disso, o morfismo
$Y'\to Y$ é finito, pois
os $A'_i$ são $A$-álgebras
finitas; logo, $p$ é um
morfismo finito
$(\mathbf{II},6.1.5)$ e,
portanto, fechado
$(\mathbf{II},6.1.10)$.
Existe, assim, uma
vizinhança afim $U$ de
$x$ em $X$ tal que
$p^{-1}(U)\subset U'$.
Sejam $B$ e $A'$ os anéis
dos esquemas $U$ e $Y'$,
respectivamente (sendo
$A'$ o produto direto dos
$A'_i$). Substituindo
$X$ por $U$, temos
$\sh{F}=\widetilde M$,
onde $M$ é um $B$-módulo
e, por hipótese,
$M'=M\otimes_A A'$ é um
$A'$-módulo plano
(2.1.2). Como o
homomorfismo $A\to A'$ é
injetivo por construção,
podemos aplicar
\sref{IV.11.4.3}, o que
demonstra que $M$ é um
$A$-módulo plano.
\end{proof}

\begin{proposition}[11.4.12]
\label{IV.11.4.12}
Sejam $A$ um anel local
artiniano de corpo residual
$k$, $Y=\operatorname{Spec}(A)$,
$f:X\to Y$ um morfismo
localmente de tipo finito
e $\sh{F}$ um
$\sh{O}_X$-módulo coerente.
Sejam $(A'_\alpha)$ uma
família de anéis locais e,
para cada $\alpha$,
$u_\alpha:A\to A'_\alpha$
um homomorfismo de anéis.
Ponhamos
$Y'_\alpha=\operatorname{Spec}(A'_\alpha)$,
$X'_\alpha=X\times_Y Y'_\alpha$,
$f'_\alpha=f_{(Y'_\alpha)}$
e $\sh{F}'_\alpha=
\sh{F}\otimes_A A'_\alpha$.
Seja $x$ um ponto de $X$,
e suponhamos satisfeitas
as seguintes condições:
\begin{enumerate}[label=(\roman*)]
\item A interseção dos
núcleos dos $u_\alpha$
se reduz a $0$.
\item Para todo $\alpha$,
$\sh{F}'_\alpha$ é
$f'_\alpha$-plano em todos
os pontos
$x'_\alpha\in X'_\alpha$
cujas projeções respectivas
em $X$ e $Y'_\alpha$
são $x$ e o ponto fechado
$y'_\alpha$ de
$Y'_\alpha$.
\end{enumerate}
Então $\sh{F}$ é
$f$-plano no ponto $x$.
\end{proposition}

\begin{proof}
\oldpage[IV-3]{150}

Por hipótese, a interseção
dos núcleos dos $u_\alpha$
se reduz a $0$; como $A$
é artiniano, já há um
número finito desses
núcleos cuja interseção é
$0$. Podemos, portanto,
nos restringir ao caso em
que a família
$(A'_\alpha)$ é finita.
Seja $k'$ um fecho
algébrico de $k$; sabe-se
$(\mathbf{0}_{\mathrm{III}},
10.3.1)$ que existe um
homomorfismo local
$A\to B$ que faz de $B$
um $A$-módulo plano, tal
que $B$ seja um anel local
artiniano, integral sobre
$A$, e que $B\otimes_A k$
seja isomorfo a $k'$.

Pela planicidade, os
núcleos dos homomorfismos
$B\to B'_\alpha=
B\otimes_A A'_\alpha$
se obtêm dos núcleos dos
$u_\alpha$ por tensorização
com $B$ e, como são em
número finito, sua
interseção se reduz a $0$
$(\mathbf{0}_{\mathrm I},
6.1.3)$. Consideremos os
anéis $B'_{\alpha\beta}$,
localizações de $B'_\alpha$
em seus ideais maximais;
sabe-se (Bourbaki,
\emph{Alg. comm.}, cap.~II,
\S~3, n$^\circ$~3,
cor.~2 do teorema~1) que
a interseção dos núcleos
dos homomorfismos
$B'_\alpha\to
B'_{\alpha\beta}$, para
um $\alpha$ dado, se
reduz a $0$. Conclui-se,
portanto, que a interseção
dos núcleos dos
homomorfismos compostos
$v_{\alpha\beta}:B\to
B'_\alpha\to
B'_{\alpha\beta}$, com
$\alpha$ e $\beta$
variáveis, se reduz a $0$.
Por outro lado, como $B$
é integral sobre $A$,
$B'_\alpha$ é integral
sobre $A'_\alpha$; seus
ideais maximais estão,
portanto, sobre o ideal
maximal de $A'_\alpha$.
Pondo
$Z'_{\alpha\beta}=
\operatorname{Spec}(B'_{\alpha\beta})$,
$T'_{\alpha\beta}=
X'_\alpha\times_{Y'_\alpha}
Z'_{\alpha\beta}$,
$f'_{\alpha\beta}=
f_{(Z'_{\alpha\beta})}$,
$\sh{G}=\sh{F}\otimes_A B$
e $\sh{G}'_{\alpha\beta}=
\sh{G}\otimes_B
B'_{\alpha\beta}$, vemos
que as hipóteses (i) e
(ii) continuam satisfeitas
quando substituímos,
respectivamente, $A$,
$A'_\alpha$, $u_\alpha$,
$\sh{F}$ e $x$ por $B$,
$B'_{\alpha\beta}$,
$v_{\alpha\beta}$,
$\sh{G}$ e um ponto $t$
de $T=X\otimes_A B$
acima de $x$. Como o
corpo residual de $B$ é
separavelmente fechado,
deduzimos de
\sref{IV.11.4.11} que
$\sh{G}$ é plano sobre
$B$ no ponto $t$. Mas,
como $B$ é um $A$-módulo
fielmente plano, concluímos
de (2.1.4) que $\sh{F}$
é plano sobre $A$ no
ponto $x$, o que demonstra
(11.4.12).
\end{proof}

\begin{corollary}[11.4.13]
\label{IV.11.4.13}
Sejam $A$ um anel local
artiniano,
$Y=\operatorname{Spec}(A)$,
$f:X\to Y$ um morfismo
localmente de tipo finito
e $\sh{F}$ um
$\sh{O}_X$-módulo coerente.
Seja $(Y'_\alpha)$ uma
família de
$Y$-pré-esquemas e, para
todo $\alpha$, ponhamos
$X'_\alpha=X\times_Y Y'_\alpha$,
$f'_\alpha=f_{(Y'_\alpha)}$
e $\sh{F}'_\alpha=
\sh{F}\otimes_Y Y'_\alpha$.
Seja $x$ um ponto de $X$,
e suponhamos satisfeitas
as seguintes hipóteses:
\begin{enumerate}[label=(\roman*)]
\item A interseção dos
núcleos dos homomorfismos
$A\to\Gamma(Y'_\alpha,
\sh{O}_{Y'_\alpha})$
correspondentes aos
morfismos estruturais se
reduz a $0$.
\item Para todo $\alpha$,
$\sh{F}'_\alpha$ é
$f'_\alpha$-plano em todos
os pontos
$x'_\alpha\in X'_\alpha$
cuja projeção em $X$ é
$x$.
\end{enumerate}
Então $\sh{F}$ é
$f$-plano no ponto $x$.
\end{corollary}

\begin{proof}
Com efeito, para todo
$y'_\alpha\in Y'_\alpha$,
consideremos o esquema
local
$Y''_{\alpha,y'_\alpha}=
\operatorname{Spec}(
\sh{O}_{Y'_\alpha,y'_\alpha})$;
em virtude de (2.1.4),
$\sh{F}'_\alpha
\otimes_{Y'_\alpha}
Y''_{\alpha,y'_\alpha}$
é $Y''_{\alpha,y'_\alpha}$-plano
nos pontos cujas projeções
em $X$ e
$Y''_{\alpha,y'_\alpha}$
são $x$ e o ponto fechado
de $Y''_{\alpha,y'_\alpha}$.
Além disso, o núcleo do
homomorfismo
$A\to\Gamma(Y'_\alpha,
\sh{O}_{Y'_\alpha})$ é a
interseção dos núcleos dos
homomorfismos
$A\to\Gamma(
Y''_{\alpha,y'_\alpha},
\sh{O}_{Y''_{\alpha,y'_\alpha}})$,
pois nos reduzimos de
imediato ao caso em que
$Y'_\alpha$ é afim, e
basta então aplicar
Bourbaki, \emph{Alg. comm.},
cap.~II, \S~3, n$^\circ$~3,
cor.~2 do teorema~1.
Substituindo a família
$(Y'_\alpha)$ pela família
dos $Y''_{\alpha,y'_\alpha}$,
ficamos, portanto, reduzidos
a \sref{IV.11.4.12}.
\end{proof}

\subsection{Descida da planicidade por morfismos quaisquer: caso geral}
\label{subsection:IV.11.5}

\begin{theorem}[11.5.1]
\label{IV.11.5.1}
Sejam $Y$ um pré-esquema
localmente noetheriano,
$f:X\to Y$ um morfismo
localmente de tipo finito,
$\sh{F}$ um
$\sh{O}_X$-módulo coerente,
$x$ um ponto de $X$ e
$y=f(x)$. Seja
$(Y'_\alpha)$ uma família
de $Y$-pré-esquemas locais
$Y'_\alpha=\operatorname{Spec}(A'_\alpha)$
tal que as imagens dos
pontos fechados
$y'_\alpha$ dos
$Y'_\alpha$ sejam todas
iguais a $y$. Para todo
$\alpha$, sejam
$\mathfrak m'_\alpha$ o
ideal maximal de
$A'_\alpha$ e
$u_\alpha:\sh{O}_y\to
A'_\alpha$
\oldpage[IV-3]{151}
o homomorfismo canônico
$(\mathbf I,2.4.4)$;
suponhamos que as
interseções finitas dos
ideais
$u_\alpha^{-1}(
\mathfrak m'_\alpha{}^{\,n_\alpha})$
formem um sistema
fundamental de vizinhanças
de $0$ em $\sh{O}_y$.
Ponhamos
$X'_\alpha=X\times_Y Y'_\alpha$,
$f'_\alpha=f_{(Y'_\alpha)}$
e $\sh{F}'_\alpha=
\sh{F}\otimes_Y Y'_\alpha$,
e suponhamos satisfeita
uma das seguintes hipóteses:
\begin{enumerate}[label=(\roman*)]
\item Para todo $\alpha$,
$\sh{F}'_\alpha$ é
$f'_\alpha$-plano em todos
os pontos cuja projeção
em $X$ é igual a $x$ e
cuja projeção em
$Y'_\alpha$ é igual a
$y'_\alpha$.
\item Para todo $\alpha$,
existe um
$x'_\alpha\in X'_\alpha$
cuja projeção em $X$ é
$x$ e cuja projeção em
$Y'_\alpha$ é
$y'_\alpha$, tal que
$\sh{F}'_\alpha$ é
$f'_\alpha$-plano no ponto
$x'_\alpha$, e $k(x)$
é uma extensão primária
de $k(y)$.
\end{enumerate}
Nessas condições,
$\sh{F}$ é $f$-plano no
ponto $x$.
\end{theorem}

\begin{proof}
Seja $\mathfrak m_y$ o
ideal maximal de
$\sh{O}_y$; como
$\sh{O}_y$ e $\sh{O}_x$
são noetherianos e
$\sh{O}_y\to\sh{O}_x$ é
um homomorfismo local,
basta, em virtude de
$(\mathbf{0}_{\mathrm{III}},
10.2.2)$, provar que,
para todo $n>0$,
$\sh{F}_x/\mathfrak m_y^n
\sh{F}_x$ é um módulo
plano sobre
$\sh{O}_y/\mathfrak m_y^n$.
Designemos por
$(\mathfrak J_\lambda)$
a família das interseções
finitas dos
$u_\alpha^{-1}(
\mathfrak m'_\alpha{}^{\,n_\alpha})$;
por hipótese, existe
$\lambda$ tal que
$\mathfrak J_\lambda\subset
\mathfrak m_y^n$, e, como
\[
\sh{F}_x\otimes_{\sh{O}_y}
(\sh{O}_y/\mathfrak m_y^n)
=\bigl(\sh{F}_x
\otimes_{\sh{O}_y}
(\sh{O}_y/\mathfrak J_\lambda)\bigr)
\otimes_{\sh{O}_y/\mathfrak J_\lambda}
(\sh{O}_y/\mathfrak m_y^n),
\]
bastará provar que
$\sh{F}_x/\mathfrak J_\lambda
\sh{F}_x$ é um módulo
plano sobre
$\sh{O}_y/\mathfrak J_\lambda$.
Ora, para cada $\alpha$
tal que
$\mathfrak J_\lambda\subset
u_\alpha^{-1}(
\mathfrak m'_\alpha{}^{\,n_\alpha})$,
a passagem aos quocientes
dá um homomorfismo
$u'_\alpha:
\sh{O}_y/\mathfrak J_\lambda
\to A'_\alpha/
\mathfrak m'_\alpha{}^{\,n_\alpha}$,
e a interseção dos núcleos
dos $u'_\alpha$ se reduz
a $0$. Levando em conta
$(\mathbf I,3.6.1)$, vemos
que estamos reduzidos a
\sref{IV.11.4.11} no caso
(ii) e a
\sref{IV.11.4.12} no caso
(i).
\end{proof}

\begin{corollary}[11.5.2]
\label{IV.11.5.2}
Sejam $Y$ um pré-esquema
localmente noetheriano,
$f:X\to Y$ um morfismo
localmente de tipo finito,
$\sh{F}$ um
$\sh{O}_X$-módulo coerente,
$x$ um ponto de $X$ e
$y=f(x)$; ponhamos
$A=\sh{O}_y$. Seja
$u:A\to A'$ um
homomorfismo de $A$ em
um anel de Zariski $A'$
tal que a imagem inversa
$u^{-1}(\mathfrak r')$ do
radical $\mathfrak r'$ de
$A'$ seja o ideal maximal
$\mathfrak m$ de $A$;
suponhamos ainda que o
homomorfismo
$\widehat u:\widehat A\to
\widehat{A'}$ seja
injetivo. Ponhamos
$Y'=\operatorname{Spec}(A')$,
$X'=X\times_Y Y'$,
$f'=f_{(Y')}$ e
$\sh{F}'=\sh{F}\otimes_Y Y'$.
Para que $\sh{F}$ seja
$f$-plano no ponto $x$,
é necessário e suficiente
que $\sh{F}'$ seja
$f'$-plano em todos os
pontos cuja projeção em
$X$ seja $x$ e cuja
projeção em $Y'$ seja um
ponto fechado de $Y'$.

Se, além disso, $A'$ é uma
$A$-álgebra finita,
podemos substituir, no
enunciado anterior, a
hipótese de injetividade
de $\widehat u$ pela de
injetividade de $u$.
\end{corollary}

\begin{proof}
Como $A$ (respectivamente,
$A'$) se identifica com
um subanel de $\widehat A$
(respectivamente, de
$\widehat{A'}$)
(Bourbaki, \emph{Alg. comm.},
cap.~III, \S~3,
n$^\circ$~3, prop.~6),
vemos primeiro que o
próprio $u$ é injetivo e
que $\widehat u$ é seu
prolongamento por
continuidade a
$\widehat A$.

Seja $(\mathfrak m'_\alpha)$
a família dos ideais
maximais de $A'$; como
temos
\[
\mathfrak m'_\alpha{}^{\,n}
\widehat{A'}
=\widehat{\mathfrak m}'_\alpha{}^{\,n},
\qquad\text{e}\qquad
\widehat{\mathfrak m}'_\alpha{}^{\,n}
\cap A'
=\mathfrak m'_\alpha{}^{\,n},
\]
e como os
$\mathfrak m'_\alpha{}^{\,n}$
são abertos em $A'$,
temos
$u^{-1}(\mathfrak m'_\alpha{}^{\,n})
=A\cap\widehat u^{-1}(
\widehat{\mathfrak m}'_\alpha{}^{\,n})$.
Bastará provar que, em
$\widehat A$, as
interseções finitas dos
$\widehat u^{-1}(
\widehat{\mathfrak m}'_\alpha{}^{\,n_\alpha})$
formam um sistema
fundamental de vizinhanças
de $0$. Isso permitirá
aplicar \sref{IV.11.5.1}
aos homomorfismos
compostos
$v_\alpha\circ u:A\to
A'_{\mathfrak m'_\alpha}$,
em que $v_\alpha:A'\to
A'_{\mathfrak m'_\alpha}$
é o homomorfismo canônico.
Como $\widehat A$ é
completo, bastará provar
que a interseção dos
$\widehat u^{-1}(
\widehat{\mathfrak m}'_\alpha{}^{\,n_\alpha})$
se reduz a $0$
(Bourbaki, \emph{Alg. comm.},
cap.~III, \S~2,
n$^\circ$~7, prop.~8,
em cuja demonstração se
pode substituir a sequência
decrescente por um conjunto
filtrante qualquer). Ora,
para cada $\alpha$
\oldpage[IV-3]{152}
fixo, a interseção dos
ideais
$\widehat{\mathfrak m}'_\alpha{}^{\,n}
\widehat{A'}_{\widehat{\mathfrak m}'_\alpha}$,
para $n>0$, se reduz a
$0$ no anel local
noetheriano
$\widehat{A'}_{\widehat{\mathfrak m}'_\alpha}$.
Por outro lado, os
$\widehat{\mathfrak m}'_\alpha$
são os ideais maximais de
$\widehat{A'}$; portanto,
o homomorfismo canônico
$\widehat{A'}\to
\prod_\alpha
\widehat{A'}_{\widehat{\mathfrak m}'_\alpha}$
é injetivo (Bourbaki,
\emph{Alg. comm.}, cap.~II,
\S~3, n$^\circ$~3,
cor.~2 do teorema~1), e,
como por hipótese
$\widehat u:\widehat A\to
\widehat{A'}$ também é
injetivo, isso conclui a
demonstração no caso geral.
A última afirmação decorre
de que $\widehat A$ é um
$A$-módulo fielmente plano
$(\mathbf{0}_{\mathrm I},
7.3.5)$ e
$\widehat{A'}=
A'\otimes_A\widehat A$,
pois $A'$ é, por hipótese,
um $A$-módulo de tipo
finito (Bourbaki,
\emph{Alg. comm.}, cap.~III,
\S~3, n$^\circ$~4,
teorema~3, e cap.~IV,
\S~2, n$^\circ$~5,
cor.~3 da prop.~9).
\end{proof}

\begin{proposition}[11.5.3]
\label{IV.11.5.3}
Sejam $f:X\to Y$ um morfismo
localmente de apresentação
finita, $\sh{F}$ um
$\sh{O}_X$-módulo
quase-coerente de apresentação
finita, $x$ um ponto de $X$,
$y=f(x)$ e
$g=(\psi,\theta):Y'\to Y$
um morfismo próprio de
apresentação finita.
Suponhamos que:
\begin{enumerate}[label=(\roman*)]
\item O homomorfismo
$\theta_y:\sh{O}_y\to
(g_*\sh{O}_{Y'})_y$
seja injetivo.
\item Para todo
$x'\in X'=X\times_Y Y'$
cuja projeção em $X$ é
$x$, $\sh{F}'=\sh{F}
\otimes_Y Y'$ é
$Y'$-plano no ponto $x'$.
\end{enumerate}
Então $\sh{F}$ é
$f$-plano no ponto $x$.
\end{proposition}

\begin{proof}
Como a questão é local
em $X$, podemos supor
que $f$ seja de
apresentação finita.
Seja $p:X'\to X$ a
projeção canônica. Como
$f'=f_{(Y')}:X'\to Y'$
é de apresentação finita
(1.6.2), e $\sh{F}'$ é um
$\sh{O}_{X'}$-módulo de
apresentação finita
$(\mathbf{0}_{\mathrm I},
5.2.5)$, resulta de
\sref{IV.11.3.1} que o
conjunto $U'$ dos pontos
de $X'$ nos quais
$\sh{F}'$ é $f'$-plano
é aberto. Como $U'$ contém
$p^{-1}(x)$ por hipótese
e $p$ é próprio, portanto
fechado, $U'$ contém um
conjunto da forma
$p^{-1}(U)$, onde $U$ é
uma vizinhança de $x$.
Substituindo $X$ por $U$,
podemos, portanto, supor
desde já que $\sh{F}'$
é $f'$-plano.

Por outro lado, levando em
conta $(\mathbf I,3.6.5)$,
$(\mathbf{II},5.4.2)$ e
(2.1.4), podemos substituir
$Y$ por
$\operatorname{Spec}(\sh{O}_y)$,
isto é, supor que
$Y=\operatorname{Spec}(A)$,
onde $A$ é um anel local.
Nessas condições, provaremos
que $\sh{F}$ é $f$-plano.
Em virtude de (5.13.3),
$A$ é limite indutivo de
subanéis locais noetherianos
$A_\lambda$ tais que a
injeção canônica
$A_\lambda\to A$ seja um
homomorfismo local. Em
virtude de (8.9.1), podemos
supor que
$X=X_\lambda\otimes_{A_\lambda}A$,
$f=f_\lambda\otimes1_A$ e
$\sh{F}=\sh{F}_\lambda
\otimes_{A_\lambda}A$,
para um $\lambda$ conveniente,
com
$f_\lambda:X_\lambda\to
Y_\lambda=
\operatorname{Spec}(A_\lambda)$
um morfismo de tipo finito
e $\sh{F}_\lambda$ um
$\sh{O}_{X_\lambda}$-módulo
coerente. Do mesmo modo,
podemos supor que
$Y'=Y'_\lambda
\otimes_{A_\lambda}A$ e
$g=g_\lambda\otimes1_A$,
onde $g_\lambda:Y'_\lambda
\to Y_\lambda$ é um
morfismo de apresentação
finita; além disso,
(8.10.5, (xii)) permite
supor que $g_\lambda$ é
próprio. Como, por hipótese,
o homomorfismo
$A\to\Gamma(Y',\sh{O}_{Y'})$
é injetivo e $A_\lambda$
é um subanel de $A$,
também é injetivo o
homomorfismo
$A_\lambda\to
\Gamma(Y',\sh{O}_{Y'})$.
Por fim, em virtude de
\sref{IV.11.2.6}, podemos
escolher $\lambda$ de modo
que $\sh{F}'_\lambda=
\sh{F}_\lambda\otimes_{Y_\lambda}
Y'_\lambda$ seja
$f'_\lambda$-plano, pois
$\sh{F}'=\sh{F}'_\lambda
\otimes_{Y'_\lambda}Y'$.
Essas observações mostram
que podemos, doravante,
supor que o anel $A$ é
noetheriano, mantendo-se
as demais hipóteses de
\sref{IV.11.5.3}.

Ponha-se então
$B=\widehat A$ e
$Z=\operatorname{Spec}(B)$;
como $B$ é um $A$-módulo
fielmente plano
$(\mathbf{0}_{\mathrm I},
7.3.5)$, dizer que
$\sh{F}$ é $f$-plano
equivale a dizer que
$\sh{F}\otimes_Y Z$ é
$Z$-plano (2.1.4).
Do mesmo modo, se pusermos
$Z'=Y'\times_Y Z$, o
morfismo $Z'\to Y'$ é
fielmente plano (2.2.13);
portanto, dizer que
$\sh{F}'$ é $f'$-plano
equivale a dizer que
$\sh{F}'\otimes_{Y'}Z'$
é $Z'$-plano. Por fim,
$h=g_{(Z)}:Z'\to Z$ é
próprio (\textbf{II},
5.4.2) e de tipo finito
(1.5.2), $\widehat A$ é
noetheriano e, se $z$ é
seu ponto
\oldpage[IV-3]{153}
fechado, o homomorfismo
$\sh{O}_z\to
(h_*\sh{O}_{Z'})_z$
é injetivo, pois decorre
de (2.3.1) que
$h_*\sh{O}_{Z'}=
g_*\sh{O}_{Y'}\otimes_Y Z$,
e nossa afirmação resulta
da hipótese (i) e da
definição dos módulos planos
$(\mathbf{0}_{\mathrm I},
6.1.1)$.

Podemos, portanto, supor
agora que o anel local
noetheriano $A$ é
\emph{completo}; a
demonstração estará concluída
se provarmos que a
interseção dos núcleos dos
homomorfismos
$A=\sh{O}_y\to\sh{O}_{y'}$,
quando $y'$ percorre
$g^{-1}(y)$, se reduz a $0$.
Com efeito, os
$\sh{O}_{y'}$ são anéis
locais noetherianos; logo,
para cada $y'$, a interseção
dos $\mathfrak m_{y'}^n$
$(n\geq0)$ se reduz a $0$.
Se $\mathfrak a_{n,y'}$ é a
imagem inversa de
$\mathfrak m_{y'}^n$ em
$A$, as interseções finitas
dos $\mathfrak a_{n,y'}$
são vizinhanças de $0$
em $A$, e a interseção de
todos os $\mathfrak a_{n,y'}$
se reduz a $0$; portanto,
essas interseções finitas
formam um sistema fundamental
de vizinhanças de $0$ em
$A$ (Bourbaki,
\emph{Alg. comm.}, cap.~III,
\S~2, n$^\circ$~7, prop.~8,
em cuja demonstração se pode
substituir a sequência
decrescente por um conjunto
filtrante qualquer). Poderemos,
então, aplicar
\sref{IV.11.5.1}.
Seja agora $s\in A$ um
elemento pertencente ao
núcleo de cada homomorfismo
$A\to\sh{O}_{y'}$; sua imagem
$s'$ em
$\Gamma(Y',\sh{O}_{Y'})$
é, portanto, uma seção de
$\sh{O}_{Y'}$ sobre $Y'$
tal que $s'_{y'}=0$ para
todo $y'\in g^{-1}(y)$.
Existe, pois, uma vizinhança
$V'$ de $g^{-1}(y)$ em $Y'$
tal que $s'|_{V'}=0$.
Mas, como $g$ é fechado,
$V'$ contém um conjunto da
forma $g^{-1}(V)$, onde
$V$ é uma vizinhança aberta
de $y$ em $Y$. Ora, $Y$ é
um esquema local e $y$ é
seu ponto fechado; portanto,
$V=Y$ e, consequentemente,
$V'=Y'$ e $s'=0$. Como
$A\to\Gamma(Y',\sh{O}_{Y'})$
é injetivo por hipótese,
temos $s=0$.
\end{proof}

O caso particular seguinte
de \sref{IV.11.5.3} nos
será útil no cap.~V:

\begin{corollary}[11.5.4]
\label{IV.11.5.4}
Seja $f=(\psi,\theta):X\to Y$
um morfismo próprio de
apresentação finita, e
seja $p:X\times_Y X\to X$
a primeira projeção.
Suponhamos que
$\theta:\sh{O}_Y\to
f_*(\sh{O}_X)$ seja
injetivo. Então, para que
$f$ seja plano, é necessário
e suficiente que $p$ o seja.
\end{corollary}

\begin{proposition}[11.5.5]
\label{IV.11.5.5}
Sejam $A$ um anel,
$Y=\operatorname{Spec}(A)$,
$f:X\to Y$ um morfismo
localmente de apresentação
finita, $\sh{F}$ um
$\sh{O}_X$-módulo
quase-coerente de apresentação
finita e $x$ um ponto de
$X$. Seja $A\to A'$ um
homomorfismo injetivo tal
que $A'$ seja integral sobre
$A$. Ponhamos
$Y'=\operatorname{Spec}(A')$,
$X'=X\times_Y Y'$,
$f'=f\times1$ e
$\sh{F}'=\sh{F}\otimes_Y Y'$.
Então, se $\sh{F}'$ é
$f'$-plano em todo ponto
de $X'$ cuja projeção em
$X$ é $x$, $\sh{F}$ é
$f$-plano no ponto $x$.
\end{proposition}

\begin{proof}
Suponhamos primeiro que
$A'$ seja uma $A$-álgebra
\emph{finita} e de
\emph{apresentação finita};
então o morfismo $Y'\to Y$
é próprio (\textbf{II},
6.1.11) e de apresentação
finita, de modo que as
hipóteses de
\sref{IV.11.5.3} estão
satisfeitas, donde a conclusão.
No caso geral, a proposição
decorre desse caso particular,
do fato de que $A'$ é limite
indutivo de suas
sub-$A$-álgebras \emph{finitas}
$A'_\lambda$ e dos dois
lemas seguintes:
\end{proof}

\begin{lemma}[11.5.5.1]
\label{IV.11.5.5.1}
Toda $A$-álgebra finita
$A'$ é limite indutivo
de $A$-álgebras
$A'_\lambda$ que são
finitas e de apresentação
finita.
\end{lemma}

\begin{proof}
Raciocinamos como em
(1.9.3.1). Com efeito,
$A'=B/\mathfrak J$, onde
$B$ é uma $A$-álgebra
finita que é um
$A$-módulo livre, e
$\mathfrak J$ é um ideal
de $B$ (1.4.7.1). Ora,
$\mathfrak J$ é limite
indutivo dos ideais
$\mathfrak J_\lambda
\subset\mathfrak J$ de
$B$ que são de tipo finito
(e, \emph{a fortiori},
$A$-módulos de tipo finito);
portanto, pela exatidão do
funtor $\varinjlim$,
$A'$ é limite indutivo das
$A$-álgebras
$B/\mathfrak J_\lambda$.
Ora, $B/\mathfrak J_\lambda$
é, por definição, um
$A$-módulo de apresentação
finita e, portanto, também
(1.4.7) uma $A$-álgebra
de apresentação finita.
\end{proof}

\oldpage[IV-3]{154}

Para aplicar este lema à
situação de
\sref{IV.11.5.5}, observemos
ainda que, se o homomorfismo
$A\to A'$ é injetivo,
também o é, \emph{a fortiori},
$A\to A'_\lambda$ para
todo $\lambda$.

\begin{lemma}[11.5.5.2]
\label{IV.11.5.5.2}
Sejam $A$ um anel,
$A'$ uma $A$-álgebra e
$(A'_\lambda)$ um sistema
indutivo de $A$-álgebras
tal que
$A'=\varinjlim A'_\lambda$;
ponhamos
$Y=\operatorname{Spec}(A)$,
$Y'=\operatorname{Spec}(A')$,
$Y'_\lambda=
\operatorname{Spec}(A'_\lambda)$.
Sejam $f:X\to Y$ um morfismo
de apresentação finita e
$\sh{F}$ um
$\sh{O}_X$-módulo
quase-coerente de apresentação
finita; ponhamos
$X'=X\times_Y Y'$,
$f'=f_{(Y')}$,
$\sh{F}'=\sh{F}\otimes_Y Y'$,
$X'_\lambda=X\times_Y Y'_\lambda$,
$f'_\lambda=f_{(Y'_\lambda)}$,
$\sh{F}'_\lambda=
\sh{F}\otimes_Y Y'_\lambda$.
Seja $x$ um ponto de $X$
tal que $\sh{F}'$ seja
$f'$-plano em todos os
pontos $x'\in X'$ acima de
$x$. Então existe $\lambda$
tal que $\sh{F}'_\lambda$
é $f'_\lambda$-plano em
todos os pontos de
$X'_\lambda$ acima de $x$.
\end{lemma}

\begin{proof}
Seja $U'$ o conjunto dos
$x'\in X'$ nos quais
$\sh{F}'$ é $f'$-plano no
ponto $x'$; sabemos por
\sref{IV.11.3.1} que $U'$
é aberto em $X'$, pois
$f'$ é de apresentação
finita (1.6.2). Da mesma
forma, o conjunto
$U'_\lambda$ dos pontos de
$X'_\lambda$ nos quais
$\sh{F}'_\lambda$ é
$f'_\lambda$-plano é aberto
em $X'_\lambda$; sabemos,
além disso, por
\sref{IV.11.2.6}, que
$U'=\bigcup_\lambda
v_\lambda^{-1}(U'_\lambda)$,
onde $v_\lambda:X'\to
X'_\lambda$ é a projeção
canônica. Consideremos o
esquema
$T=\operatorname{Spec}(k(x))$;
ponhamos
$T'=T\times_Y Y'$ e
$T'_\lambda=T\times_Y
Y'_\lambda$, e sejam
$w_\lambda:T'\to T'_\lambda$,
$g':T'\to X'$ e
$g'_\lambda:T'_\lambda\to
X'_\lambda$ as projeções
canônicas. Ponhamos
$V'_\lambda=(g'_\lambda)^{-1}
(U'_\lambda)$ e
$V'=(g')^{-1}(U')=
\bigcup_\lambda w_\lambda^{-1}
(V'_\lambda)$. Por hipótese,
tendo em conta
$(\mathbf I,3.6.1)$, temos
$V'=T'$. Como $T'$ é
quase-compacto, existe
$\lambda$ tal que
$w_\lambda^{-1}(V'_\lambda)
=T'$. Deduzimos então de
(8.3.3), aplicado aos
subconjuntos fechados
quase-compactos
$T'_\lambda-V'_\lambda$
de $T'_\lambda$, que existe
$\mu\geq\lambda$ tal que
$T'_\mu=V'_\mu$; isso
significa que
$\sh{F}'_\mu$ é
$f'_\mu$-plano em todos os
pontos de $X'_\mu$ cuja
projeção em $X$ é $x$.
\end{proof}

\subsection{Descida da planicidade por morfismos quaisquer: caso de um pré-esquema de base unirramificado}
\label{subsection:IV.11.6}

\begin{theorem}[11.6.1]
\label{IV.11.6.1}
Sejam $A$ um anel local
íntegro geometricamente
unirramificado
$(\mathbf{0},23.2.1)$,
$Y=\operatorname{Spec}(A)$,
$f:X\to Y$ um morfismo
localmente de apresentação
finita e $\sh{F}$ um
$\sh{O}_X$-módulo
quase-coerente de apresentação
finita. Sejam $A'$ um anel
local e $u:A\to A'$ um
homomorfismo local injetivo;
ponhamos
$Y'=\operatorname{Spec}(A')$,
$X'=X\times_Y Y'$,
$f'=f_{(Y')}$ e
$\sh{F}'=\sh{F}\otimes_Y Y'$.
Sejam $x$ um ponto de $X$
cuja imagem $f(x)=y$ é
o ponto fechado de $Y$,
e $x'$ um ponto de $X'$
cujas projeções em $X$ e
$Y'$ são, respectivamente,
$x$ e o ponto fechado
$y'$ de $Y'$. Então, se
$\sh{F}'$ é $f'$-plano no
ponto $x'$, $\sh{F}$ é
$f$-plano no ponto $x$.
\end{theorem}

\begin{proof}
Como a questão é local em
$X$, podemos nos limitar
ao caso em que $f$ é de
apresentação finita.
Procederemos em várias
etapas.

\medskip
\noindent\textbf{I)}
\emph{Redução ao caso em que
$A$ e $A'$ são anéis locais
integralmente fechados.}

Como $u$ é injetivo e $A$
é íntegro, existe um ideal
primo $\mathfrak p'$ de
$A'$ tal que
$u^{-1}(\mathfrak p')=0$
$(\mathbf{0}_{\mathrm I},
1.5.8)$; o homomorfismo
composto
$A\xrightarrow{u}A'\to
A''=A'/\mathfrak p'$ é,
portanto, injetivo e local.
Se $Y''=\operatorname{Spec}(A'')$,
$X''=X'\otimes_{A'}A''=
X\times_Y Y''$ e
$\sh{F}''=\sh{F}'\otimes_{A'}
A''$, então $\sh{F}''$
é $Y''$-plano nos pontos
de $X''$ acima de $x'$
(2.1.4). Substituindo,
se necessário, $A'$ por
$A''$ e levando em conta
$(\mathbf I,3.4.7)$,
podemos supor, em primeiro
lugar, que $A'$ é íntegro.
Se $K'$ é o corpo de frações
de $A'$, existe então um anel
de valoração $B'$ em $K'$
que domina $A'$; sendo o
homomorfismo composto
$A\xrightarrow{u}A'\to B'$
injetivo e local, o mesmo
raciocínio que
\oldpage[IV-3]{155}
anteriormente permite substituir
$A'$ por $B'$; podemos,
portanto, supor que o anel
local $A'$ é integralmente
fechado, sendo $A$ um subanel
local de $A'$ dominado por
$A'$. Seja $A_1$ o fecho
integral de $A$; é claro que
$A\subset A_1\subset A'$ e,
por hipótese, $A_1$ é um
anel local. Se $\mathfrak m$,
$\mathfrak m_1$ e $\mathfrak m'$
são os ideais maximais de
$A$, $A_1$ e $A'$, temos
$\mathfrak m\subset\mathfrak m_1
\subset\mathfrak m'$; com efeito,
$\mathfrak m_1$ é o \emph{único}
ideal primo de $A_1$ acima
de $\mathfrak m$, pois $A_1$
é um anel local (Bourbaki,
\emph{Alg. comm.}, cap.~V,
\S~2, n.º~1, prop.~1); como
$\mathfrak m'\cap A=\mathfrak m$,
temos $\mathfrak m'\cap A_1
\cap A=\mathfrak m$, logo
$\mathfrak m'\cap A_1
=\mathfrak m_1$.

Ponhamos $Y_1=\operatorname{Spec}(A_1)$,
$X_1=X\times_Y Y_1$,
$f_1=f_{(Y_1)}$ e
$\sh{F}_1=\sh{F}\otimes_Y Y_1$;
seja $x_1$ a projeção de $x'$
em $X_1$. Denotemos por $y_1$
o único ponto fechado de $Y_1$,
de modo que $y_1=f_1(x_1)$.
Por hipótese, o morfismo
$\operatorname{Spec}(k(y_1))
\to\operatorname{Spec}(k(y))$
é radicial; pela transitividade
das fibras $(\mathbf I,3.6.4)$
e $(\mathbf I,3.5.7)$, concluímos
que o morfismo
$f_1^{-1}(y_1)\to f^{-1}(y)$
é radicial e, em particular,
que $x_1$ é o \emph{único}
ponto de $X_1$ cujas projeções
em $X$ e $Y_1$ são,
respectivamente, $x$ e $y_1$.
Além disso, vimos que $y_1$
é o único ponto de $Y_1$ cuja
projeção em $Y$ é $y$; portanto,
$x_1$ é o único ponto de $X_1$
cuja projeção em $X$ é $x$.
Se provarmos que $\sh{F}_1$
é $f_1$-plano no ponto $x_1$,
poderemos aplicar
\sref{IV.11.5.5}, donde resultará
a conclusão. Ficamos, assim,
reduzidos ao caso em que $A$
também é \emph{integralmente
fechado}.

\medskip
\noindent\textbf{II)} \emph{Redução
ao caso em que $A$ e $A'$
são anéis locais de
$\mathbf Z$-álgebras de tipo
finito integralmente fechados.}

Podemos considerar $A'$ como
limite indutivo filtrante de
suas sub-$\mathbf Z$-álgebras
(íntegras) de tipo finito
$B'_{\lambda}$; além disso,
como $A'$ é integralmente
fechado e o fecho integral de
uma $\mathbf Z$-álgebra de tipo
finito também é uma
$\mathbf Z$-álgebra de tipo
finito (7.8.3), vemos que
$A'$ é limite indutivo de suas
sub-$\mathbf Z$-álgebras de
tipo finito $B'_{\lambda}$
\emph{integralmente fechadas}.
Se $\mathfrak m'_{\lambda}
=\mathfrak m'\cap B'_{\lambda}$,
$A'$ também é limite indutivo
dos subanéis locais
$(B'_{\lambda})_{\mathfrak m'_{\lambda}}
=A'_{\lambda}$ dominados por
$A'$ (5.13.3). Para todo
$\lambda$, $B'_{\lambda}\cap A$
também é limite indutivo de
suas sub-$\mathbf Z$-álgebras
de tipo finito $B_{\alpha\lambda}$;
portanto,
$A=\varinjlim_{\alpha,\lambda}
B_{\alpha\lambda}$ e, como
acima, podemos substituir
$B_{\alpha\lambda}$ nessa
fórmula por seu fecho integral
(contido, por hipótese, em
$B'_{\lambda}\cap A$), e depois
pelo anel local
$A_{\alpha\lambda}
=(B_{\alpha\lambda})_{\mathfrak m_{\alpha\lambda}}$,
onde $\mathfrak m_{\alpha\lambda}
=\mathfrak m\cap B_{\alpha\lambda}
=\mathfrak m'_{\lambda}\cap
B_{\alpha\lambda}$; assim,
$A_{\alpha\lambda}$ é dominado
por $A'_{\lambda}$ e por $A$.

Ponhamos
$Y_{\alpha\lambda}
=\operatorname{Spec}(A_{\alpha\lambda})$.
Decorre de (8.9.1) que existe
um par $(\alpha,\lambda)$
suficientemente grande, um
morfismo de tipo finito
$f_{\alpha\lambda}:
X_{\alpha\lambda}\to
Y_{\alpha\lambda}$ e um
$\sh{O}_{X_{\alpha\lambda}}$-módulo
coerente $\sh{F}_{\alpha\lambda}$
tais que
\[
X=X_{\alpha\lambda}
\times_{Y_{\alpha\lambda}}Y,
\qquad
f=f_{\alpha\lambda}\times1_Y,
\qquad
\sh{F}=\sh{F}_{\alpha\lambda}
\otimes_{Y_{\alpha\lambda}}Y;
\]
se $x_{\alpha\lambda}$ é a
projeção de $x$ em
$X_{\alpha\lambda}$, bastará
mostrar que $\sh{F}_{\alpha\lambda}$
é $f_{\alpha\lambda}$-plano
no ponto $x_{\alpha\lambda}$.
Como $A'$ é limite indutivo
dos $A'_{\mu}$ para
$\mu\geq\lambda$, $X'$ é
limite projetivo dos
$X_{\alpha\lambda}
\times_{Y_{\alpha\lambda}}Y'_{\mu}
=X'_{\mu}$, e temos também
$\sh{F}'=\sh{F}'_{\alpha\mu}
\otimes_{Y'_{\mu}}Y'$, onde
$\sh{F}'_{\alpha\mu}
=\sh{F}_{\alpha\lambda}
\otimes_{Y_{\alpha\lambda}}Y'_{\mu}$.
Aplicando (11.2.6), vemos
que podemos tomar $\mu$
suficientemente grande para
que $\sh{F}'_{\alpha\mu}$
seja $Y'_{\mu}$-plano no ponto
$x'_{\mu}$, projeção de $x'$
em $X'_{\mu}$; além disso,
pela construção dos $A'_{\mu}$,
a projeção de $x'_{\mu}$ em
$Y'_{\mu}$ é o ponto fechado
de $Y'_{\mu}$.

\medskip
\noindent\textbf{III)} \emph{Redução
ao caso em que o corpo residual
$k'$ de $A'$ é uma extensão
finita e radicial do corpo
residual $k$ de $A$.}

Podemos, em primeiro lugar,
repetir o raciocínio da parte I)
da demonstração de
\sref{IV.11.4.11}, levando em
conta que $\mathbf Z$ é um
anel de
\oldpage[IV-3]{156}
Jacobson; ficamos, assim,
reduzidos ao caso em que $k'$
é uma extensão \emph{finita}
de $k$, o que suporemos a
seguir. Seja $k''$ a maior
extensão separável de $k$
contida em $k'$, e seja $k_1$
uma extensão finita de Galois
de $k$ que contém $k''$; assim,
$k''\otimes_k k_1$ é um produto
direto de corpos isomorfos a
$k_1$. Como $k'$ é uma extensão
radicial de $k''$, também
$k'\otimes_k k_1$ é um produto
direto de extensões radiciais
de $k_1$.

Existe um anel local $A_1$
que é uma $A$-álgebra e um
$A$-módulo \emph{livre de tipo
finito}, tal que
$A_1/\mathfrak m A_1$ é
$k$-isomorfo a $k_1$
$(\mathbf 0_{\mathrm{III}},10.3.1.2)$.
Mais precisamente, podemos
supor que $k_1=k[T]/(r)$,
em que $r$ é um polinômio
mônico irredutível e separável
de $k[T]$ de grau $n$; se
$R$ é um polinômio mônico de
$A[T]$ cuja imagem canônica
é $r$ (e que, portanto, também
tem grau $n$), podemos tomar
$A_1=A[T]/(R)$. Ora, se $K$
é o corpo de frações de $A$,
é claro que $R$ é um polinômio
irredutível e separável de
$K[T]$; deduzimos, em primeiro
lugar, que $A_1$ é um anel
íntegro cujo corpo de frações
$K_1=K[T]/(R)$ é uma extensão
separável de $K$. Além disso,
se $t$ é a imagem canônica de
$T$ em $A_1$, os $t^j$
$(0\leq j<n)$ formam uma base
do $A$-módulo $A_1$, e suas
imagens em $k_1$ formam uma
base sobre $k$; daí deduzimos
que
\[
d=\det\bigl(\operatorname{Tr}_{A_1/A}
(t^{i+j})\bigr)
\]
é um elemento de $A$ cuja
classe em $k$ é diferente de
zero e que, por conseguinte,
é invertível. O mesmo
raciocínio que em (6.12.4.1, I))
prova então que o morfismo
$\operatorname{Spec}(A_1)
\to\operatorname{Spec}(A)$
é plano e tem fibras regulares;
concluímos, portanto, de
(6.5.4, (ii)), que $A_1$ é
\emph{integralmente fechado}.

Consideremos agora o anel
$A'_1=A'\otimes_A A_1$; ele
é um $A'$-módulo livre de tipo
finito e, portanto, um anel
semilocal (Bourbaki,
\emph{Alg. comm.}, cap.~IV,
\S~2, n.º~5, cor.~3 da prop.~9).
Além disso, os ideais maximais
dessa $A'$-álgebra finita estão
todos acima do ideal maximal
$\mathfrak m'$ de $A'$ e,
\emph{a fortiori}, contêm
$\mathfrak m A'_1$. Mas
\[
A'_1/\mathfrak m A'_1
=(A'/\mathfrak m A')\otimes_k k_1,
\]
e, como $k_1$ é uma extensão
separável e finita de $k$,
o radical de
$A'_1/\mathfrak m A'_1$ é igual
a $(\mathfrak m'/\mathfrak m A')
\otimes_k k_1$ (Bourbaki,
\emph{Alg.}, cap.~VIII,
\S~7, n.º~2, cor.~2 da prop.~3).
Se $\mathfrak n_i$
$(1\leq i\leq r)$ são os ideais
maximais de $A'_1$, os corpos
$A'_1/\mathfrak n_i$ são,
portanto, os corpos que
compõem a álgebra
$k'\otimes_k k_1$; em outras
palavras, são extensões
finitas \emph{radiciais} de
$k_1$.

Além disso, como $A\to A'$
é um homomorfismo injetivo,
$A_1\to A'_1$ também o é,
pois $A_1$ é um $A$-módulo
plano. O homomorfismo canônico
\[
A'_1\longrightarrow
\prod_{i=1}^{r}(A'_1)_{\mathfrak n_i}
\]
também é injetivo (Bourbaki,
\emph{Alg. comm.}, cap.~II,
\S~3, n.º~3, teorema~1),
e o mesmo vale para o composto
de $A_1$ nesse produto. Mas
$A_1$ é íntegro e os núcleos
dos homomorfismos
$A_1\to(A'_1)_{\mathfrak n_i}$
são em número finito; como
sua interseção é nula, um
deles já é nulo. Em outras
palavras, existe um
$B_1=(A'_1)_{\mathfrak n_i}$
tal que o homomorfismo
$A_1\to B_1$ é \emph{injetivo}
e \emph{local}.

Ponhamos
$Y'_1=\operatorname{Spec}(B_1)$
e $X'_1=X'\times_{Y'}Y'_1$.
Então
$\sh{F}'_1=\sh{F}'\otimes_{Y'}Y'_1$
é $Y'_1$-plano em todos os
pontos de $X'_1$ acima de
$x'$; além disso, o ideal
maximal de $B_1$ é o único
acima de $\mathfrak m'$,
de modo que todos esses pontos
têm como projeção em $Y'_1$
o ponto fechado $y'_1$. Seja
$x'_1$ um desses pontos.
Por outro lado, ponhamos
$Y_1=\operatorname{Spec}(A_1)$,
$X_1=X\times_Y Y_1$ e
$\sh{F}_1=\sh{F}\otimes_Y Y_1$;
se $x_1$ é a projeção de $x'_1$
em $X_1$, sua projeção em
$X$ é $x$ e sua projeção em
$Y_1$ é o ponto fechado $y_1$.
Se provarmos que
$(\sh{F}_1)_{x_1}$ é um
$\sh{O}_{y_1}$-módulo plano,
decorrerá que $\sh{F}_x$ é um
$\sh{O}_y$-módulo plano; com
efeito, $\sh{O}_{y_1}$ é um
$\sh{O}_y$-módulo plano,
logo $(\sh{F}_1)_{x_1}$ é um
$\sh{O}_y$-módulo plano
$(\mathbf 0_{\mathrm I},6.2.1)$.
\oldpage[IV-3]{157}
Mas $(\sh{F}_1)_{x_1}
=\sh{F}_x\otimes_{\sh{O}_x}
\sh{O}_{x_1}$, e $\sh{O}_{x_1}$
é um $\sh{O}_x$-módulo
\emph{fielmente plano}; logo,
por (2.2.11, (iii)),
$\sh{F}_x$ é um
$\sh{O}_y$-módulo plano.
Como $X'_1=X_1\times_{Y_1}Y'_1$
e $\sh{F}'_1=\sh{F}_1
\otimes_{Y_1}Y'_1$, estamos
efetivamente reduzidos à
situação do enunciado de
\sref{IV.11.6.1}, substituindo
$A$ e $A'$ por $A_1$ e $B_1$,
respectivamente.

\medskip
\noindent\textbf{IV)} \emph{Fim
da demonstração.} --- Estamos,
por fim, reduzidos a provar
\sref{IV.11.6.1} sob as
seguintes hipóteses adicionais:
\begin{enumerate}[label=(\roman*)]
  \item $A$ e $A'$ são anéis
  locais de $\mathbf Z$-álgebras
  de tipo finito (portanto,
  anéis \emph{excelentes}
  (7.8.3));
  \item $A$ é integralmente
  fechado;
  \item o corpo residual $k'$
  de $A'$ é uma extensão finita
  e radicial do corpo residual
  $k$ de $A$.
\end{enumerate}
Sabe-se então, por (7.8.3)
e (2.3.8), que, sob as condições
(i) e (ii), se $\mathfrak m$ e
$\mathfrak m'$ são os ideais
maximais de $A$ e $A'$,
respectivamente, a topologia
$\mathfrak m$-ádica sobre $A$
é induzida pela topologia
$\mathfrak m'$-ádica de $A'$.
O homomorfismo dos completados
$\widehat u:\widehat A
\to\widehat{A'}$ é, portanto,
injetivo. Por outro lado,
como o morfismo
$\operatorname{Spec}(k')
\to\operatorname{Spec}(k)$
é radicial, o mesmo ocorre
com $f'^{-1}(y')\to f^{-1}(y)$
$(\mathbf I,3.5.7)$; logo,
existe um único ponto $x'$
de $X'$ cujas projeções em
$X$ e $Y'$ são, respectivamente,
$x$ e $y'$. Podemos, assim,
aplicar o resultado de
\sref{IV.11.5.2}.
\end{proof}

\begin{corollary}[11.6.2]
\label{IV.11.6.2}
Sejam $A$ um anel local
íntegro unirramificado,
$Y=\operatorname{Spec}(A)$,
$f:X\to Y$ um morfismo
localmente de apresentação
finita e $\sh{F}$ um
$\sh{O}_X$-módulo quase-coerente
de apresentação finita.
Sejam $A'$ um anel local e
$A\to A'$ um homomorfismo
local injetivo; ponhamos
$Y'=\operatorname{Spec}(A')$,
$X'=X\times_Y Y'$, $f'=f_{(Y')}$
e $\sh{F}'=\sh{F}\otimes_Y Y'$.
Seja $x$ um ponto de $X$ cuja
projeção $f(x)=y$ é o ponto
fechado de $Y$; suponhamos
que $\sh{F}'$ é $f'$-plano em
todos os pontos $x'$ de $X'$
cujas projeções em $X$ e $Y'$
são, respectivamente, $x$ e
o ponto fechado $y'$ de $Y'$.
Então $\sh{F}$ é $f$-plano
no ponto $x$.
\end{corollary}

\begin{proof}
Podemos, com efeito, retomar
a parte I) da demonstração de
\sref{IV.11.6.1}, que prova,
com as mesmas notações, que,
se $\sh{F}_1$ é $f_1$-plano
em \emph{todos} os pontos
$x_1$ de $X_1$ cujas projeções
em $X$ e $Y_1$ são,
respectivamente, $x$ e $y_1$,
então $\sh{F}$ é $f$-plano
no ponto $x$. Ficamos, assim,
reduzidos ao caso em que $A$
é integralmente fechado e,
portanto, geometricamente
unirramificado; a conclusão
decorre então de
\sref{IV.11.6.1}.
\end{proof}

\subsection{Contraexemplos}
\label{subsection:IV.11.7}

\begin{env}[11.7.1]
\label{IV.11.7.1}
Consideremos primeiro o caso
em que $A$ é um anel local
\emph{artiniano} e em que se
satisfazem as hipóteses de
\sref{IV.11.4.11}, \emph{exceto}
a condição~(ii) relativa ao
corpo residual $k$ de $A$.
Veremos que, nesse caso,
a conclusão de \sref{IV.11.4.11}
pode falhar. Seja $k$ um corpo
que admite uma extensão de
Galois $k'$ de grau $[k':k]>1$,
e seja $\Gamma$ seu grupo de
Galois. Seja $A$ uma
$k$-álgebra com base formada
pelos três elementos $1,a,b$
e tabela de multiplicação
$a^2=ab=ba=b^2=0$; assim,
$A$ é um anel local artiniano
cujo ideal maximal
$\mathfrak m=ka+kb$ tem
quadrado nulo. Seja
$A'=A\otimes_k k'$, uma
$k'$-álgebra com base
$1,a,b$ e um anel local
artiniano cujo ideal maximal
$\mathfrak m'=k'a+k'b
=\mathfrak m A'$ tem quadrado
nulo; identificamos $A$
canonicamente com um subanel
de $A'$. Seja $\mathfrak J$
o subespaço vetorial sobre
$k'$ de $\mathfrak m'$ gerado
por $a+\gamma b$, em que
$\gamma\in k'$ não pertence
a $k$; é claro que
$\mathfrak J$ é um ideal de
$A'$.

Ponhamos $B=A'/\mathfrak J$;
trata-se de um anel artiniano
que é um $A$-módulo
\emph{não plano}. Caso
contrário (Bourbaki,
\emph{Alg. comm.}, cap.~II,
\S~3, n.º~2, cor.~2 da prop.~5),
$B$ seria um $A$-módulo
\emph{livre}. Como $A'$ também
é um $A$-módulo livre e o
homomorfismo canônico
$A'/\mathfrak m A'
\to B/\mathfrak m B$ é
bijetivo, o homomorfismo
canônico $A'\to B=A'/\mathfrak J$
também seria
\oldpage[IV-3]{158}
bijetivo (\emph{loc. cit.},
n.º~2, cor.~da prop.~6),
o que é absurdo. Em outras
palavras, se pomos
$Y=\operatorname{Spec}(A)$
e $X=\operatorname{Spec}(B)$,
$\sh{O}_X$ não é $Y$-plano
no único ponto $x$ de $X$.

Mas seja $A_1=B$ e
$Y_1=\operatorname{Spec}(A_1)$,
e consideremos o homomorfismo
canônico $A\to A_1$, que é
local e injetivo, pois
$\mathfrak J\cap A=0$. Se
$B_1=B\otimes_A A_1
=B\otimes_A B$, veremos
que existe um ponto de
$X_1=X\times_Y Y_1
=\operatorname{Spec}(B_1)$
no qual $\sh{O}_{X_1}$ é
$Y_1$-plano. Para isso,
observemos que
\[
B\otimes_A B
=\frac{A'\otimes_A A'}
{\operatorname{Im}(\mathfrak J
\otimes_A A')
+\operatorname{Im}(A'\otimes_A
\mathfrak J)}.
\]
Ora, a estrutura de
$C'=A'\otimes_A A'$ obtém-se
facilmente. Consideremos a
$A$-álgebra produto
$C=\prod_{\sigma\in\Gamma}
A'_{\sigma}$, em que todos os
$A'_{\sigma}$ são iguais a
$A'$, e a aplicação canônica
$\varphi:A'\otimes_A A'\to C$
definida por
$\varphi(x\otimes y)
=(\sigma(x)y)_{\sigma\in\Gamma}$,
com $\Gamma$ agindo
canonicamente sobre $A'$.
Passando aos quocientes,
obtemos um homomorfismo
$C'/\mathfrak m C'
\to C/\mathfrak m C$ que é
precisamente o homomorfismo
canônico $k'\otimes_k k'
\to\prod_{\sigma}k'_{\sigma}$,
com $k'_{\sigma}=k'$ para
todo $\sigma\in\Gamma$.
Sabemos que este último é
bijetivo (Bourbaki,
\emph{Alg.}, cap.~VIII,
\S~8, prop.~4); portanto,
$\varphi$ também é bijetivo,
pois $C'$ e $C$ são
$A$-módulos livres (Bourbaki,
\emph{Alg. comm.}, cap.~II,
\S~3, n.º~2, cor.~da prop.~6).

Do que precede, resulta que
$B\otimes_A B$ é uma
$A$-álgebra semilocal,
produto direto das
$A$-álgebras locais
$A'/(\sigma(\mathfrak J)
+\mathfrak J)$.
O fator correspondente à
identidade de $\Gamma$ é
isomorfo a $A'/\mathfrak J$
e, portanto, é um
$A_1$-módulo plano. Contudo,
como $\gamma\notin k$, existe
pelo menos um
$\sigma\in\Gamma$ tal que
$\sigma(\mathfrak J)
\ne\mathfrak J$; então
$\sigma(\mathfrak J)
+\mathfrak J=\mathfrak m'$,
e $A'/\mathfrak m'$ não é um
$(A'/\mathfrak J)$-módulo
plano, pois é o quociente
de $A'/\mathfrak J$ por um
ideal não nulo.
\end{env}

\begin{env}[11.7.2]
\label{IV.11.7.2}
Mostraremos agora que o
resultado de \sref{IV.11.5.4}
deixa de ser válido quando
já não se supõe que $f$ seja
um morfismo \emph{próprio}
(e, \emph{a fortiori},
\sref{IV.11.5.3} deixa de
ser exato quando já não se
supõe que $g$ seja próprio).
Sejam $k$ um corpo,
$A_0=k[S,T]$ o anel de
polinômios e
$A=A_0/A_0ST$ o anel quociente.
Assim, $Y=\operatorname{Spec}(A)$
é a curva redutível formada
pelos dois ``eixos de
coordenadas'' no plano afim
$V_k^2=\operatorname{Spec}(A_0)$.
O anel $A$ admite dois ideais
primos minimais
$\mathfrak p_1=A_0S/A_0ST$ e
$\mathfrak p_2=A_0T/A_0ST$;
como $A$ é reduzido, ele se
mergulha canonicamente em
$B=B_1\oplus B_2$, em que
$B_1=A/\mathfrak p_1$ e
$B_2=A/\mathfrak p_2$.
Além disso, $B_1$ se identifica
com $k[T]$ e $B_2$ com $k[S]$;
portanto, são anéis íntegros
integralmente fechados, e
$Z=\operatorname{Spec}(B)$
é precisamente a normalização
do pré-esquema $Y$ em relação
a $R(Y)$ $(\mathbf{II},6.3.8)$,
soma disjunta dos dois esquemas
$Z_1=\operatorname{Spec}(B_1)$
e $Z_2=\operatorname{Spec}(B_2)$.

Denotemos por $y$ o ``ponto
duplo'' de $Y$, correspondente
ao ideal maximal
$\mathfrak p_1+\mathfrak p_2
=\mathfrak m$ de $A$, e por
$z_1$ e $z_2$ os pontos de
$Z$ que se projetam em $y$,
correspondentes, respectivamente,
aos ideais maximais
$\mathfrak m_1=(T)$ e
$\mathfrak m_2=(S)$ de $B_1$
e $B_2$. Designemos por $X$
o subpré-esquema de $Z$
induzido no complementar de
$z_2$ em $Z$; temos, portanto,
$X=\operatorname{Spec}(B_1
\oplus(B_2)_S)$. É imediato
que o homomorfismo
$A\to A'=B_1\oplus(B_2)_S$
é injetivo, mas o morfismo
correspondente $f:X\to Y$
não é fechado, pois
$f(Z_2-\{z_2\})$ não é fechado
em $Y$, embora $Z_2-\{z_2\}$
seja fechado em $X$;
\emph{a fortiori}, $f$ não é
próprio.

Veremos agora que $f$ não é
plano no ponto $z_1$. Basta
mostrar $(\mathbf 0_{\mathrm I},6.6.2)$
que $\sh{O}_{z_1}$ não é um
$\sh{O}_y$-módulo fielmente
plano; para isso, basta ver
que o homomorfismo canônico
$\sh{O}_y\to\sh{O}_{z_1}$
não é injetivo. Isso é
imediato, pois $\sh{O}_{z_1}$
é um anel íntegro, enquanto
$\sh{O}_y$ admite divisores
de zero. Entretanto, a
primeira projeção
$p:X\times_Y X\to X$ é um
isomorfismo: com efeito,
\[
B_1\otimes_A B_1
=(A/\mathfrak p_1)
\otimes_A(A/\mathfrak p_1)
=A/\mathfrak p_1,
\]
$(B_2)_S\otimes_A(B_2)_S
=(B_2\otimes_A B_2)_S
=(B_2)_S$ pela mesma razão,
e, por fim,
$B_1\otimes_A(B_2)_S=0$,
pois a imagem canônica de
$S$ em $B_1$ é nula.
\end{env}

\begin{env}[11.7.3]
\label{IV.11.7.3}
O exemplo precedente pode
ser generalizado. Consideremos,
sobre um corpo $k$, uma curva
algébrica reduzida $Y$ que
admite um único ``ponto duplo
ordinário'' $y$ (noção que
será definida mais adiante de
modo geral), e sua normalização
$Z$, de modo que o morfismo
$g:Z\to Y$ seja finito, que
a restrição de $g$ a
$Z-g^{-1}(y)$ seja um
isomorfismo sobre $Y-\{y\}$
e que $g^{-1}(y)$ se reduza
a dois pontos ``simples''
$z_1,z_2$. Além disso, o
pré-esquema $g^{-1}(y)$ é a
soma dos dois pré-esquemas
$\operatorname{Spec}(k(z_1))$
e $\operatorname{Spec}(k(z_2))$,
canonicamente isomorfos a
$\operatorname{Spec}(k(y))$.

Seja $X$ o subpré-esquema de
$Z$ induzido no aberto
$Z-\{z_2\}$; o morfismo
$f:X\to Y$, restrição de
$g$ a $X$, não é próprio,
pois, caso contrário,
$(\mathbf{II},5.4.3)$ implicaria
o mesmo para a inclusão
canônica $j:X\to Z$, que
não é fechada. O morfismo
$f$ é radicial: para todo
$y'\in Y$, a fibra
$f^{-1}(y')$ contém um único
ponto $x'$, e
$k(y')\to k(x')$ é um
isomorfismo. Daí concluímos,
primeiro, que o morfismo
diagonal $\Delta_f:X\to
X\times_Y X$ é bijetivo
(1.7.7.1); por outro
lado, como $f$ é não
ramificado (17.4.2, d$'$),
$\Delta_f$ é uma imersão
aberta (17.4.2, b$'$).
Consequentemente, $\Delta_f$
é um isomorfismo, e a primeira
projeção $p:X\times_Y X\to X$
é seu isomorfismo inverso.

Entretanto, $f$ não é plano
no ponto $z_1$; caso contrário,
$\sh{O}_{z_1}$ seria um
$\sh{O}_y$-módulo fielmente
plano $(\mathbf 0_{\mathrm I},6.6.2)$.
Como $\sh{O}_y$ contém dois
ideais primos minimais
distintos $\mathfrak p_1$
e $\mathfrak p_2$, correspondentes
aos dois ``ramos'' de $Y$ no
ponto $y$, existiriam em
$\sh{O}_{z_1}$ dois ideais
primos distintos cujas imagens
inversas por
$u:\sh{O}_y\to\sh{O}_{z_1}$
seriam $\mathfrak p_1$ e
$\mathfrak p_2$
$(\mathbf 0_{\mathrm I},6.5.1)$.
Isso é absurdo, pois
$\sh{O}_{z_1}$ tem apenas
dois ideais primos distintos,
$0$ e o ideal maximal
$\mathfrak m_1$, enquanto
$u^{-1}(\mathfrak m_1)$ é o
ideal maximal $\mathfrak m$
de $\sh{O}_y$.
\end{env}

\begin{env}[11.7.4]
\label{IV.11.7.4}
Observemos que, no exemplo
precedente, o homomorfismo
$u$ é \emph{injetivo} quando
$Y$ é \emph{irredutível}
(podemos tomar, por exemplo,
$Y=\operatorname{Spec}
\bigl(k[S,T]/(S(S^2+T^2)
-(S^2-T^2))\bigr)$, uma
``cúbica com ponto duplo'').
Com efeito, substituindo $Y$
por uma vizinhança afim de
$y$, podemos então supor
que $Y=\operatorname{Spec}(A)$,
onde $A$ é íntegro, e,
portanto,
$Z=\operatorname{Spec}(B)$,
onde $B$ é o fecho integral
de $A$. Como $B$, $\sh{O}_y$
e $\sh{O}_{z_1}$ se identificam
então com subanéis do corpo
\oldpage[IV-3]{159}
de frações de $A$, $u$ é
evidentemente injetivo.
Observemos, por outro lado,
que o homomorfismo
$\widehat u:\widehat{\sh{O}}_y
\to\widehat{\sh{O}}_{z_1}$
\emph{não é injetivo}, pois
$\widehat{\sh{O}}_{z_1}$ é um
anel local íntegro (sendo
$z_1$ um ponto simples),
ao passo que
$\widehat{\sh{O}}_y$ tem dois
ideais primos minimais
distintos (correspondentes
aos dois ``ramos'' de $Y$)
e, portanto, admite divisores
de zero. Isso fornece um
exemplo que mostra que, no
enunciado de \sref{IV.11.5.2},
não se pode substituir a
hipótese de injetividade de
$\widehat u$ pela de
injetividade do próprio $u$,
mesmo quando $A'$ é um anel
local. Com efeito, basta
tomar, com as notações
precedentes, $A=\sh{O}_y$,
$A'=\sh{O}_{z_1}$,
$Y=\operatorname{Spec}(A)$ e
$X=\operatorname{Spec}(A')$;
o raciocínio de \sref{IV.11.7.3}
ainda prova que a primeira
projeção $p:X\times_Y X\to X$
é um isomorfismo, embora $f$
não seja plano no ponto $z_1$.
\end{env}

\begin{env}[11.7.5]
\label{IV.11.7.5}
Os exemplos de \sref{IV.11.7.2}
e \sref{IV.11.7.3} explicam
a restrição aos anéis locais
\emph{unirramificados} em
\sref{IV.11.6.1} e
\sref{IV.11.6.2}. Veremos agora
que, em \sref{IV.11.6.1},
não se pode enfraquecer a
hipótese sobre $A$ supondo
apenas que $A$ seja
unirramificado. Consideremos,
com efeito, o anel local
íntegro completo
$A=\mathbf R[[U,V]]/(U^2+V^2)$,
que é unirramificado, mas
não geometricamente
unirramificado
\sref{IV.6.15.11}.
Sabe-se (\emph{loc. cit.})
que, se $u,v$ são as imagens
de $U,V$ em $A$, o fecho
integral de $A$ é
$\overline A=A[t]$, com
$t=v/u$ e $t^2=-1$;
portanto, $\overline A$ é
isomorfo a $\mathbf C[[U]]$.
Ponhamos
$Y=\operatorname{Spec}(A)$
e $X=\operatorname{Spec}(\overline A)$,
a normalização de $Y$
$(\mathbf{II},6.3.8)$, e sejam
$y$ e $x$ os pontos fechados
de $Y$ e $X$, respectivamente.
Mostraremos que, para uma
$A$-álgebra local adequada
$A'$, se pomos
$Y'=\operatorname{Spec}(A')$
e $X'=X\times_Y Y'$ e se
$y'$ designa o ponto fechado
de $Y'$, então $\sh{O}_{X'}$
é $Y'$-plano em um ponto de
$X'$ cujas projeções em $X$
e $Y'$ são, respectivamente,
$x$ e $y'$, mas não é
$Y'$-plano em \emph{todos}
os pontos com essas projeções.
Daí resultará evidentemente
(2.1.4) que $\sh{O}_X$ não
é $Y$-plano no ponto $x$
(o que, aliás, é trivial
\emph{a priori}, pois
$\overline A$ não é um
$A$-módulo livre).

Seja $B=A\otimes_{\mathbf R}
\mathbf C$, isomorfo a
$\mathbf C[[U,V]]/
((U+iV)(U-iV))$; $B$ tem
dois ideais primos minimais
$\mathfrak p'$ e $\mathfrak p''$,
gerados, respectivamente,
por $u+iv$ e $u-iv$, e
$\mathfrak n=\mathfrak p'
+\mathfrak p''$ é o ideal
maximal do anel local completo
$B$. Seja
$\overline B=\overline A
\otimes_{\mathbf R}\mathbf C$;
$\overline B$ é o produto
direto de duas álgebras
isomorfas a $\mathbf C[[U]]$,
geradas pelos idempotentes
$e'=(1\otimes1+t\otimes i)/2$
e $e''=(1\otimes1-t\otimes i)/2$.
Como o homomorfismo
$A\to\overline A$ é injetivo,
o mesmo vale para
$B\to\overline B$; as imagens
de $u+iv$ e $u-iv$ por essa
injeção são, respectivamente,
$2ue'$ e $2ue''$\footnote{O
impresso francês e os
testemunhos espanhol e inglês
registram $ue'$ e $ue''$.
Pelas definições
$e'=(1+it)/2$ e $e''=(1-it)/2$,
as imagens são
$u(1+it)=2ue'$ e
$u(1-it)=2ue''$. A leitura
impressa fica preservada
como testemunho, não como
igualdade nesta edição.};
concluímos imediatamente
que $\overline B$ se identifica
canonicamente com
$(B/\mathfrak p')\oplus
(B/\mathfrak p'')$.

Tomemos, então, para $A'$
o anel local $B/\mathfrak p'$;
assim, $\overline A\otimes_A A'$
se identifica com
$\overline B\otimes_B A'$.
Mas temos
$(B/\mathfrak p')\otimes_B
(B/\mathfrak p')=B/\mathfrak p'$
e $(B/\mathfrak p')\otimes_B
(B/\mathfrak p'')=B/\mathfrak n$;
logo $\overline A\otimes_A A'$
é isomorfo a
$A'\oplus(B/\mathfrak n)$.
Isso prova nossa afirmação,
pois $B/\mathfrak n
=A'/(\mathfrak n/\mathfrak p')$
não é um $A'$-módulo plano
(caso contrário, seria um
$A'$-módulo livre (Bourbaki,
\emph{Alg. comm.}, cap.~II,
\S~3, n.º~2, cor.~2 da
prop.~5), o que é absurdo).
\end{env}

\subsection{Um critério valuativo de planicidade}
\label{subsection:IV.11.8}

\begin{theorem}[11.8.1]
\label{IV.11.8.1}
Sejam $f:X\to Y$ um morfismo
localmente de apresentação
finita, $\sh{F}$ um
$\sh{O}_X$-módulo
quase-coerente de apresentação
finita, $x$ um ponto de $X$
e $y=f(x)$. Suponhamos que
o anel local $\sh{O}_y$ seja
íntegro (resp. reduzido e
noetheriano). Para que
$\sh{F}$ seja $f$-plano no
ponto $x$, é necessário e
suficiente que, para todo
anel de valuação (resp.
todo anel de valuação
discreta) $A'$ e todo
homomorfismo local
$\sh{O}_y\to A'$, seja
satisfeita a seguinte
condição: pondo
$Y'=\operatorname{Spec}(A')$,
$X'=X\times_Y Y'$,
$f'=f_{(Y')}$ e
$\sh{F}'=\sh{F}\otimes_Y Y'$,
o $\sh{O}_{X'}$-módulo
$\sh{F}'$ é $f'$-plano em
todos os pontos $x'$ de
$X'$ cujas projeções em
$X$ e $Y'$ são,
respectivamente, $x$ e o
ponto fechado $y'$ de $Y'$.
\end{theorem}

\begin{proof}
A necessidade da condição
é evidente \sref{IV.2.1.4};
resta provar sua suficiência.
Podemos nos limitar
(\sref[I]{I.3.6.5} e
\sref[I]{I.2.4.4}) ao caso
em que
$Y=\operatorname{Spec}(A)$
é o espectro de um anel
\emph{local} $A$ e $y$ é o
ponto fechado de $Y$.

\medskip
\noindent\textbf{(i)}
\emph{Caso em que $A$ é
íntegro.} --- Sejam $K$ o
corpo de frações de $A$ e
$A_1$ o fecho integral de
$A$; se pomos
$Y_1=\operatorname{Spec}(A_1)$,
$X_1=X\times_Y Y_1$ e
$f_1=f_{(Y_1)}$, basta,
em virtude de
\sref{IV.11.5.5}, mostrar
que
$\sh{F}_1=\sh{F}\otimes_Y Y_1$
é $f_1$-plano em \emph{todos}
os pontos $x_1$ de $X_1$
cuja projeção em $X$ é $x$.
Ora, se $f_1(x_1)=y_1$,
temos
$\mathfrak i_{y_1}=\mathfrak m_1$,
em que $\mathfrak m_1$ é um
ideal primo de $A_1$ cuja
contração a $A$ é
$\mathfrak m=\mathfrak i_y$
($\mathfrak m_1$ é, aliás,
necessariamente maximal).
Seja então $A'$ um anel de
valuação de $K$ que domina
$\sh{O}_{y_1}=(A_1)_{\mathfrak m_1}$;
como o homomorfismo
$A\to\sh{O}_{y_1}$ é local,
o mesmo vale para
$A\to A'$. Existe então
ao menos um
\oldpage[IV-3]{160}
ponto $x'\in X'$ cujas
projeções em $X_1$ e $Y'$
são, respectivamente,
$x_1$ e $y'$
\sref[I]{I.3.4.7}; como,
por hipótese, $\sh{F}'$
é $f'$-plano no ponto $x'$
e $\sh{O}_{y_1}$ é
integralmente fechado,
podemos aplicar
\sref{IV.11.6.1} e deduzir
que $\sh{F}_1$ é $f_1$-plano
no ponto $x_1$, donde
resulta o teorema neste caso.

\medskip
\noindent\textbf{(ii)}
\emph{Caso em que $A$ é
reduzido e noetheriano.}
--- Sejam $\mathfrak p_i$
$(1\leq i\leq m)$ os ideais
primos minimais de $A$;
como $A$ é reduzido,
identifica-se canonicamente
com um subanel do produto
dos
$A_i=A/\mathfrak p_i$,
que são anéis locais
noetherianos. Pondo
$Y_i=\operatorname{Spec}(A_i)$,
$X_i=X\times_Y Y_i$ e
$f_i=f_{(Y_i)}$,
segue de \sref{IV.11.5.2}
que basta mostrar, para
cada $i$, que
$\sh{F}_i=\sh{F}\otimes_Y Y_i$
é $f_i$-plano em todo
ponto $x_i$ de $X_i$ cujas
projeções em $X$ e $Y_i$
são, respectivamente,
$x$ e o ponto fechado
$y_i$ de $Y_i$.

Ora, como $A_i$ é um anel
local noetheriano íntegro,
existe em seu corpo de
frações $K_i$ um subanel
$A''_i$ que contém $A_i$,
é uma $A_i$-álgebra
\emph{finita} (portanto, um
anel semilocal noetheriano)
e cujos anéis locais são
\emph{geometricamente
unirramificados}
(\sref{IV.6.15.5} e
$(\mathbf 0,23.2.5)$).
Como os ideais maximais de
$A''_i$ estão necessariamente
sobre o ideal maximal de
$A_i$, deduzimos ainda de
\sref[I]{I.3.4.7} e
\sref{IV.11.5.2} que basta,
pondo
$Y''_i=\operatorname{Spec}(A''_i)$,
$X''_i=X\times_Y Y''_i$ e
$f''_i=f_{(Y''_i)}$,
provar que
$\sh{F}''_i=\sh{F}\otimes_Y Y''_i$
é $f''_i$-plano em todo
ponto $x''_i$ de $X''_i$
cujas projeções em $X$ e
$Y''_i$ são,
respectivamente, $x$ e o
ponto fechado $y''_i$ de
$Y''_i$.

Seja agora $A'$ um anel
de valuação discreta de
$K_i$ que domina $A''_i$,
e seja $x'$ um ponto de
$X'$ cujas projeções em
$X''_i$ e $Y'$ são,
respectivamente, $x''_i$
e $y'$ \sref[I]{I.3.4.7};
como $\sh{O}_{y''_i}$ é
geometricamente
unirramificado, podemos
aplicar novamente
\sref{IV.11.6.1} e deduzir
que $\sh{F}''_i$ é de fato
$f''_i$-plano no ponto
$x''_i$.
\end{proof}

\begin{remarks}[11.8.2]
\label{IV.11.8.2}
\begin{enumerate}[label=(\roman*)]
  \item No enunciado de
  \sref{IV.11.8.1}, podemos
  nos limitar a supor que
  a condição sobre $\sh{F}'$
  se verifica para os anéis
  de valuação $A'$
  \emph{completos} cujo corpo
  residual é
  \emph{algebricamente
  fechado}. De fato, sabe-se
  que todo anel de valuação
  $A'$ é dominado por um
  anel $A''$ dessa classe
  \sref[II]{II.7.1.2} e
  que, se $A'$ é um anel
  de valuação discreta,
  também podemos tomar
  $A''$ discreto
  $(\mathbf 0_{\mathrm{III}},
  10.3.1)$.

  \item A demonstração de
  \sref{IV.11.8.1} se
  simplifica quando supomos
  não apenas que $A$ é
  íntegro e noetheriano,
  mas também que seu
  completado $\widehat A$
  é \emph{íntegro}.
  Substituindo $X$ por
  $X\otimes_A\widehat A$
  e raciocinando como na
  demonstração de
  \sref{IV.11.5.3},
  podemos, de fato,
  nos reduzir a provar
  \sref{IV.11.8.1} quando
  $A=\sh{O}_y$ é íntegro,
  noetheriano e
  \emph{completo}. Ora,
  sabe-se \sref[II]{II.7.1.7}
  que tal anel $A$ é
  dominado por um anel
  de valuação discreta
  completo; a conclusão
  decorre, portanto,
  diretamente de
  \sref{IV.11.5.2}.
\end{enumerate}
\end{remarks}

\subsection{Famílias separantes e universalmente separantes de homomorfismos de feixes de módulos}
\label{subsection:IV.11.9}

\begin{env}[11.9.1]
\label{IV.11.9.1}
Sejam $X$ um pré-esquema,
$(f_\lambda)_{\lambda\in L}$
uma família de morfismos
$f_\lambda:Z_\lambda\to X$
e $\sh{F}$ um
$\sh{O}_X$-módulo
quase-coerente. Para cada
$\lambda\in L$, suponhamos
dados um
$\sh{O}_{Z_\lambda}$-módulo
quase-coerente
$\sh{G}_\lambda$ e um
homomorfismo
\begin{equation}
u_\lambda:\sh{F}
\longrightarrow
(f_\lambda)_*(\sh{G}_\lambda).
\label{IV.11.9.1.1}
\tag{11.9.1.1}
\end{equation}
\end{env}

\oldpage[IV-3]{161}

Dizemos que a família
$(u_\lambda)$ (ou a família
correspondente dos
$u_\lambda^\#:
(f_\lambda)^*(\sh{F})
\to\sh{G}_\lambda$) é
\emph{separante} se a
interseção dos núcleos
dos $u_\lambda$ é
\emph{nula}. Em outros
termos, isso significa que,
para \emph{todo} aberto
$U$ de $X$ e toda seção
$t$ de $\sh{F}$ sobre
$U$, se para \emph{todo}
$\lambda$ a seção
$u_\lambda(t)$ (que, por
definição, é uma seção de
$\sh{G}_\lambda$ sobre
$f_\lambda^{-1}(U)$) é
nula, então a própria
$t$ é nula.

\begin{env}[11.9.2]
\label{IV.11.9.2}
Com as notações de
\sref{IV.11.9.1}, seja
$M$ um segundo conjunto
de índices; para cada
$\lambda\in L$, seja
$(g_{\lambda\mu})_{\mu\in M}$
uma família de morfismos
$g_{\lambda\mu}:
Z'_{\lambda\mu}\to Z_\lambda$.
Para cada par
$(\lambda,\mu)$,
suponhamos dados um
$\sh{O}_{Z'_{\lambda\mu}}$-módulo
quase-coerente
$\sh{H}_{\lambda\mu}$ e um
homomorfismo
$v_{\lambda\mu}:
\sh{G}_\lambda\to
(g_{\lambda\mu})_*
(\sh{H}_{\lambda\mu})$.
Ponhamos
$h_{\lambda\mu}
=f_\lambda\circ g_{\lambda\mu}:
Z'_{\lambda\mu}\to X$
e consideremos o
homomorfismo composto
\[
w_{\lambda\mu}:
\sh{F}
\xrightarrow{u_\lambda}
(f_\lambda)_*
(\sh{G}_\lambda)
\xrightarrow{(f_\lambda)_*
(v_{\lambda\mu})}
(h_{\lambda\mu})_*
(\sh{H}_{\lambda\mu}).
\]
Suponhamos que, para
cada $\lambda\in L$,
a família
$(v_{\lambda\mu})_{\mu\in M}$
seja separante; então
o mesmo vale para a
família
$((f_\lambda)_*
(v_{\lambda\mu}))_{\mu\in M}$,
como se vê imediatamente.
Concluímos que, para que
a família $(u_\lambda)$
seja separante, é
necessário e suficiente
que o seja a família
$(w_{\lambda\mu})_
{(\lambda,\mu)\in L\times M}$.
\end{env}

\begin{env}[11.9.3]
\label{IV.11.9.3}
Para verificar que a família
$(u_\lambda)$ é separante
(com as notações de
\sref{IV.11.9.1}), podemos
primeiro nos reduzir ao
caso em que $X$ é
\emph{afim}, pois a
propriedade é local em
$X$. Podemos, além disso,
supor que $Z_\lambda=X$
\emph{para todo}
$\lambda\in L$. De fato,
seja $M$ um conjunto de
índices, soma de uma
família
$(M_\lambda)_{\lambda\in L}$,
e, para cada
$\lambda\in L$, seja
$(Y_{\lambda\mu})_
{\mu\in M_\lambda}$
uma cobertura aberta afim
de $Z_\lambda$; seja
$j_{\lambda\mu}:
Y_{\lambda\mu}\to Z_\lambda$
a inclusão canônica e
ponhamos
$\sh{H}_{\lambda\mu}
=j_{\lambda\mu}^*
(\sh{G}_\lambda)
=\sh{G}_\lambda|
Y_{\lambda\mu}$.
Considerando o
homomorfismo canônico
$v_{\lambda\mu}
=\rho_{\sh{G}_\lambda}:
\sh{G}_\lambda\to
(j_{\lambda\mu})_*
(j_{\lambda\mu}^*
(\sh{G}_\lambda))
=(j_{\lambda\mu})_*
(\sh{H}_{\lambda\mu})$
relativo a
$j_{\lambda\mu}$
$(\mathbf 0_{\mathrm I},
4.4.3.2)$, vemos
imediatamente que, para
cada $\lambda\in L$,
a família
$(v_{\lambda\mu})_
{\mu\in M_\lambda}$
é separante. Em virtude
de \sref{IV.11.9.2},
ficamos, portanto,
reduzidos a provar que
a família dos
homomorfismos compostos
$((f_\lambda)_*
(v_{\lambda\mu}))
\circ u_\lambda$ é
separante; em outras
palavras, reduzimo-nos
ao caso em que os
$Z_\lambda$ são afins.
Mas, então, os
$(f_\lambda)_*
(\sh{G}_\lambda)$ são
$\sh{O}_X$-módulos
quase-coerentes
\sref[I]{I.1.6.2} e,
pela definição, podemos
substituir os $Z_\lambda$
por $X$ e os
$\sh{G}_\lambda$ por
$(f_\lambda)_*
(\sh{G}_\lambda)$,
o que prova nossa
afirmação.

Observemos ainda que,
se $L$ é \emph{finito}
e os $f_\lambda$ são
\emph{quase-compactos},
podemos, na redução
precedente, supor que
os $M_\lambda$ também
são \emph{finitos};
nesse caso, reduzimo-nos
a verificar que uma
família \emph{finita}
de homomorfismos de
$\sh{F}$ em
$\sh{O}_X$-módulos
quase-coerentes é
separante.
\end{env}

\begin{env}[11.9.4]
\label{IV.11.9.4}
Consideremos, portanto, o
caso em que
$Z_\lambda=X$ para todo
$\lambda$ e em que
$X=\operatorname{Spec}(A)$
é afim. Temos, então,
$\sh{F}=\widetilde M$ e
$\sh{G}_\lambda=
\widetilde N_\lambda$,
em que $M$ e $N_\lambda$
são $A$-módulos, e
$u_\lambda=
\widetilde\varphi_\lambda$,
em que
$\varphi_\lambda:
M\to N_\lambda$ são
$A$-homomorfismos.
Dizer que a família
$(u_\lambda)$ é separante
significa, então, que,
para todo $s\in A$, a
interseção dos núcleos
dos
$(\varphi_\lambda)_s:
M_s\to(N_\lambda)_s$
é nula. Dizemos também,
nesse caso, que a família
$(\varphi_\lambda)$ é
\emph{separante}.
Observemos que, se $L$
é \emph{finito}, isso
equivale a dizer que a
interseção dos núcleos
dos $\varphi_\lambda$
é nula, pois então
\[
\bigcap_{\lambda\in L}
\operatorname{Ker}
((\varphi_\lambda)_s)
=\left(
\bigcap_{\lambda\in L}
\operatorname{Ker}
(\varphi_\lambda)
\right)_s
\quad
(\mathbf 0_{\mathrm I},1.3.2).
\]
Mas essa relação deixa
de valer em geral quando
$L$ é \emph{infinito},
e o fato de a interseção
dos núcleos dos
$\varphi_\lambda$ ser
nula \emph{não implica,
portanto},
\oldpage[IV-3]{162}
em geral, que a família
$(\varphi_\lambda)$ seja
separante. Por exemplo,
suponhamos que $A$ seja
um anel de valoração
discreta, de ideal
maximal $\mathfrak m$,
e consideremos a família
de homomorfismos
$\varphi_k:
A\to A/\mathfrak m^k$,
cujos núcleos
$\mathfrak m^k$ têm
interseção nula. Essa
família, entretanto,
não é separante, pois
as fibras de todos os
$\mathfrak m^k$ no
ponto genérico $x$ de
$X=\operatorname{Spec}(A)$
(que é aberto em $X$)
são iguais a
$\sh{O}_x=k(x)$, o corpo
de frações de $A$, e sua
interseção, portanto,
não é nula.
\end{env}

\begin{env}[11.9.5]
\label{IV.11.9.5}
No que se segue, vamos nos
ocupar principalmente do
problema da \emph{mudança
de base} para as famílias
separantes. Com as
notações de
\sref{IV.11.9.1},
consideremos um morfismo
$g:X'\to X$ e ponhamos
$\sh{F}'=
\sh{F}\otimes_{\sh{O}_X}
\sh{O}_{X'}=g^*(\sh{F})$
e, para cada $\lambda$,
$Z'_\lambda=
Z_\lambda\times_X X'$,
$f'_\lambda=f_\lambda
\times 1$ e
$\sh{G}'_\lambda=
\sh{G}_\lambda
\otimes_{\sh{O}_{Z_\lambda}}
\sh{O}_{Z'_\lambda}
=(g'_\lambda)^*
(\sh{G}_\lambda)$,
em que
$g'_\lambda:
Z'_\lambda\to Z_\lambda$
é a projeção canônica.
Para cada $\lambda$,
designemos por
$u'_\lambda:
\sh{F}'\to
(f'_\lambda)_*
(\sh{G}'_\lambda)$
o homomorfismo obtido
da seguinte maneira:
seja
\[
\psi:
(f_\lambda)_*
(\sh{G}_\lambda)
\longrightarrow
(f_\lambda)_*
((g'_\lambda)_*
((g'_\lambda)^*
(\sh{G}_\lambda)))
=g_*((f'_\lambda)_*
(\sh{G}'_\lambda))
\]
o homomorfismo
$(f_\lambda)_*
(\rho_{\sh{G}_\lambda})$,
em que
$\rho_{\sh{G}_\lambda}$
é o homomorfismo
canônico
$(\mathbf 0_{\mathrm I},
4.4.3.2)$ correspondente
a $g'_\lambda$.
Então $u'_\lambda$
é definido como o
composto
\[
g^*(\sh{F})
\xrightarrow{g^*(u_\lambda)}
g^*((f_\lambda)_*
(\sh{G}_\lambda))
\xrightarrow{g^*(\psi)}
g^*g_*((f'_\lambda)_*
(\sh{G}'_\lambda))
\xrightarrow{\sigma}
(f'_\lambda)_*
(\sh{G}'_\lambda),
\]
em que
$\sigma=
\sigma_{(f'_\lambda)_*
(\sh{G}'_\lambda)}$
é o homomorfismo
canônico
$(\mathbf 0_{\mathrm I},
4.4.3.3)$ correspondente
a $g$. Diremos, por
abreviação, que
$u'_\lambda$ é
\emph{deduzido de
$u_\lambda$ pela mudança
de base} $g$. Quando
$f_\lambda:
Z_\lambda\to X$ é
um morfismo \emph{afim},
temos
$(f'_\lambda)_*
(\sh{G}'_\lambda)
=g^*((f_\lambda)_*
(\sh{G}_\lambda))$;
nesse caso, portanto,
temos simplesmente
$u'_\lambda=g^*(u_\lambda)$.

Também podemos interpretar
$u'_\lambda$ da seguinte
forma: basta conhecer o
valor de $u'_\lambda(t')$
quando $t'$ é uma seção
de $\sh{F}'$ sobre um
aberto $U'$ de $X'$ do
seguinte tipo particular:
$t'$ é a restrição a
$U'$ da imagem canônica,
pelo homomorfismo
$\Gamma(U,\sh{F})\to
\Gamma(g^{-1}(U),
\sh{F}')$, de uma seção
$t$ de $\sh{F}$ sobre
um aberto $U$ de $X$
que contém $g(U')$
(essas seções geram,
de fato, o
$\sh{O}_{X'}$-módulo
$\sh{F}'$
$(\mathbf 0_{\mathrm I},
3.7.1)$). Consideremos
a seção $u_\lambda(t)$
de $\sh{G}_\lambda$
sobre $f_\lambda^{-1}(U)$
e sua imagem canônica
$t''$ por
\[
\Gamma(f_\lambda^{-1}(U),
\sh{G}_\lambda)
\longrightarrow
\Gamma((g'_\lambda)^{-1}
(f_\lambda^{-1}(U)),
\sh{G}'_\lambda);
\]
então $u'_\lambda(t')$
é a restrição de $t''$
a $(f'_\lambda)^{-1}(U')$.

Consideremos, em
particular, o caso em
que $Z_\lambda$ é o
subpré-esquema induzido
sobre um aberto de $X$,
$\sh{G}_\lambda=
j_\lambda^*(\sh{F})$,
em que $j_\lambda:
Z_\lambda\to X$ é a
inclusão canônica, e
$u_\lambda$ é o
homomorfismo canônico
$\rho_{\sh{F}}:
\sh{F}\to
(j_\lambda)_*
(j_\lambda^*(\sh{F}))
=(j_\lambda)_*
(\sh{G}_\lambda)$
$(\mathbf 0_{\mathrm I},
4.4.3.2)$ correspondente
a $j_\lambda$. Então
$Z'_\lambda$ é induzido
sobre um aberto de $X'$,
e a interpretação anterior
mostra que $u'_\lambda$
não é senão o
homomorfismo canônico
$\rho_{\sh{F}'}:
\sh{F}'\to
(j'_\lambda)_*
((j'_\lambda)^*
(\sh{F}'))
=(j'_\lambda)_*
(\sh{G}'_\lambda)$
correspondente à
inclusão canônica
$j'_\lambda:
Z'_\lambda\to X'$.
\end{env}

\begin{env}[11.9.6]
\label{IV.11.9.6}
Nas condições de
\sref{IV.11.9.5},
suponhamos que $X$ e
$X'$ sejam afins e que
queiramos provar que
toda seção $t'$ de
$\sh{F}'$ sobre $X'$
cujas imagens por todos
os $u'_\lambda$ são
nulas é ela própria
nula. Podemos então
nos limitar também ao
caso em que
$Z_\lambda=X$ para
todo $\lambda\in L$.
De fato, com as notações
de \sref{IV.11.9.3}
e \sref{IV.11.9.5},
se pomos
$Y'_{\lambda\mu}
=(g'_\lambda)^{-1}
(Y_{\lambda\mu})$,
o homomorfismo
$v'_{\lambda\mu}$
deduzido de
$v_{\lambda\mu}$ pela
mudança de base $g$
não é senão o
homomorfismo canônico
\oldpage[IV-3]{163}
$\rho_{\sh{G}'_\lambda}:
\sh{G}'_\lambda\to
(j'_{\lambda\mu})_*
((j'_{\lambda\mu})^*
(\sh{G}'_\lambda))$
correspondente à
inclusão canônica
$j'_{\lambda\mu}:
Y'_{\lambda\mu}\to
Z'_\lambda$, como vimos
em \sref{IV.11.9.5}.
A afirmação decorre,
então, dos raciocínios
de \sref{IV.11.9.2}
e \sref{IV.11.9.3},
substituindo
$Y_{\lambda\mu}$ e
$v_{\lambda\mu}$ por
$Y'_{\lambda\mu}$ e
$v'_{\lambda\mu}$.

Temos uma redução
semelhante quando
queremos provar que a
família $(u'_\lambda)$
é \emph{separante}
($X$ e $X'$ sendo
afins): isso decorre
novamente de
\sref{IV.11.9.2}
e \sref{IV.11.9.3}.
\end{env}

\begin{proposition}[11.9.7]
\label{IV.11.9.7}
Com as notações de
\sref{IV.11.9.1} e
\sref{IV.11.9.5},
suponhamos que
$X=\operatorname{Spec}(A)$
e
$X'=\operatorname{Spec}(A')$
sejam afins e, além
disso, que $A'$ seja
um $A$-módulo projetivo.
Então, se a família
$(u_\lambda)$ é
separante, toda seção
$t'$ de $\sh{F}'$ sobre
$X'$ cujas imagens por
todos os $u'_\lambda$
são nulas é ela própria
nula.
\end{proposition}

\begin{proof}
Vimos em \sref{IV.11.9.6}
que podemos nos limitar
ao caso em que todos
os $Z_\lambda$ são
iguais a $X$. A proposição
é então consequência
do lema seguinte:
\end{proof}

\begin{lemma}[11.9.7.1]
\label{IV.11.9.7.1}
Sejam $A$ um anel,
$(M_\lambda)_{\lambda\in L}$
uma família de
$A$-módulos, $M$ um
$A$-módulo e, para cada
$\lambda$, seja
$u_\lambda:
M\to M_\lambda$ um
homomorfismo.
Suponhamos que a
interseção dos núcleos
dos $u_\lambda$ seja
nula. Então, para todo
$A$-módulo projetivo
$N$, a interseção dos
núcleos dos
homomorfismos
$u_\lambda\otimes 1:
M\otimes_A N\to
M_\lambda\otimes_A N$
é nula.
\end{lemma}

\begin{proof}
Com efeito, $N$ é um
somando direto de um
$A$-módulo livre $L$,
e basta evidentemente
provar que a interseção
dos núcleos dos
homomorfismos
$v_\lambda:
M\otimes_A L\to
M_\lambda\otimes_A L$
é nula, pois
$u_\lambda\otimes 1:
M\otimes_A N\to
M_\lambda\otimes_A N$
é a restrição de
$v_\lambda$. Mas essa
afirmação decorre então
trivialmente da hipótese.
\end{proof}

\begin{remark}[11.9.8]
\label{IV.11.9.8}
Não sabemos se, sob as
hipóteses da proposição
\sref{IV.11.9.7}, a
família $(u'_\lambda)$
é \emph{separante}:
seria necessário, com
efeito, em virtude de
\sref{IV.11.9.4},
provar que uma seção
$t'$ de $\sh{F}'$
\emph{sobre um aberto}
$D(h')\subset X'$
(em que $h'\in A'$)
cujas imagens por todos
os $u'_\lambda$ são
nulas é ela própria
nula. Ora, não podemos
aplicar a proposição
\sref{IV.11.9.7}\footnote{O
impresso francês e os
testemunhos espanhol e
inglês trazem aqui a
expressão proposição
(11.9.4). O n.º~11.9.4
não é uma proposição
e não supõe
projetividade; a
proposição cuja aplicação
é impedida é 11.9.7.}
a
$D(h')=
\operatorname{Spec}(A'_{h'})$,
pois de $A'$ ser um
$A$-módulo projetivo
(mesmo livre) não
decorre que
$A'_{h'}$ seja um
$A$-módulo projetivo.
Por exemplo, podemos
tomar para $A$ um anel
de valoração discreta,
para $A'$ um anel de
valoração discreta
que seja um $A$-módulo
livre de posto $2$,
e para $A'_{h'}$ o
corpo de frações de
$A'$.
\end{remark}

Temos, contudo, o
seguinte resultado:

\begin{corollary}[11.9.9]
\label{IV.11.9.9}
Sejam $X$ um pré-esquema
artiniano e $g:X'\to X$
um morfismo plano
(observemos que essas
duas condições são
satisfeitas quando $X$
é o espectro de um
\emph{corpo} e $g$ é
um morfismo
\emph{qualquer}).
Então, com as notações
de \sref{IV.11.9.1}
e \sref{IV.11.9.5},
se a família
$(u_\lambda)$ é
separante, o mesmo
vale para a família
$(u'_\lambda)$.
\end{corollary}

\begin{proof}
Podemos nos limitar ao
caso em que
$X=\operatorname{Spec}(A)$
é o espectro de um
anel local artiniano
\sref[I]{I.6.2.2};
observamos, então,
que para todo aberto
afim
$U'=\operatorname{Spec}(A')$
de $X'$, $A'$ é um
$A$-módulo plano,
portanto \emph{projetivo}
$(\mathbf 0_{\mathrm{III}},
10.1.3)$. Basta,
assim, aplicar
\sref{IV.11.9.7}
a cada aberto afim de
$X'$ para obter o
corolário.
\end{proof}

\begin{theorem}[11.9.10]
\label{IV.11.9.10}
Sejam $X$ um pré-esquema,
$(u_\lambda)$ uma família
de homomorfismos
\sref{IV.11.9.1.1},
$g:X'\to X$ um morfismo
e $(u'_\lambda)$ a
família de homomorfismos
deduzida de
$(u_\lambda)$ pela
mudança de base $g$
\sref{IV.11.9.5}.
\begin{enumerate}[label=(\roman*)]
  \item Se $g$ é um
  morfismo fielmente
  plano e a família
  $(u'_\lambda)$ é
  separante, então
  a família
  $(u_\lambda)$ é
  separante.
  \oldpage[IV-3]{164}
  \item Suponhamos que
  $g$ seja um morfismo
  plano e que, além
  disso, uma das duas
  condições seguintes
  seja satisfeita:
  \begin{enumerate}[label=\emph{\alph*)}]
    \item $L$ é finito
    e os $f_\lambda$
    são quase-compactos.
    \item O morfismo
    $g$ é localmente
    de apresentação
    finita.
  \end{enumerate}
  Então, se a família
  $(u_\lambda)$ é
  separante, o mesmo
  vale para a família
  $(u'_\lambda)$.
\end{enumerate}
\end{theorem}

\begin{proof}
\begin{enumerate}[label=(\roman*)]
  \item Em virtude de
  \sref{IV.2.2.8},
  basta mostrar que,
  se uma seção $t$ de
  $\sh{F}$ sobre um
  aberto $U$ de $X$
  pertence ao núcleo
  de cada $u_\lambda$,
  sua imagem $t'$ pelo
  homomorfismo canônico
  \[
  \Gamma(\rho):
  \Gamma(U,\sh{F})
  \longrightarrow
  \Gamma(g^{-1}(U),
  g^*(\sh{F}))
  =\Gamma(U,
  g_*(g^*(\sh{F})))
  \]
  é nula. Ora, as
  imagens de $t'$ pelos
  $g^*(u_\lambda)$ são
  as imagens dos
  $u_\lambda(t)$ pelo
  homomorfismo
  \[
  \Gamma(U,
  (f_\lambda)_*
  (\sh{G}_\lambda))
  \longrightarrow
  \Gamma(g^{-1}(U),
  g^*((f_\lambda)_*
  (\sh{G}_\lambda))),
  \]
  portanto são nulas;
  \emph{a fortiori},
  $u'_\lambda(t')=0$
  para todo $\lambda$,
  donde $t'=0$ pela
  hipótese. Isso
  prova (i).

  \item Como a questão
  é local em $X$ e
  $X'$, podemos nos
  limitar ao caso em
  que
  $X=\operatorname{Spec}(A)$
  e
  $X'=\operatorname{Spec}(A')$
  são afins, sendo
  $A'$ um $A$-módulo
  plano, e a provar
  que toda seção $z'$
  de $\sh{F}'$ sobre
  $X'$ cujas imagens
  por todos os
  $u'_\lambda$ são
  nulas é ela própria
  nula. Vimos ainda em
  \sref{IV.11.9.6}
  que podemos então
  supor que
  $Z_\lambda=X$ para
  todo $\lambda$.

  Distinguamos agora
  os dois casos.

  \emph{a)} Se $L$
  é finito e os
  $f_\lambda$ são
  quase-compactos,
  vimos em
  \sref{IV.11.9.3}
  que podemos ainda
  nos reduzir ao
  caso em que
  $Z_\lambda=X$
  para todo
  $\lambda\in L$ e,
  além disso, $L$
  é finito. Equivale
  então dizer que a
  interseção dos
  núcleos dos
  $u_\lambda:
  \sh{F}\to
  \sh{G}_\lambda$
  é nula, ou que o
  homomorfismo
  $u=(u_\lambda):
  \sh{F}\to
  \sh{G}=
  \bigoplus_\lambda
  \sh{G}_\lambda$
  é injetivo. Como
  $u'=g^*(u)$ é
  injetivo, pois
  $g$ é plano, a
  proposição fica
  demonstrada nesse
  caso.

  \emph{b)} Sejam
  $\sh{F}=\widetilde M$,
  $\sh{G}_\lambda=
  \widetilde N_\lambda$,
  e ponhamos
  $M'=M\otimes_A A'$,
  $N'_\lambda=
  N_\lambda\otimes_A A'$,
  de modo que
  $\sh{F}'=\widetilde{M'}$
  e
  $\sh{G}'_\lambda=
  \widetilde{N'_\lambda}$;
  por abuso de linguagem,
  continuaremos a denotar
  por $u_\lambda$ o
  homomorfismo
  $M\to N_\lambda$
  e por $u'_\lambda$
  o homomorfismo
  $u_\lambda\otimes1:
  M'\to N'_\lambda$.
  Escolhamos um elemento
  $z'\in M'$ tal que
  $u'_\lambda(z')=0$
  para todo $\lambda$;
  resta provar que
  $z'=0$. Ora, a hipótese
  de que $g$ é plano e
  de apresentação finita
  implica, segundo
  \sref{IV.11.3.15},
  que existe uma sequência
  finita
  $(s_i)_{1\leq i\leq n}$
  de elementos de $A$
  tal que, pondo
  $\mathfrak J_k=
  s_1A+\cdots+s_kA$
  e
  $A_i=A_{s_i}/
  \mathfrak J_{i-1}A_{s_i}$,
  o anel
  $A'_i=A'\otimes_A A_i$
  seja um $A_i$-módulo
  \emph{livre} para
  $1\leq i\leq n$,
  e $\mathfrak J_n=A$.
  A proposição estará
  demonstrada se provarmos,
  para $1\leq i\leq n$,
  a afirmação seguinte:
  \begin{enumerate}[label=$(*_i)$]
    \item \emph{Existe um
    inteiro $m_i>0$ tal
    que $s_j^{m_i}z'=0$
    para $j\leq i$.}
  \end{enumerate}
  Com efeito, pondo então
  $k=m_n$ e observando
  que os $s_i^k$
  $(1\leq i\leq n)$
  também geram o ideal
  unitário de $A$, a
  afirmação $(*_n)$
  mostrará que $z'$,
  combinação linear
  dos $s_i^kz'$, é nulo.

  Provemos $(*_i)$ por
  indução em $i$, sendo
  a afirmação vazia
  para $i=0$. Suponhamos
  $i>0$ e seja $m$ um
  múltiplo comum dos
  $m_j$ para $j\leq i-1$.
  Observemos que, para
  $1\leq h\leq n$, se
  $\mathfrak J'_h$ é o
  ideal gerado pelos
  $s_j^m$
  $(1\leq j\leq h)$\footnote{O
  francês transcrito e
  os testemunhos espanhol
  e inglês trazem
  $0\leq j\leq h$.
  A sequência dos $s_j$
  foi definida apenas
  para $1\leq j\leq n$;
  não há $s_0$ definido.
  O texto corrente usa
  $1\leq j\leq h$,
  preservando aqui a
  leitura testemunhal.},
  então
  $\mathfrak J_h/
  \mathfrak J'_h$ é
  nilpotente. Substituir
  os $s_j$ por $s_j^m$
  para $1\leq j\leq n$
  equivale, portanto,
  a substituir, para
  $1\leq h\leq n$,
  $A_h$ por
  $B_h=A_{s_h}/
  \mathfrak J'_{h-1}
  A_{s_h}$, de modo que
  $A_h=B_h/
  ((\mathfrak J_{h-1}/
  \mathfrak J'_{h-1})B_h)$;
  como $A'$ é um
  $A$-módulo plano,
  decorre de
  $(\mathbf 0_{\mathrm{III}},
  10.1.2)$ que
  $B'_h=A'\otimes_A B_h$
  ainda é um
  $B_h$-módulo livre.
  Podemos, portanto,
  \oldpage[IV-3]{165}
  substituir todos os
  $s_j$
  $(1\leq j\leq n)$
  por $s_j^m$ sem
  alterar as propriedades
  dos $\mathfrak J_h$
  e dos $A'_h$,
  e supor doravante
  que $m=1$. Então,
  como $z'$ é anulado
  por $\mathfrak J_{i-1}$,
  identifica-se com um
  elemento de
  $\operatorname{Hom}_{A'}
  (A'/\mathfrak J_{i-1}A',
  M')$; como
  $\mathfrak J_{i-1}A'$
  é um ideal de tipo
  finito de $A'$,
  esse módulo de
  homomorfismos se
  identifica, por sua
  vez, com
  $\operatorname{Hom}_A
  (A/\mathfrak J_{i-1},M)
  \otimes_A A'$
  $(\mathbf 0_{\mathrm I},
  6.2.2)$. Seja
  \[
  v_\lambda=
  \operatorname{Hom}(1,u_\lambda):
  \operatorname{Hom}_A
  (A/\mathfrak J_{i-1},M)
  \longrightarrow
  \operatorname{Hom}_A
  (A/\mathfrak J_{i-1},
  N_\lambda)
  \]
  o homomorfismo deduzido
  de $u_\lambda$; a
  família $(v_\lambda)$
  também é
  \emph{separante}.
  Com efeito, para
  todo $t\in A$,
  \[
  (\operatorname{Hom}_A
  (A/\mathfrak J_{i-1},M))_t
  =
  \operatorname{Hom}_{A_t}
  (A_t/(\mathfrak J_{i-1})_t,
  M_t)
  \]
  e o mesmo vale com
  $M$ substituído por
  $N_\lambda$, pois
  o ideal
  $\mathfrak J_{i-1}$
  é de tipo finito
  $(\mathbf 0_{\mathrm I},
  1.3.5)$; como, por
  hipótese, a interseção
  dos núcleos dos
  $(u_\lambda)_t$ é
  nula, o mesmo vale
  para a interseção
  dos núcleos dos
  $(v_\lambda)_t=
  \operatorname{Hom}
  (1,(u_\lambda)_t)$,
  donde a afirmação
  \sref{IV.11.9.4}.
  Substituindo $A$ por
  $A/\mathfrak J_{i-1}$,
  $M$ por
  $\operatorname{Hom}_A
  (A/\mathfrak J_{i-1},M)$,
  $N_\lambda$ por
  $\operatorname{Hom}_A
  (A/\mathfrak J_{i-1},
  N_\lambda)$
  (que são
  $(A/\mathfrak J_{i-1})$-módulos),
  $u_\lambda$ por
  $v_\lambda$ e,
  enfim, $A'$ por
  $A'/\mathfrak J_{i-1}A'$,
  vemos que podemos
  nos reduzir ao caso
  em que, na situação
  inicial, o elemento
  $s=s_i\in A$ é tal
  que $A'_s$ é um
  $A_s$-módulo livre.
  Ora, a família dos
  $(u_\lambda)_s:
  M_s\to(N_\lambda)_s$
  é separante por
  hipótese; como
  $(u'_\lambda)_s
  (z'/1)=0$, decorre
  de \sref{IV.11.9.7}
  que $z'/1=0$ em
  $M'_s$. Isso significa
  que existe um inteiro
  $r$ tal que
  $s^rz'=0$ em $M'$,
  encerrando a prova
  de $(*_i)$ por indução.
\end{enumerate}
\end{proof}

\begin{remark}[11.9.11]
\label{IV.11.9.11}
Limitamo-nos, para
simplificar, ao caso
em que $Z_\lambda=X$
para todo $\lambda$.
Convém notar que,
se a família dos
homomorfismos
$u_\lambda:
\sh{F}\to\sh{G}_\lambda$
é separante, disso
\emph{não decorre
necessariamente} que,
para todo $x\in X$,
a interseção dos
núcleos dos
homomorfismos
$(u_\lambda)_x:
\sh{F}_x\to
(\sh{G}_\lambda)_x$
se reduza a $0$.

Por exemplo, sejam
$X$ um pré-esquema
de Jacobson localmente
noetheriano, de
dimensão $\geq1$,
e $\sh{F}$ um
$\sh{O}_X$-módulo
coerente. Para cada
ponto fechado $x$
de $X$ e cada inteiro
$n\geq0$,
$\sh{F}/\mathfrak m_x^n
\sh{F}$ é um
$\sh{O}_X$-módulo
coerente cujo suporte
está contido em
$\{x\}$.\footnote{O
impresso francês e
as versões espanhola
e inglesa dizem aqui
$\mathfrak m_x^n\sh{F}$,
sem o quociente. O
suporte dessa
subestrutura não
precisa estar contido
em $\{x\}$; para
$n=0$, ela é
$\sh{F}$. A família
de quocientes da frase
seguinte fixa a leitura
corrigida.}
A família dos
homomorfismos canônicos
$\sh{F}\to
\sh{F}/\mathfrak m_x^n
\sh{F}$ (em que $x$
percorre o conjunto
$X_0$ dos pontos
fechados de $X$
e $n$ percorre os
inteiros $\geq0$)
é \emph{separante}.
Com efeito, seja $t$
uma seção de $\sh{F}$
sobre um aberto $U$
de $X$ cujas imagens
em todos os
$\Gamma(U,\sh{F}/
\mathfrak m_x^n\sh{F})$
sejam nulas. Segue-se
imediatamente que
$t_x=0$ para todo
ponto fechado
$x\in U$. Como o
conjunto dos pontos
fechados contidos
em $U$ é muito denso
em $U$, daí decorre
$t=0$
\sref{IV.10.2.1}.

Entretanto, se
$\sh{F}=\sh{O}_X$
e $y\in X$ é um
ponto \emph{não
fechado}, então
$(\sh{O}_X/
\mathfrak m_x^n
\sh{O}_X)_y=0$
para todo ponto
fechado $x$ de $X$,
ao passo que
$\sh{O}_{X,y}\ne0$.
Assim, a interseção
dos núcleos dos
homomorfismos
$\sh{O}_{X,y}\to
(\sh{O}_X/
\mathfrak m_x^n
\sh{O}_X)_y$
é igual a
$\sh{O}_{X,y}$.
\end{remark}

\begin{lemma}[11.9.12]
\label{IV.11.9.12}
Com as notações de
\sref{IV.11.9.1} e
\sref{IV.11.9.5},
suponhamos que a família
$(u_\lambda)$ seja
separante. Suponhamos,
além disso, que $X$
seja um $S$-pré-esquema,
com
$S=\operatorname{Spec}(A)$
afim, e que
$X'=X\times_S S'$,
com
$S'=\operatorname{Spec}(A')$
e $A'$ uma $A$-álgebra.
Suponhamos, por fim, que
os $\sh{G}_\lambda$
sejam planos sobre $S$.
Seja
$(A'_\alpha)_{\alpha\in I}$
a família filtrante das
sub-$A$-álgebras de
tipo finito de $A'$.
Para cada $\alpha\in I$,
ponhamos
$S''_\alpha=
\operatorname{Spec}(A'_\alpha)$,
$X''_\alpha=
X\times_S S''_\alpha$
e
$Z''_{\alpha\lambda}=
Z_\lambda\times_X
X''_\alpha$.
Sejam
$f''_{\alpha\lambda}:
Z''_{\alpha\lambda}
\to X''_\alpha$,
$\sh{F}''_\alpha$,
$\sh{G}''_{\alpha\lambda}$
e $u''_{\alpha\lambda}$
os morfismos, módulos
e homomorfismos de
módulos obtidos de
$f_\lambda$, $\sh{F}$,
$\sh{G}_\lambda$ e
$u_\lambda$ mediante
a mudança de base
$X''_\alpha\to X$.
Se, para todo
$\alpha\in I$, a família
$(u''_{\alpha\lambda})
_{\lambda\in L}$ é
separante, então
$(u'_\lambda)_{\lambda\in L}$
também é separante.
\end{lemma}

\begin{proof}
Devemos provar que,
se $t'$ é uma seção de
$\sh{F}'$ sobre um
aberto afim $U'$ de
$X'$ e
$u'_\lambda(t')=0$
para todo
$\lambda\in L$,
então $t'=0$.
Se
$h_\alpha:X'\to
X''_\alpha$ é a
projeção canônica,
\oldpage[IV-3]{166}
segue de
\sref{IV.8.2.11}
que existem um índice
$\alpha$ e um aberto
quase-compacto
$U''_\alpha\subset
X''_\alpha$ tais que
$U'=h_\alpha^{-1}
(U''_\alpha)$.
Além disso, em virtude
de \sref{IV.8.5.2},~(i),
podemos supor que $t'$
é a imagem canônica
de uma seção
$t''_\alpha$ de
$\sh{F}''_\alpha$
sobre $U''_\alpha$.

Substituindo, se
necessário, $S$, $X$,
$Z_\lambda$,
$f_\lambda$,
$\sh{F}$,
$\sh{G}_\lambda$ e
$u_\lambda$ por
$S''_\alpha$,
$U''_\alpha$,
$(f''_{\alpha\lambda})
^{-1}(U''_\alpha)$
e pelas restrições
correspondentes de
$\sh{F}''_\alpha$,
$\sh{G}''_{\alpha\lambda}$
e
$u''_{\alpha\lambda}$,
podemos, portanto,
supor que $U'=X'$,
que $t'$ é a imagem
canônica de uma seção
$t$ de $\sh{F}$
sobre $X$ e que o
homomorfismo
$A\to A'$ é injetivo.
Equivalentemente,
se $p:S'\to S$ é o
morfismo correspondente,
o homomorfismo
$\sh{O}_S\to
p_*(\sh{O}_{S'})$
que intervém na
definição de $p$
é injetivo.

Da planicidade de
$\sh{G}_\lambda$
sobre $S$ decorre
imediatamente, por
exemplo reduzindo-nos
ao caso em que
$Z_\lambda$ é afim
sobre $S$
(\sref[I]{I.1.6.3}
e \sref[I]{I.1.6.5}),
que o homomorfismo
canônico
$\rho:\sh{G}_\lambda
\to(g'_\lambda)_*
(\sh{G}'_\lambda)$
também é injetivo.
Mas o homomorfismo
composto
\[
\Gamma(X,\sh{F})
\xrightarrow{\Gamma(u_\lambda)}
\Gamma(Z_\lambda,\sh{G}_\lambda)
\xrightarrow{\Gamma(\rho)}
\Gamma(Z'_\lambda,\sh{G}'_\lambda)
\]
é, por definição,
igual ao composto
\[
\Gamma(X,\sh{F})
\xrightarrow{\Gamma(\rho)}
\Gamma(X',\sh{F}')
\xrightarrow{\Gamma(u'_\lambda)}
\Gamma(Z'_\lambda,\sh{G}'_\lambda).
\]
Logo, a imagem de
$t$ por esses
homomorfismos
compostos é
$u'_\lambda(t')=0$.
Pela injetividade de
$\Gamma(Z_\lambda,
\sh{G}_\lambda)
\xrightarrow{\Gamma(\rho)}
\Gamma(Z'_\lambda,
\sh{G}'_\lambda)$,
concluímos que
$u_\lambda(t)=0$
para todo
$\lambda\in L$.
Da hipótese segue
$t=0$ e, por fim,
$t'=0$.
\end{proof}

\begin{proposition}[11.9.13]
\label{IV.11.9.13}
Com as notações de
\sref{IV.11.9.1} e
\sref{IV.11.9.5},
suponhamos que $X$
seja um pré-esquema
sobre um corpo $k$ e
que, pondo
$S=\operatorname{Spec}(k)$,
tenhamos
$X'=X\times_S S'$,
onde $S'$ é um
$k$-pré-esquema
arbitrário. Se a
família
$(u_\lambda)_{\lambda\in L}$
é separante, então
$(u'_\lambda)$
também é separante.
\end{proposition}

\begin{proof}
Podemos limitar-nos
ao caso em que
$S'=\operatorname{Spec}(A')$
é afim. Se $A'$ é uma
$k$-álgebra de tipo
finito, o morfismo
$g:X'\to X$ é plano
e de apresentação
finita; aplicamos,
portanto,
\sref{IV.11.9.10},~(ii),~$b$).
No caso geral,
consideramos $A'$
como limite indutivo
de suas sub-$k$-álgebras
$A'_\alpha$ de tipo
finito e aplicamos
a cada $A'_\alpha$
o resultado de
\sref{IV.11.9.10},~(ii),~$b$).
Concluímos então
pelo lema
\sref{IV.11.9.12},
pois os
$\sh{G}_\lambda$
são planos sobre $S$.
\end{proof}

\begin{env}[11.9.14]
\label{IV.11.9.14}
Conservemos as
notações de
\sref{IV.11.9.1} e
\sref{IV.11.9.5},
e suponhamos que
$X$ seja um
$S$-pré-esquema.
Se, para toda mudança
de base
$g:X\times_S S'\to X$,
em que $S'\to S$ é
um morfismo
arbitrário, a família
correspondente
$(u'_\lambda)$ é
separante, diremos
que $(u_\lambda)$ é
\emph{universalmente
separante relativamente
a} $S$. Quando a
família $(u_\lambda)$
tem um único elemento
$u$, diremos também
que $u$ é
\emph{universalmente
injetivo relativamente
a} $S$.

É claro que, para
todo morfismo
$h:S'\to S$, a família
correspondente
$(u'_\lambda)$ é
universalmente
separante relativamente
a $S'$. Reciprocamente,
se $h$ é
\emph{fielmente plano}
e $(u'_\lambda)$ é
universalmente
separante relativamente
a $S'$, então
$(u_\lambda)$ é
universalmente
separante relativamente
a $S$. Isso decorre
imediatamente de
\sref{IV.11.9.10},~(i)
e do fato de que,
para todo morfismo
$S''\to S$, o morfismo
correspondente
$S''\times_S S'\to S''$
é fielmente plano.
\end{env}

\oldpage[IV-3]{167}

\begin{proposition}[11.9.15]
\label{IV.11.9.15}
Com as notações de
\sref{IV.11.9.1},
suponhamos que $X$
seja um $S$-pré-esquema
e que os
$\sh{G}_\lambda$
sejam planos sobre $S$.
Seja $S_0$ um
subpré-esquema fechado
de $S$, definido por
um ideal quase-coerente
nilpotente $\sh{J}$
de $\sh{O}_S$, tal que
os
$(\sh{O}_S/\sh{J})$-módulos
$\sh{J}^k/
\sh{J}^{k+1}$ sejam
todos localmente
livres. Seja
$(u_{\lambda0})$
a família de
homomorfismos obtida
de $(u_\lambda)$
pela mudança de base
$X\leftarrow X_0=
X\times_S S_0$.
Se $(u_{\lambda0})$
é separante
(respectivamente,
universalmente
separante relativamente
a $S_0$), então
$(u_\lambda)$ é
separante
(respectivamente,
universalmente
separante relativamente
a $S$).
\end{proposition}

\begin{proof}
Notemos que, para
qualquer mudança de
base $S'\to S$, se
$S'_0=S'\times_S S_0$,
então $S'_0$ é um
subpré-esquema fechado
de $S'$ definido por
um ideal quase-coerente
nilpotente $\sh{J}'$
de $\sh{O}_{S'}$.
Além disso,
$\sh{J}'^k/
\sh{J}'^{k+1}$
é um
$(\sh{O}_{S'}/
\sh{J}')$-módulo
localmente livre para
todo $k$
\sref{IV.2.1.8},~(i).
Como os
$\sh{G}'_\lambda=
\sh{G}_\lambda
\otimes_{\sh{O}_S}
\sh{O}_{S'}$
são planos sobre
$S'$, a afirmação
para famílias
universalmente
separantes decorre
da afirmação para
famílias separantes.

Para provar esta
última, podemos,
por \sref{IV.11.9.3},
reduzir-nos ao caso
em que
$S=\operatorname{Spec}(A)$
e
$X=\operatorname{Spec}(B)$
são afins,
$Z_\lambda=X$ para
todo $\lambda$,
$\sh{F}=\widetilde M$
e
$\sh{G}_\lambda=
\widetilde N_\lambda$,
com $M$ e
$N_\lambda$
$B$-módulos e
$N_\lambda$ planos
como $A$-módulos.
Como a questão é
local em $X$ e em
$S$, basta mostrar
que, se $t\in M$
satisfaz
$u_\lambda(t)=0$
para todo $\lambda$,
então $t=0$.

Temos
$\sh{J}=
\widetilde{\mathfrak J}$,
onde $\mathfrak J$
é um ideal nilpotente
de $A$ cujos quocientes
$\mathfrak J^k/
\mathfrak J^{k+1}$
são
$(A/\mathfrak J)$-módulos
livres. Por hipótese,
os homomorfismos
$u_{\lambda0}:
M/\mathfrak JM\to
N_\lambda/
\mathfrak JN_\lambda$
formam uma família
separante. Suponhamos
$\mathfrak J^{n+1}=0$,
com $n\geq0$ inteiro,
e procedamos por
indução em $n$; para
$n=0$ a afirmação
é imediata.

Se $\bar t$ é a
classe de $t$ em
$M/\mathfrak J^nM$,
a classe de
$u_\lambda(t)$ em
$N_\lambda/
\mathfrak J^nN_\lambda$
é nula para todo
$\lambda$. Pela
hipótese de indução,
$\bar t=0$; em outras
palavras,
$t\in\mathfrak J^nM$.
Ora,
$\mathfrak J^n=
\mathfrak J^n/
\mathfrak J^{n+1}$
é um
$(A/\mathfrak J)$-módulo
livre. Se
$(e_\alpha)$ é uma
base desse módulo,
podemos escrever
$t=\sum_\alpha
e_\alpha t^0_\alpha$,
com
$t^0_\alpha\in
M/\mathfrak JM$ e
$t^0_\alpha=0$ salvo
para um número finito
de índices.

Por outro lado,
como $N_\lambda$
é um $A$-módulo
plano,
$\mathfrak J^n
N_\lambda$
identifica-se com
$\mathfrak J^n
\otimes_{A/\mathfrak J}
(N_\lambda/
\mathfrak JN_\lambda)$.
Podemos, portanto,
escrever
\[
u_\lambda(t)=
\sum_\alpha
e_\alpha
u_{\lambda0}(t^0_\alpha)
=
\sum_\alpha
e_\alpha\otimes
u_{\lambda0}(t^0_\alpha).
\]
Como, por hipótese,
$u_\lambda(t)=0$,
segue que
$u_{\lambda0}(t^0_\alpha)
=0$ para todo
$\alpha$ e todo
$\lambda$. A família
$(u_{\lambda0})$
sendo separante,
$t^0_\alpha=0$
para todo $\alpha$
e, por conseguinte,
$t=0$.
\end{proof}

\begin{theorem}[11.9.16]
\label{IV.11.9.16}
Com as notações de
\sref{IV.11.9.1},
suponhamos que $X$
seja um $S$-pré-esquema
localmente noetheriano,
que $\sh{F}$ seja um
$\sh{O}_X$-módulo
coerente e que os
$\sh{G}_\lambda$
sejam planos sobre
$S$. Para cada
$s\in S$, seja
$((u_\lambda)_s)
_{\lambda\in L}$
a família obtida de
$(u_\lambda)$ pela
mudança de base
$X_s=X\times_S
\operatorname{Spec}(k(s))
\to X$. Então a
família $(u_\lambda)$
é universalmente
separante relativamente
a $S$ se, e somente
se, para todo $s\in S$,
a família
$((u_\lambda)_s)$
é separante.
\end{theorem}

\begin{proof}
A necessidade decorre
imediatamente das
definições.
Reciprocamente,
suponhamos satisfeita
a condição do
enunciado e provemos
primeiro que
$(u_\lambda)$ é
separante. Por
\sref{IV.11.9.3},
podemos reduzir-nos
ao caso afim
$S=\operatorname{Spec}(A)$,
$X=\operatorname{Spec}(B)$,
com $Z_\lambda=X$
para todo $\lambda$,
$\sh{F}=\widetilde M$
e
$\sh{G}_\lambda=
\widetilde N_\lambda$,
onde $B$ é
noetheriano, $M$ é
um $B$-módulo de
tipo finito e os
$N_\lambda$ são
$A$-módulos planos.
Basta provar que,
se $t\in M$ e
$u_\lambda(t)=0$
para todo $\lambda$,
então $t=0$.
Para isso, basta
mostrar que, para
todo ideal maximal
$\mathfrak p$ de $B$,
a imagem
$t_\mathfrak p$
de $t$ em
$M_\mathfrak p$
é nula (Bourbaki,
\emph{Alg.\ comm.},
cap.~II, \S~3,
n.~3, cor.~1 do
teorema~1).
Podemos, portanto,
limitar-nos a provar
que a interseção
\oldpage[IV-3]{168}
dos núcleos dos
$(u_\lambda)_\mathfrak p:
M_\mathfrak p\to
(N_\lambda)_\mathfrak p$,
obtidos de
$(u_\lambda)$ pela
mudança de base
$\operatorname{Spec}
(B_\mathfrak p)
\to\operatorname{Spec}(B)$,
é nula.

Em outras palavras,
reduzimo-nos ao caso
em que $B$ é um
anel local noetheriano.
Considerando em $A$
a imagem inversa do
ideal maximal de
$B$, podemos também
supor que $A$ é
um anel local, com
ideal maximal
$\mathfrak m$.
Como $\mathfrak mB$
está contido no
ideal maximal de
$B$ e $M$ é um
$B$-módulo de tipo
finito, a interseção
dos
$\mathfrak m^nM$
é nula
\sref[0]{0.7.3.5}.
Basta, pois, provar
que, para cada $n$,
a imagem de $t$
em
$M/\mathfrak m^{n+1}M$
é nula. Assim,
basta mostrar que
a família obtida
de $(u_\lambda)$
pela mudança de base
$\operatorname{Spec}
(B/\mathfrak m^{n+1}B)
\to\operatorname{Spec}(B)$
é separante; podemos
então limitar-nos
ao caso em que $A$
é local e seu ideal
maximal $\mathfrak m$
é nilpotente. Nesse
caso, os
$\mathfrak m^k/
\mathfrak m^{k+1}$
são
$(A/\mathfrak m)$-módulos
livres. Pela hipótese
sobre as famílias
$((u_\lambda)_s)$,
aplicamos precisamente
\sref{IV.11.9.15},
obtendo a conclusão
anunciada.

Seja agora
$h:S'\to S$ um
morfismo de mudança
de base, e seja
$(u'_\lambda)$
a família obtida
de $(u_\lambda)$
pela mudança de base
$h':X\times_S S'\to X$.
Provemos que
$(u'_\lambda)$
também é separante.
Suponhamos primeiro
que $h$ seja
\emph{localmente de
tipo finito}. Então
$h'$ também o é,
e
$X'=X\times_S S'$
é localmente
noetheriano. Além
disso, se $s'\in S'$
está sobre $s\in S$,
segue de
\sref{IV.11.9.13},
aplicada a $X_s$ e
a
$X_{s'}=
X_s\otimes_{k(s)}k(s')$,
que
$((u'_\lambda)_{s'})$
é separante para
todo $s'\in S'$.
Pela primeira parte
da demonstração,
$(u'_\lambda)$
é separante nesse
caso.

Por fim, se $h$ é
arbitrário, podemos
limitar-nos ao caso
afim
$S=\operatorname{Spec}(A)$,
$S'=\operatorname{Spec}(A')$
e considerar $A'$
como limite indutivo
de suas sub-$A$-álgebras
de tipo finito.
Como os
$\sh{G}_\lambda$
são planos sobre
$S$, aplicamos
o resultado anterior
e o lema
\sref{IV.11.9.12}
\footnote{O impresso
francês remete aqui
ao lema 11.9.5;
11.9.5, porém, não
é um lema de passagem
ao limite. O lema
pertinente é 11.9.12,
como já leem as
versões espanhola
e inglesa.}
para concluir a
demonstração.
\end{proof}

\begin{proposition}[11.9.17]
\label{IV.11.9.17}
Sejam $f:X\to S$
um morfismo localmente
de apresentação finita,
$\sh{F}$ um
$\sh{O}_X$-módulo
quase-coerente de
apresentação finita
e $f$-plano,
$U$ um aberto de
$X$ e $j:U\to X$
a imersão canônica.
Seja
$u:\sh{F}\to
j_*(j^*(\sh{F}))$
o homomorfismo
canônico
\sref[0]{0.4.4.3.2}.
Para cada $s\in S$,
sejam
$X_s=X\times_S
\operatorname{Spec}(k(s))$
a fibra,
$U_s=U\cap X_s$
o aberto de $X_s$,
$\sh{F}_s=
\sh{F}\otimes_{\sh{O}_S}
k(s)$,
$j_s:U_s\to X_s$
a imersão canônica
e
$u_s:\sh{F}_s\to
(j_s)_*((j_s)^*
(\sh{F}_s))$
o homomorfismo
canônico. Então
$u$ é universalmente
injetivo relativamente
a $S$ se, e somente
se, $u_s$ é injetivo
para todo $s\in S$.
\end{proposition}

\begin{proof}
Resta provar a
suficiência da
condição. Quando
$X$ é localmente
noetheriano, a
proposição decorre
imediatamente de
\sref{IV.11.9.16}.
Reduziremos o caso
geral a esse em
duas etapas,
limitando-nos, como
é permitido, ao caso
em que
$S=\operatorname{Spec}(A)$
e $X$ são afins.

\medskip
\emph{A) Caso em que
$U$ é quase-compacto.}
Usaremos o lema
seguinte.

\begin{lemma}[11.9.17.1]
\label{IV.11.9.17.1}
Sob as hipóteses gerais
de \sref{IV.11.9.17},
supondo ainda que $S$
e $X$ sejam afins e
que $U$ seja
quase-compacto, o
conjunto $E$ dos
$s\in S$ para os quais
$u_s$ é injetivo é
construtível.
\end{lemma}

\begin{proof}
Com efeito, as fibras
$X_s$ são
pré-esquemas
localmente
noetherianos.
Portanto, por
\sref{IV.5.10.2},
$E$ também é o
conjunto dos
$s\in S$ tais que
$\operatorname{Ass}
(\sh{F}_s)\subset U_s$.
Além disso, $U$,
por ser quase-compacto
em um esquema afim,
é construtível.
Assim, a condição
\sref{IV.9.2.1},~(i)
decorre imediatamente
de \sref{IV.4.2.7}.
Por outro lado, a
\oldpage[IV-3]{169}
condição
\sref{IV.9.2.1},~(ii)
é verificada mediante
o estudo dos ciclos
primos associados
na vizinhança da
fibra genérica
\sref{IV.9.8.3},
bem como
\sref{IV.9.5.2}
e \sref{IV.9.5.3}.
\end{proof}

Estabelecido o lema,
podemos, em virtude
de \sref{IV.8.9.1}
e \sref{IV.8.2.11},
supor que existem
um subanel noetheriano
$A_0$ de $A$,
um pré-esquema
$X_0$ de tipo
finito sobre
$S_0=\operatorname{Spec}(A_0)$,
um aberto $U_0$
de $X_0$ e um
$\sh{O}_{X_0}$-módulo
coerente $\sh{F}_0$
tais que
$X=X_0\times_{S_0}S$
e, se
$p_0:X\to X_0$
é a projeção
canônica, temos
$U=p_0^{-1}(U_0)$
e
$\sh{F}=p_0^*
(\sh{F}_0)$.

Seja $(A_\alpha)$
a família filtrante
dos subanéis de
$A$ que são
$A_0$-álgebras
de tipo finito,
e ponhamos
$S_\alpha=
\operatorname{Spec}(A_\alpha)$,
$X_\alpha=
X_0\times_{S_0}S_\alpha$,
$U_\alpha=
U_0\times_{S_0}S_\alpha$
e
$\sh{F}_\alpha=
\sh{F}_0\otimes_{S_0}
S_\alpha$.
Seja $u_\alpha$
o homomorfismo
canônico relativo
a $\sh{F}_\alpha$
e $U_\alpha$,
definido como em
\sref{IV.11.9.17}.
Para cada $s\in S$,
a injetividade de
$u_s$ implica a de
$(u_\alpha)_{s_\alpha}$
\sref{IV.11.9.10},~(i),
onde
$s_\alpha=q_\alpha(s)$
é a imagem de $s$
pelo morfismo
$q_\alpha:S\to
S_\alpha$.
Se $E_\alpha$ é
o conjunto dos
$s_\alpha\in S_\alpha$
para os quais
$(u_\alpha)_{s_\alpha}$
é injetivo, temos
$S=q_\alpha^{-1}
(E_\alpha)$, e
os $E_\alpha$
formam um sistema
projetivo de conjuntos.
Pelo lema
\sref{IV.11.9.17.1},
aplicado a $u_\alpha$,
$E_\alpha$ é
construtível em
$S_\alpha$. Logo,
por \sref{IV.8.3.4},
existe um índice
$\alpha$ tal que
$E_\alpha=S_\alpha$.
Como $X_\alpha$
é noetheriano,
$u_\alpha$ é
universalmente
injetivo por
\sref{IV.11.9.16};
o mesmo vale,
portanto, para $u$.

\medskip
\emph{B) Caso geral.}
O aberto $U$ é
uma reunião filtrante
crescente de abertos
quase-compactos
$U_\lambda$. Se
$j_\lambda:
U_\lambda\to X$
é a imersão canônica
e
$u_\lambda:
\sh{F}\to
(j_\lambda)_*
(j_\lambda^*(\sh{F}))$
o homomorfismo
correspondente,
o lema
\sref{IV.11.9.17.1},
aplicado a
$u_\lambda$, mostra
que o conjunto
$E_\lambda$ dos
$s\in S$ para os quais
$(u_\lambda)_s$ é
injetivo é
construtível.

Por outro lado,
para cada $s\in S$,
$u_s$ ser injetivo
significa que
$\operatorname{Ass}
(\sh{F}_s)
\subset U\cap X_s$;
de modo análogo,
$(u_\lambda)_s$
ser injetivo
significa que
$\operatorname{Ass}
(\sh{F}_s)
\subset U_\lambda
\cap X_s$
\sref{IV.5.10.2}.
Como
$\operatorname{Ass}
(\sh{F}_s)$ é
finito
\sref{IV.3.1.6},
a injetividade de
$u_s$ equivale à
existência de
$\lambda$ com
$s\in E_\lambda$.
Assim, a hipótese
significa que
$S=\bigcup_\lambda
E_\lambda$. Em
virtude de
\sref{IV.1.9.9},
existe um índice
$\lambda$ tal que
$S=E_\lambda$.
Pela primeira parte
do argumento,
$u_\lambda$ é
universalmente
injetivo. Por fim,
a fatoração de
$u_\lambda$
\[
\sh{F}
\xrightarrow{u}
j_*(\sh{F}|_U)
\longrightarrow
(j_\lambda)_*
(\sh{F}|_{U_\lambda})
\]
mostra que $u$
também é
universalmente
injetivo.
\end{proof}

\begin{remark}[11.9.18]
\label{IV.11.9.18}
Sejam $X$ um
pré-esquema,
$\sh{F}$ e
$\sh{G}$ dois
$\sh{O}_X$-módulos
quase-coerentes de
apresentação finita,
com $\sh{G}$ também
\emph{localmente
livre}, e seja
$u:\sh{F}\to\sh{G}$
um homomorfismo.
As condições seguintes
são equivalentes:
\begin{enumerate}[label=\emph{\alph*)}]
\item Para todo morfismo
$g:X'\to X$, o
homomorfismo
$g^*(u):
g^*(\sh{F})\to
g^*(\sh{G})$
é injetivo.
\item Para todo
$x\in X$, o
homomorfismo
$u\otimes1_{k(x)}:
\sh{F}\otimes_{\sh{O}_X}
k(x)\to
\sh{G}\otimes_{\sh{O}_X}
k(x)$
é injetivo.
\item Para todo
$x\in X$, existe
uma vizinhança aberta
$U$ de $x$ tal que
$u|_U:
\sh{F}|_U\to
\sh{G}|_U$
é um isomorfismo
de $\sh{F}|_U$
sobre um somando
direto de
$\sh{G}|_U$.
\end{enumerate}

É claro que
\emph{c)} implica
\emph{a)}, e
\emph{a)} implica
\emph{b)}. A
implicação
\emph{b)} para
\emph{c)} decorre
de \sref[0]{0.19.1.12},
\sref[0]{0.5.2.5}
e \sref[0]{0.5.5.5}.

Quando essas
condições equivalentes
são satisfeitas,
dizemos que $u$
é \emph{universalmente
injetivo}.
\end{remark}

\oldpage[IV-3]{170}

\subsection{Famílias esquematicamente dominantes de morfismos e famílias
esquematicamente densas de subpré-esquemas}
\label{subsection:IV.11.10}

\begin{proposition}[11.10.1]
\label{IV.11.10.1}
Sejam $X$ um
pré-esquema,
$(Z_\lambda)_{\lambda\in L}$
uma família de
pré-esquemas e,
para cada
$\lambda\in L$,
$f_\lambda=
(\psi_\lambda,
\theta_\lambda):
Z_\lambda\to X$
um morfismo.
As condições
seguintes são
equivalentes:
\begin{enumerate}[label=\emph{\alph*)}]
\item A família de
homomorfismos
$\theta_\lambda:
\sh{O}_X\to
(f_\lambda)_*
(\sh{O}_{Z_\lambda})$
é separante; em
outras palavras
\sref{IV.11.9.1},
a interseção dos
núcleos dos
$\theta_\lambda$
é nula.
\item Para todo
aberto $U$ de $X$,
toda seção $t$
de $\sh{O}_X$
sobre $U$ cujas
imagens por todos
os homomorfismos
canônicos
\begin{equation}
(\theta_\lambda)_U:
\Gamma(U,\sh{O}_X)
\longrightarrow
\Gamma(f_\lambda^{-1}(U),
\sh{O}_{Z_\lambda})
\label{IV.11.10.1.1}
\tag{11.10.1.1}
\end{equation}
são nulas é ela
própria nula.
\item Para todo
aberto $U$ de $X$
e todo
subpré-esquema
fechado $Y$ de $U$
tal que, para cada
$\lambda\in L$,
exista uma fatoração
\begin{equation}
f_\lambda^{-1}(U)
\xrightarrow{g_\lambda}
Y\xrightarrow{j}U
\label{IV.11.10.1.2}
\tag{11.10.1.2}
\end{equation}
da restrição
$f_\lambda^{-1}(U)
\to U$ de
$f_\lambda$, onde
$j$ é a imersão
canônica, temos
$Y=U$.
\end{enumerate}

Se, além disso,
$X$ é um
$S$-pré-esquema,
essas condições
também são
equivalentes à
seguinte:
\begin{enumerate}[label=\emph{\alph*)}]
\setcounter{enumi}{3}
\item Para todo
morfismo separado
$g:X'\to S$ e
todo par de
$S$-morfismos
$u,v:U\to X'$
de um aberto
$U$ de $X$,
se, para cada
$\lambda$, os
compostos de
$u$ e $v$ com
$f_\lambda^{-1}(U)
\to U$, restrição
de $f_\lambda$,
são iguais,
então $u=v$.
\end{enumerate}
\end{proposition}

\begin{proof}
A equivalência de
\emph{a)} e
\emph{b)} decorre
das definições.
Para mostrar que
\emph{b)} implica
\emph{c)}, basta
considerar o ideal
quase-coerente
$\sh{J}$ de
$\sh{O}_U$ que
define $Y$ e
observar que,
para todo aberto
$V\subset U$,
a hipótese faz
$(\theta_\lambda)_V$
fatorar-se como
\[
\Gamma(V,\sh{O}_U)
\longrightarrow
\Gamma(V,\sh{O}_U/\sh{J})
\longrightarrow
\Gamma(f_\lambda^{-1}(V),
\sh{O}_{Z_\lambda}).
\]
Portanto, toda
seção $t$ de
$\sh{J}$ sobre
$V$ tem imagem
nula em todos os
$\Gamma(
f_\lambda^{-1}(V),
\sh{O}_{Z_\lambda})$.
Por \emph{b)},
$\sh{J}=0$ e
$Y=U$.

Reciprocamente,
se \emph{c)}
vale, aplicamo-la,
para provar
\emph{b)}, ao
subpré-esquema
fechado $Y$ de
$U$ definido
pelo ideal
$t\sh{O}_U$.
Se todas as
imagens de $t$
pelos
$(\theta_\lambda)_U$
são nulas, as
fatorações
\sref{IV.11.10.1.2}
existem para
todo $\lambda$
\sref[I]{I.4.1.9}.

Para provar que
\emph{c)} implica
\emph{d)}, aplicamos
\emph{c)} ao
subpré-esquema
fechado $Y$
de $U$ obtido
como imagem
inversa, por
$(u,v)_S$, da
diagonal de
$X'$ em
$X'\times_S X'$,
e usamos
\sref[I]{I.4.4.1}.
Reciprocamente,
\emph{d)} implica
\emph{b)}:
consideramos o
$S$-esquema
$X'=S[T]=
S\otimes_{\mathbf Z}
\mathbf Z[T]$,
com $T$
indeterminada.
As seções de
$\sh{O}_U$ sobre
$U$ correspondem
biunivocamente
aos
$S$-morfismos
$U\to S[T]$
\sref[I]{I.3.3.15}.
Duas seções de
$\sh{O}_U$ têm
as mesmas imagens
por todos os
$(\theta_\lambda)_U$
se, e somente se,
os dois morfismos
correspondentes
têm compostos
iguais com todos
os morfismos
$f_\lambda^{-1}(U)
\to U$.
\end{proof}

\oldpage[IV-3]{171}

\begin{definition}[11.10.2]
\label{IV.11.10.2}
Quando as condições
equivalentes de
\sref{IV.11.10.1}
são satisfeitas,
dizemos que a família
$(f_\lambda)$ é
\emph{esquematicamente
dominante}. Quando os
$f_\lambda$ são as
imersões canônicas
em $X$ de uma família
$(Z_\lambda)$ de
subpré-esquemas de
$X$, dizemos também
que a família
$(Z_\lambda)$ é
\emph{esquematicamente
densa}.
\end{definition}

\begin{remarks}[11.10.3]
\label{IV.11.10.3}
\begin{enumerate}[label=(\roman*)]
\item A noção de família
esquematicamente
dominante é
\emph{local} em $X$,
como mostra, por
exemplo, a condição
\emph{b)} de
\sref{IV.11.10.1}.
Se $(W_\alpha)$
é uma cobertura
aberta de $X$,
$(f_\lambda)$ é
esquematicamente
dominante se, e
somente se, cada
família de restrições
$f_\lambda^{-1}
(W_\alpha)\to
W_\alpha$ o é.
\item Se $Z$ é o
pré-esquema soma dos
$Z_\lambda$ e
$f:Z\to X$ coincide
com $f_\lambda$
em cada
$Z_\lambda$, a
família
$(f_\lambda)$ é
esquematicamente
dominante se, e
somente se, a
família com o único
elemento $f$ o é.
\item Seja $M$ um
segundo conjunto
de índices e,
para cada
$\lambda\in L$,
seja
$(g_{\lambda\mu})
_{\mu\in M}$
uma família de
morfismos
$g_{\lambda\mu}:
Z'_{\lambda\mu}
\to Z_\lambda$.
Se cada família
$(g_{\lambda\mu})
_{\mu\in M}$ é
esquematicamente
dominante, então
$(f_\lambda)
_{\lambda\in L}$
é esquematicamente
dominante se, e
somente se, a família
$(f_\lambda\circ
g_{\lambda\mu})
_{(\lambda,\mu)
\in L\times M}$
o é
\sref{IV.11.9.2}.
\item Seja
$f:Z\to X$ um
morfismo tal que
$f_*(\sh{O}_Z)$
é um
$\sh{O}_X$-módulo
quase-coerente;
isso ocorre, por
exemplo, se $f$
é quase-compacto e
quase-separado
\sref{IV.1.7.4}.
Então $f$ ser
esquematicamente
dominante significa
que a
\emph{imagem fechada
de $Z$ por $f$}
\sref[I]{I.9.5.3}
é \emph{idêntica
a $X$}.
\end{enumerate}
\end{remarks}

\begin{proposition}[11.10.4]
\label{IV.11.10.4}
Se a família de
morfismos
$f_\lambda:
Z_\lambda\to X$
é esquematicamente
dominante, a reunião
dos
$f_\lambda(Z_\lambda)$
é densa em $X$.
Reciprocamente,
se essa reunião é
densa em $X$ e
$X$ é
\emph{reduzido},
então a família
$(f_\lambda)$ é
esquematicamente
dominante.
\end{proposition}

\begin{proof}
A primeira afirmação
decorre imediatamente
de
\sref{IV.11.10.1},~\emph{b)}.
Por outro lado,
suponhamos que $X$
seja reduzido e
que a reunião dos
$f_\lambda(Z_\lambda)$
seja densa em $X$.
Se existem as
fatorações
\sref{IV.11.10.1.2}
para todo
$\lambda\in L$,
também existem
as fatorações
\[
(f_\lambda^{-1}(U))
_{\mathrm{red}}
\xrightarrow{(g_\lambda)
_{\mathrm{red}}}
Y_{\mathrm{red}}
\xrightarrow{j_{\mathrm{red}}}
U.
\]
A hipótese implica
que o espaço
subjacente a $Y$
é idêntico a $U$.
Logo,
$Y=Y_{\mathrm{red}}
=U$, pois $U$
é reduzido.
\end{proof}

Os resultados de
\S~11.9 sobre
famílias separantes
traduzem-se em
resultados sobre
famílias
esquematicamente
dominantes.

\begin{theorem}[11.10.5]
\label{IV.11.10.5}
Sejam
$(f_\lambda)$
uma família de
morfismos
$f_\lambda:
Z_\lambda\to X$
e $g:X'\to X$
um morfismo.
Ponhamos
$Z'_\lambda=
Z_\lambda\times_X X'$
e
$f'_\lambda=
(f_\lambda)_{(X')}:
Z'_\lambda\to X'$.
\begin{enumerate}[label=(\roman*)]
\item Se $g$ é
fielmente plano
e a família
$(f'_\lambda)$ é
esquematicamente
dominante, então
$(f_\lambda)$ é
esquematicamente
dominante.
\item Suponhamos
que $g$ seja
plano e que,
além disso,
se verifique
uma das condições
seguintes:
\begin{enumerate}[label=\emph{\alph*)}]
\item $L$ é finito
e os $f_\lambda$
são quase-compactos.
\item O morfismo
$g$ é localmente
de apresentação
finita.
\end{enumerate}
Então, se
$(f_\lambda)$ é
esquematicamente
dominante, também
$(f'_\lambda)$ o é.
\end{enumerate}
\end{theorem}

\oldpage[IV-3]{172}

\begin{proposition}[11.10.6]
\label{IV.11.10.6}
Com as notações de
\sref{IV.11.10.5},
suponhamos que $X$
seja um pré-esquema
sobre um corpo $k$
e que, pondo
$S=\operatorname{Spec}(k)$,
se tenha
$X'=X\times_S S'$,
onde $S'$ é um
$k$-pré-esquema
qualquer. Então,
se a família
$(f_\lambda)$ é
esquematicamente
dominante, o mesmo
ocorre com
$(f'_\lambda)$.
\end{proposition}

\begin{corollary}[11.10.7]
\label{IV.11.10.7}
Sejam $X$ um
pré-esquema sobre
um corpo $k$ e
$(Z_\lambda)_{\lambda
\in L}$ uma família
de $k$-pré-esquemas
geometricamente
reduzidos sobre $k$;
para cada
$\lambda\in L$,
seja
$f_\lambda:
Z_\lambda\to X$
um $k$-morfismo.
Designemos por $Y$
o subpré-esquema
reduzido de $X$
cujo espaço
subjacente é o fecho
de
$\bigcup_{\lambda
\in L}f_\lambda
(Z_\lambda)$.
Sejam $k'$ uma
extensão de $k$ e
$X'$, $Z'_\lambda$
e
$f'_\lambda:
Z'_\lambda\to X'$
os pré-esquemas e
morfismos obtidos
de $X$,
$Z_\lambda$ e
$f_\lambda$ por
extensão da base
a $k'$. Então,
se $Y'$ é o
subpré-esquema
reduzido de $X'$
cujo espaço
subjacente é o fecho
de
$\bigcup_{\lambda
\in L}f'_\lambda
(Z'_\lambda)$,
tem-se
$Y'=Y\otimes_k k'$.
Em particular,
$Y$ é
geometricamente
reduzido sobre $k$.
\end{corollary}

\begin{proof}
Como os $Z_\lambda$
são reduzidos, os
morfismos
$f_\lambda$ se
fatoram como
\[
Z_\lambda
\xrightarrow{g_\lambda}
Y\xrightarrow{j}X,
\]
onde $j$ é a
imersão canônica
\sref[I]{I.5.2.2}.
Decorre então de
\sref{IV.11.10.4}
que $(g_\lambda)$
é uma família
esquematicamente
dominante. Ponhamos
$Y'_1=Y\otimes_k k'$,
subpré-esquema
fechado de $X'$,
e seja $g'_\lambda$
o morfismo obtido
de $g_\lambda$ por
extensão da base
a $k'$, de modo
que $f'_\lambda$
se fatora como
\[
Z'_\lambda
\xrightarrow{g'_\lambda}
Y'_1\xrightarrow{j'}X',
\]
onde $j'$ é a
imersão canônica.
Decorre de
\sref{IV.11.10.6}
que $(g'_\lambda)$
é uma família
esquematicamente
dominante. Mas,
por hipótese, os
$Z'_\lambda$ são
\emph{reduzidos};
portanto, por
\sref[I]{I.5.2.2},
os $g'_\lambda$
se fatoram como
\[
Z'_\lambda
\xrightarrow{h'_\lambda}
Y'\to Y'_1.
\]
Conclui-se então
de \sref{IV.11.10.2}
que $Y'=Y'_1$ e
$h'_\lambda=
g'_\lambda$, o
que demonstra o
corolário.
\end{proof}

\begin{definition}[11.10.8]
\label{IV.11.10.8}
Com as notações de
\sref{IV.11.10.5},
suponhamos que $X$
seja um
$S$-pré-esquema.
Diz-se que a família
$(f_\lambda)$ é
\emph{universalmente
esquematicamente
dominante}
(relativamente a $S$)
se, para toda
mudança de base
$S'\to S$, a família
$(f'_\lambda)$
correspondente à
mudança de base
$X'=X\times_S S'
\to X$ é
esquematicamente
dominante.
\end{definition}

Quando $S$ é o
espectro de um corpo,
a
\sref{IV.11.10.6}
significa, portanto,
que uma família
esquematicamente
dominante é
universalmente
esquematicamente
dominante
(relativamente a $S$).

Quando os
$f_\lambda$ são
imersões canônicas
$Z_\lambda\to X$,
diz-se também que a
família de
subpré-esquemas
$Z_\lambda$ é
\emph{universalmente
esquematicamente densa
em $X$}
(relativamente a $S$),
em vez de dizer que
a família
$(f_\lambda)$ é
universalmente
esquematicamente
dominante
(relativamente a $S$).

\begin{theorem}[11.10.9]
\label{IV.11.10.9}
Com as notações de
\sref{IV.11.10.5},
suponhamos que $X$
seja um
$S$-pré-esquema
localmente
noetheriano e que
todos os
$Z_\lambda$ sejam
planos sobre $S$.
Para cada $s\in S$,
sejam
\[
X_s=X\times_S
\operatorname{Spec}(k(s)),
\qquad
(Z_\lambda)_s=
Z_\lambda\times_S
\operatorname{Spec}(k(s)),
\qquad
(f_\lambda)_s=
f_\lambda\times1:
(Z_\lambda)_s\to X_s.
\]
Para que a família
$(f_\lambda)$ seja
universalmente
esquematicamente
dominante
relativamente a $S$,
é necessário e
suficiente que, para
todo $s\in S$, a
família
$((f_\lambda)_s)$
seja
esquematicamente
dominante.
\end{theorem}

\begin{proposition}[11.10.10]
\label{IV.11.10.10}
Sejam $f:X\to S$
um morfismo plano
e localmente de
apresentação finita
e $U$ um aberto
de $X$. Para que
$U$ seja
universalmente
esquematicamente
denso em $X$
relativamente a $S$,
é necessário e
suficiente que, para
todo $s\in S$,
$U_s=U\cap X_s$
seja
esquematicamente
denso em $X_s$
(ou, de modo
equivalente, que
$\operatorname{Ass}
(\sh{O}_{X_s})
\subset U_s$).
\end{proposition}
